% !TeX encoding = UTF-8

\documentclass[fontset=none]{exam-zh}
\usepackage{mathtools}
\usepackage{tcolorbox}
% 431 项目统一文档主题。
% 各根入口应先声明 \documentclass{exam-zh}，再输入本文件。

% 项目统一中文字体：避免依赖某台电脑上的 STSong/STHeiti 等系统字体。
\setCJKmainfont{SourceHanSerifCN-Regular.ttf}[
    Path = fonts/,
    BoldFont = SourceHanSerifCN-Bold.ttf,
    ItalicFont = AR-PL-KaitiM-GB.ttf,
    BoldItalicFont = AR-PL-KaitiM-GB.ttf,
    BoldItalicFeatures = {FakeBold = 2},
]
\setCJKsansfont{SourceHanSansCN-Regular.ttf}[
    Path = fonts/,
    BoldFont = SourceHanSansCN-Bold.ttf,
]
\setCJKmonofont{SourceHanSansCN-Regular.ttf}[
    Path = fonts/,
]

% fontset=none 时 CTeX 不创建这些常用字体族命令，显式补齐。
\setCJKfamilyfont{zhsong}{SourceHanSerifCN-Regular.ttf}[Path = fonts/]
\setCJKfamilyfont{zhhei}{SourceHanSansCN-Bold.ttf}[Path = fonts/]
\setCJKfamilyfont{zhkai}{AR-PL-KaitiM-GB.ttf}[Path = fonts/]
\setCJKfamilyfont{zhfs}{SourceHanSerifCN-Regular.ttf}[Path = fonts/]
\providecommand{\songti}{\CJKfamily{zhsong}}
\providecommand{\heiti}{\CJKfamily{zhhei}}
\providecommand{\kaishu}{\CJKfamily{zhkai}}
\providecommand{\fangsong}{\CJKfamily{zhfs}}


\usepackage{booktabs}
\usepackage{multirow}
\usepackage{threeparttable}
\usepackage{threeparttablex}
\usepackage{tikz}
\usetikzlibrary{calc}
\usetikzlibrary{arrows.meta}
\usetikzlibrary{decorations.pathmorphing}
\usetikzlibrary{decorations.pathreplacing}
\usetikzlibrary{positioning}

\tikzset{
    mindmap_root/.style={rectangle, rounded corners, draw=black, fill=gray!10, thick, align=center, minimum height=0.8cm, font=\bfseries\small},
    mindmap_level1/.style={rectangle, rounded corners, draw=blue!60, fill=blue!5, thin, align=center, minimum height=0.6cm, font=\small},
    mindmap_level2/.style={rectangle, draw=none, fill=none, align=center, font=\footnotesize, text width=3cm},
    arrow/.style={->, >=stealth, thick, gray}
}

\usepackage{tabularx}
\usepackage{xcolor}
\usepackage{makecell}
\usepackage{extarrows}
\usepackage{longtable}
\usepackage{float}
\usepackage{enumitem}
\usepackage{perpage}
\usepackage{scalerel}
\usepackage{stackengine}
\MakePerPage{footnote}
\stackMath
\tcbuselibrary{breakable}
\tcbset{breakable}

% 适用于较长估计变量名的宽帽号。
\newcommand{\reallywidehat}[1]{%
  \savestack{\tmpbox}{\stretchto{%
    \scaleto{%
      \scalerel*[\widthof{\ensuremath{#1}}]{\kern-.6pt\bigwedge\kern-.6pt}{%
        \rule[-\textheight/2]{1ex}{\textheight}%
      }%
    }{\textheight}%
  }{0.5ex}}%
  \stackon[1pt]{#1}{\tmpbox}%
}

% 列表全局设置
\setlist{nosep}
% 列表内部保持正文节奏；整个列表块与上下内容各保留半行间距。
\setlist[itemize]{
    leftmargin=*, labelsep=0.5em,
    itemsep=0pt, parsep=0pt,
    topsep=0.5\baselineskip, partopsep=0pt,
    after={}
}
\setlist[itemize,1]{label=$\bullet$}
\setlist[itemize,2]{label=\small$\bullet$}
\setlist[itemize,3]{label=\scriptsize$\bullet$}
\setlist[enumerate]{
    leftmargin=*, labelsep=0.5em,
    itemsep=0pt, parsep=0pt,
    topsep=0.5\baselineskip, partopsep=0pt,
    after={}
}
\setlist[description]{leftmargin=\parindent, labelindent=0pt}

% 章节公共样式。需要每节另起一页的资料在入口中单独覆盖 section/break。
\IfFileExists{ctex.sty}{
    \ctexset{
        section = {
            format = \Large\bfseries\raggedright,
            beforeskip = 1.5ex plus 0.5ex minus 0.2ex,
            afterskip = 1.5ex plus 0.2ex,
        },
        subsection = {
            format = \large\bfseries\raggedright,
            beforeskip = 1.0ex plus 0.5ex minus 0.2ex,
            afterskip = 1.0ex plus 0.2ex,
        },
        subsubsection = {
            format = \normalsize\bfseries\raggedright,
            beforeskip = 0.5ex plus 0.5ex minus 0.2ex,
            afterskip = 0.5ex plus 0.2ex,
        },
        paragraph = {
            runin = false,
            number = \arabic{paragraph}.,
            format = \normalsize\bfseries\raggedright,
            indent = 2em,
            beforeskip = 0.5ex plus 0.3ex minus 0.1ex,
            afterskip = 0.3ex plus 0.2ex,
        }
    }
}{}

% 按 chapter 计数器自动生成“第 X 讲”。
% exam-zh 在 \chapter 完成时推进 chapter，因此标题使用当前值加一。
\ExplSyntaxOn
\NewDocumentCommand{\lecture}{m}
  {
    \exp_args:Nx \chapter
      {第\zhnumber{\int_eval:n {\value{chapter} + 1}}讲：\exp_not:n {#1}}
  }
\ExplSyntaxOff

% 431 资料通用的 exam-zh 设置；可选参数为居中页眉，必选参数为右侧页眉。
\newcommand{\SetupFourThreeOneDocument}[2][]{
    \examsetup{
        page={
            size = a4paper,
            show-head = true,
            head-content = {
                \fancyhf{}
                \fancyhead[C]{#1}
                \fancyhead[EL]{\leftmark}
                \fancyhead[OR]{#2}
            },
            foot-content = {\hfill 第;页 / 共;页 \hfill}
        },
        style/student-version-cleanaux = true,
        fillin/type = paren,
        question = {
            show-points = auto,
            show-answer = true,
        },
        problem = {
            points-separate-par = false,
        },
        paren = {
            text-color = blue,
            show-paren = true,
            show-answer = true,
        },
        solution = {
            show-solution = show-stay,
            label-indentation = true,
            text-color = blue,
            blank-type = hide,
        },
    }

    \ExamPrintAnswerSet{
        question/show-answer = false,
        paren/show-answer = false,
        solution/show-solution = hide,
        paren/text-color = black,
        solution/text-color = black,
    }
}

\renewcommand{\chaptermark}[1]{\markboth{#1}{}}
\renewcommand{\headrulewidth}{1pt}

\newcommand{\true}{$\checkmark$}
\newcommand{\false}{$\times$}

\usepackage{hyperref}
\usepackage{tocloft}

\renewcommand{\contentsname}{\hspace*{\fill}目录\hspace*{\fill}}
\setlength{\cftbeforetoctitleskip}{2em}
\setlength{\cftaftertoctitleskip}{2em}
\setlength{\cftchapindent}{0em}
\setlength{\cftsecindent}{2.0em}
\setlength{\cftsecnumwidth}{2em}
\setlength{\cftsubsecindent}{4.0em}
\setlength{\cftsubsecnumwidth}{2.0em}
\setlength{\cftsubsubsecindent}{6.0em}
\setlength{\cftsubsubsecnumwidth}{2.0em}

\setlength{\intextsep}{1em}

\hypersetup{
    colorlinks = true,
    linkcolor = blue,
    urlcolor = blue,
    citecolor = blue,
    pdfborder = {0 0 0},
}

\newcommand{\kref}[1]{\ref{#1} \textsuperscript{(p.\pageref{#1})}}

\usepackage[group-separator={,}]{siunitx}

\tcbuselibrary{breakable}
\tcbset{breakable}
\usetikzlibrary{intersections}

\SetupFourThreeOneDocument{金融经济学习题讲义（王江）}

\title{金融经济学习题}
\subject{王江《金融经济学》习题解答}
\author{Bilibili@园丁小区詹姆斯}

% 原书题号含星号（如 2.2*），不用 exam-zh 自动序号。
\examsetup{
    problem = {
        label = {},
        label-align = left,
        show-points = false,
    }
}

\newcommand{\problemno}[1]{\textbf{#1}\quad}

\begin{document}

\FourThreeOneFrontMatter

\ifdefined\FourThreeOnePrintInterior
\else
\MakeFourThreeOneCover
    {431 FINANCE \quad/\quad FINANCIAL ECONOMICS}
    {金融经济学习题讲义}
    {王江《金融经济学》习题解答}
    {从偏好与状态价格出发，完成资产定价的推演}
    {Bilibili@园丁小区詹姆斯}
    {PREFERENCE \quad EQUILIBRIUM \quad PRICING}

\null
\thispagestyle{empty}
\fi

\MakeFourThreeOneContents[2]

\cleardoublepage

% 原书解答从第 2 章开始，无第 1 章习题。
\setcounter{chapter}{1}
% !TeX root = ../../金融经济学习题讲义.tex
% 源 PDF：references/金融经济学/史上最全最完整的金融经济学答案(王江版).pdf，页 2–6（印刷页 1–5）

\zhchapter{基本框架}

\begin{tcolorbox}[knowledge,title=2.1 经济环境]
\begin{enumerate}

\item 形式上可以用一棵状态树来表示上面描述的经济环境。用 $\omega$ 标识状态，假设总共有 $\Omega$ 个状态，那么表示所有状态可能性的集合 $\Omega$ 称为\textbf{状态空间}。记状态发生的概率为 $p(\omega)$，所有状态发生的概率为 $P=\{p_{\omega}, \omega\in\Omega\}$。

\item 假设在经济中只有一种不可储存的商品可以用作消费和生产。
\item 形式上可以用一种状态出来描述上面说的经济环境：
\begin{center}
\WJEventTree[left={$\Omega$},stem]{
    \omega=1 \text{，概率为 }p_1,
    \omega=2 \text{，概率为 }p_2,
    \vdots,
    \omega=\Omega \text{，概率为 }p_{\Omega}
}
\end{center}
\item 未来两种状态（好状态和坏状态），发生概率相等，用产出水平来描述状态：
\begin{center}
    \WJEventTree[left={$\Omega$},stem]{
        \omega=200 \text{，好状态$\omega$=1，概率为 1/2},
        \omega=50 \text{，坏状态$\omega$=2，概率为 1/2}
    }
    \end{center}
\end{enumerate}
\end{tcolorbox}


\begin{tcolorbox}[knowledge,title=2.2 经济参与者]
    把经济参与者描述为：为自己切身利益而行动的个体。一个参与者的经济特征包括：
    \begin{enumerate}
        \item 他控制的经济资源（任何有经济价值，可以用于满足经济需求的东西，可以是有形 or 无形的）。参与者的经济资源包括
        \begin{itemize}
            \item 所占有的实物商品/资本品
            \item 信息
            \item 生产技术
        \end{itemize}
        \item 他的经济需求。
    \end{enumerate}

\tcbline

参与者的经济资源：
\begin{enumerate}
    \item 禀赋：即参与者初始占有的资源，是他与生俱来的商品、资本品。给定两个时期，且第二时期有多个状态，那么他的禀赋可表示为：
    \[
    e_k=\begin{bmatrix}e_{k,0} \\ e_{k,1}\end{bmatrix}
    =\begin{bmatrix}e_{k,0} \\ \begin{bmatrix}e_{k,11} \\ e_{k,12}\end{bmatrix}
    \end{bmatrix}
    \quad k=1,\cdots,K
    \]

    在数学上，可以把参与者的禀赋看成是 $1+\Omega$ 维实空间 $\mathrm{R}^{1+\Omega}$ 中的一个元素，且假设参与者的禀赋非负，那么禀赋 $e_k \in \mathrm{R}^{1+\Omega}_{+}$。
\begin{tcolorbox}[warning,title=向量的“大于”]
    对于向量 $a=\begin{bmatrix}a_1 \\ a_2 \\ \vdots \\ a_n\end{bmatrix}$，有以下定义：
    \begin{enumerate}
        \item $a \ge 0$，即 a 的所有分量都非负，即 $a_i \ge 0, i=1,\cdots,n$。
        \item $a > 0$，即 a 的所有分量都非负，且至少有一个分量严格大于 0，即 $\exists i \in \{1,\cdots,n\}$ 使得 $a_i > 0$。
        \item $a \gg 0$，即 a 的所有分量都严格大于零，即 $a_i > 0, i \in \{1,\cdots,n\}$。
    \end{enumerate}
\end{tcolorbox}
    \item 信息：信息具备经济价值。参与者的信息可以分为两种类型：公开信息和私有信息。如果参与者仍然拥有私有信息，信息就是不对称的。信息不对称会导致很多问题，例如道德风险和逆向选择。但在本书中，将假设经济者对于经济，特别是其未来状态发生的可能性，拥有相同的信息。
    \item 生产技术：指可以把商品转化为更多商品，或者把现在的商品转化为未来的商品。在本书中，假设所有参与者不拥有任何生产技术。
\end{enumerate}

\tcbline

参与者的经济需求：
参与者的经济需求基于他的不同时期和未来不同状态下的消费。对于参与者 k，记他在第零期的消费是$c_{k,0}$，第一期的消费是 $c_{k,\omega}$，其中 $\omega \in \Omega$。消费的完整表述是：
\[
c_k \equiv \begin{bmatrix}c_{k,0} \\ c_{k,1}\end{bmatrix}=
\begin{bmatrix}c_{k,0} \\ \begin{bmatrix}c_{k,11} \\ \vdots \\ c_{k,1\omega} \\ \vdots \\ c_{k,1\Omega}\end{bmatrix}
\end{bmatrix} \quad k=1,\cdots,K
\]


\begin{enumerate}
\item 消费集：消费者可能的消费选择叫做消费计划 $c=\begin{bmatrix}
    c_0 \\ c_1
\end{bmatrix}$。所有消费计划的集合叫做消费集 $C$。和禀赋一样，消费计划属于 $\mathrm{R}^{1+\Omega}$ 中的一个元素，且假设消费计划非负，那么消费计划 $c \in \mathrm{R}^{1+\Omega}_{+}$，消费集就是 $C=\mathrm{R}^{1+\Omega}_{+}$。

\begin{tcolorbox}[knowledge]
\textbf{定义 2.1} 设 A 是 $R_n$ 的一个集合，对于 $\forall a,b \in A$ 和 $\alpha \in [0,1]$，如果 $a\alpha+(1-\alpha)b$ 也在A 中，则称集合 A 是凸的（convex）。

\textbf{定义 2.2} 对于 $R^n$的集合 A，如果其中任意序列 $a_i(i=1,\cdots,\infty)$ 有极限，且极限也在 A 中，则称集合 A 是闭的（closed）。

\textbf{假设 1} 消费集 $C=\mathrm{R}^{1+\Omega}_{+}$ 是 $\mathrm{R}^{1+\Omega}$ 中的一个闭凸子集。
\end{tcolorbox}

\item 偏好：偏好是参与者对所有可能性计划的排序这样的排序定义了他的经济需求。

\begin{tcolorbox}[knowledge]
\textbf{定义 2.3} 偏好是C 上的一个二元关系，表示为 $\succeq$，且满足以下三个条件：
\begin{itemize}
    \item 完备性：对于 $\forall a,b \in C$，有 $a\succeq b$ 或 $b\succeq a$。
    \item 传递性：对于 $\forall a,b,c \in C$，如果 $a\succeq b$ 且 $b\succeq c$，则 $a\succeq c$。
    \item 自反性：对于 $\forall a \in C$，有 $a\succeq a$。
\end{itemize}
\end{tcolorbox}

\item 偏好的基本假设：\textcolor{red}{\textbf{参与者喜多厌少}：如果两个消费计划 a 和 b 之间，在所有时期和状态下，a 都不比 b 提供的消费少，且在有些状态下提供更多的消费，那么 a 优于 b。}
    \begin{tcolorbox}[knowledge]
        \textcolor{red}{\textbf{公理 1 不满足性：}} 如果 $a>b$，那么 $a\succeq b$。

        \textbf{公理 2 连续性：} $\forall c \in C$，集合 $\{a\in C: a\succeq c\}$ 和 $\{b\in C: b\preceq c\}$ 都是闭集。

        \textbf{公理 3 凸性：} $\forall a,b \in C$ 以及 $\alpha \in [0,1]$，如果 $a \succeq b$ ，那么有 $\alpha a+(1-\alpha)b\succeq b$ 。
    \end{tcolorbox}

    \item 效用函数：从消费计划到实数的映射叫做效用函数。
    
    \begin{tcolorbox}[knowledge]
        \textbf{定义 2.4} 对应于偏好关系的效用函数 U 是从 C 到 R 的函数，使得 $\forall a,b \in C$，当且仅当 $a\succeq b$ 时，有 $U(a) \ge U(b)$。
    \end{tcolorbox}

    \begin{tcolorbox}[knowledge]
        \textbf{定理 2.1} 对于一个在闭凸消费集 C 上，有定义 2.3 定义的（完备性、传递性和自反性）、满足公理 2 的偏好（连续性），存在一个连续的效用函数 U 使得 $\forall a,b \in C, a\succeq b$，当且仅当 $U(a) \ge U(b)$。
    \end{tcolorbox}

    \begin{tcolorbox}[example]
        \begin{question}
            \textbf{例 2.3}\quad 在例 2.1 中描述的「Lucas 树」经济中，消费空间是 $\mathrm{R}^{3}_{+}$，表示在今年和明年的两个可能状态下的消费；状态 $a$ 对应于好天气，而状态 $b$ 则对应于坏天气。假设经济中的某个参与者有如下的效用函数：
            \[
                U(c_0,c_{1a},c_{1b})=\log c_0+\tfrac12(\log c_{1a}+\log c_{1b})
            \]
            在如下三个消费计划中，他将会选择哪一个？
            \begin{center}
            \begin{tabular}{@{}c@{\quad}c@{\quad}c@{}}
                \textbf{计划 A} & \textbf{计划 B} & \textbf{计划 C} \\[0.35em]
                \WJEventTree[left={$1$},stem]{1,1} &
                \WJEventTree[left={$\tfrac12$},stem]{\tfrac32,\tfrac32} &
                \WJEventTree[left={$1$},stem]{\tfrac32,\tfrac12}
            \end{tabular}
            \end{center}
        \end{question}

        \begin{solution}
            对于 A、B 和 C，效用函数给出的相应的效用值分别为 $0$、$-0.49$ 和 $-0.14$。显然，参与者对三个消费计划的排序是 $A>C>B$，因此他会选择 A。
        \end{solution}

        \begin{question}
            \textbf{例 2.4}\quad 让我们考虑同一经济中的另外一个参与者，他的效用函数是
            \[
                V(c_0,c_{1a},c_{1b})=c_0\sqrt{c_{1a}c_{1b}}
            \]
            他对上面三个消费计划的排序是什么？
        \end{question}
        \begin{solution}
            \[
                \begin{aligned}
                    V_A &= 1\cdot\sqrt{1\cdot 1}=1,\\
                    V_B &= \tfrac12\cdot\sqrt{\tfrac32\cdot\tfrac32}=\tfrac34,\\
                    V_C &= 1\cdot\sqrt{\tfrac32\cdot\tfrac12}=\tfrac{\sqrt3}{2}.
                \end{aligned}
            \]
            因此 $V_A>V_C>V_B$，参与者对三个消费计划的排序仍是 $A>C>B$。
        \end{solution}
    \end{tcolorbox}
\end{enumerate}

上面的例子说明了如下的结论：从定理 2.1 中得到的效用函数只是\textbf{序数的}（ordinal），即它只给出了消费计划的排序。效用函数的任意正单调变换不会改变这种排序。效用函数的实际数值并不重要。

利用效用函数作为偏好的一种表述形式，我们可以把偏好的三个基本假设重新表示如下：
\begin{enumerate}
    \item \textbf{不满足性：}如果 $a\ge b$，那么 $U(a)\ge U(b)$。
    \item \textbf{连续性：}$\lim_{a_n \to a}U(a_n)=U(a)$。
    \item \textbf{凸性：}如果 $U(a)>U(b)$ 且 $\alpha\in(0,1)$，那么 $U(\alpha a+(1-\alpha)b)>U(b)$。
\end{enumerate}

\end{tcolorbox}



\begin{tcolorbox}[knowledge,title=2.3 证券市场]
资源的配置是通过金融市场中的交易来完成的，市场允许参与者在不同时期和状态上配置资源。

\tcbline

\begin{enumerate}
    \item \textbf{证券及其支付：}在我们所要构建的简单框架中，金融市场由一组\textbf{证券}（security）构成。一只证券即是一份\textbf{金融要求权}（financial claim），它在 1 期给其所有者带来\textbf{支付}（payoff）；支付的数量一般依赖于当时的经济状态。令 $x_\omega$ 为证券在状态 $\omega$ 时的支付，其中 $\omega\in\Omega$，则一只证券由它在各个可能状态下的支付 $\{x_\omega,\,\omega=1,\dots,\Omega\}$ 来定义。形式上可用一棵\textbf{支付树}表示：
    \begin{center}
    \WJEventTree[left={$0$},stem]{x_1,x_2,\vdots,x_\Omega}
    \end{center}
    定义 $\mathrm{R}^{\Omega}$ 为\textbf{支付空间}（payoff space），即所有可能支付的集合；其中的向量
    \[
        x=\begin{bmatrix}x_1 \\ x_2 \\ \vdots \\ x_\Omega\end{bmatrix}
    \]
    就定义了一只证券。在定义证券时，我们未考虑任何在 0 期的可能支付；到期以前（在我们的框架中就是 0 期）证券产生的支付通常叫做\textbf{股利}（dividend）。
    
    在 Lucas 树例子中，把代表树的所有权的证券叫做\textbf{股票}（stock），它代表对实物资产（树）所产生的支付的要求权；这样的证券也叫做\textbf{金融资产}（financial asset）或\textbf{资产}，其支付来自\textbf{标的}（underlying）实物资产。

    \textbf{无风险债券}（riskless bond）是这样一种证券，它在所有状态下的支付都相等；若各状态下的支付均为 $1$，则其支付向量为 $\begin{bmatrix}1 \\ \vdots \\ 1\end{bmatrix}$。
    
    \textbf{风险证券}（如股票）的支付则随状态变化，例如支付向量为 $\begin{bmatrix}x_1 \\ \vdots \\ x_\Omega\end{bmatrix}$。
    
    金融证券具有三个基本特征：
    \begin{enumerate}
        \item 支付只依赖于未来的经济状况；
        \item 支付是外生给定的，不受参与者行为的影响；
        \item 状态集合及其概率分布是所有参与者的共同知识。
    \end{enumerate}

    \item \textbf{市场结构：}假设市场上有 $N$ 只可交易证券，用 $n=1,\dots,N$ 来标识。记第 $n$ 只证券的支付向量为
    \[
        X_n=\begin{bmatrix}x_{1,n} \\ x_{2,n} \\ \vdots \\ x_{\Omega,n}\end{bmatrix},
    \]
    把所有证券的支付向量并排，就得到 $\Omega\times N$ 的\textbf{市场结构矩阵}（也称支付矩阵）
    \[
        X\equiv[X_{\cdot,1},\,\cdots\,X_{\cdot,n},\,\cdots\,X_{\cdot,N}]=
        \begin{bmatrix}
            x_{1,1}   & x_{1,2}   & \cdots & x_{1,N}   \\
            \vdots    & \vdots    & \ddots & \vdots    \\
            x_{\omega,1}   & x_{\omega,2}   & \cdots & x_{\omega,N}   \\
            \vdots    & \vdots    & \ddots & \vdots    \\
            x_{\Omega,1} & x_{\Omega,2} & \cdots & x_{\Omega,N}
        \end{bmatrix}\equiv
        \begin{bmatrix}
            X_{1,\cdot} \\ \vdots \\ X_{\omega,\cdot} \\ \vdots \\ X_{N,\cdot}
        \end{bmatrix}
    \]
    矩阵 $X$ 的每一列是一只证券在各状态下的支付（列向量），每一行是一个状态下各证券的支付（行向量）。通常有 $\Omega>N$。

    \item \textbf{证券组合：}对每只证券的持有数量构成\textbf{证券组合}（security portfolio），简称\textbf{组合}（portfolio）。若 $\theta_n$ 为证券 $n$ 的持有数量，则\textbf{组合}由向量$\theta=\begin{bmatrix}\theta_1 \\ \theta_2 \\ \vdots \\ \theta_N\end{bmatrix}$ 定义；该组合在 1 期的支付向量为 $X\theta$。对每个参与者 $k$，记其交易前的\textbf{初始证券组合}为$\bar{\theta}_k=\begin{bmatrix}\bar{\theta}_{k,1} \\ \bar{\theta}_{k,2} \\ \vdots \\ \bar{\theta}_{k,N}\end{bmatrix}$，所有参与者初始证券组合之和称为\textbf{市场组合}（market portfolio）$\theta_M$，即
    \[
        \theta_M=\sum_{k=1}^{K}\bar{\theta}_k,
    \]
    它表示市场上所有可交易证券的总供给。

    \item \textbf{市场化：}若支付 $x\in\mathrm{R}^{\Omega}$ 能由某交易组合复制出来，即存在 $\theta\in\mathrm{R}^{N}$ 使得 $X\theta=x$（有解），则称 $x$ 是\textbf{市场化的}（marketed）。所有市场化支付的集合为
    \[
        M=\{X\theta:\,\theta\in\mathrm{R}^{N}\}.
    \]
    一般地，$M$ 是支付空间 $\mathrm{R}^{\Omega}$ 的一个子空间，其维数不超过 $N$。

    \item \textbf{交易过程：}交易可以通过\textbf{集中化交易}（centralized exchange）或\textbf{柜台交易}（over-the-counter，OTC）完成；具体的交易安排称为\textbf{市场微观结构}（market micro-structure）。记 $S=\begin{bmatrix}S_1 \\ S_2 \\ \vdots \\ S_N\end{bmatrix}$ 为 0 期各交易证券的价格向量，$\theta_k(S)=\begin{bmatrix}\theta_{k,1}(S) \\ \vdots \\ \theta_{k,N}(S)\end{bmatrix}$ 为参与者 $k$ 在价格 $S$ 下对证券的（净）需求。使各证券总需求等于总供给的价格 $S$ 称为\textbf{市场出清价格}。若记初始组合为 $\bar{\theta}_k$，则总持有量出清为 $\sum_{k=1}^{K}\theta_k(S)=\sum_{k=1}^{K}\bar{\theta}_k$；在均衡分析中常设 $\bar{\theta}_k=0$（或将初始组合并入禀赋），此时净需求出清即 $\sum_{k=1}^{K}\theta_k(S)=0$。

    交易成本、头寸限制（如卖空约束）、个体交易对价格的影响以及税收等，都会妨碍交易，统称为\textbf{市场摩擦}。\textbf{无摩擦市场}满足以下假设：
    \begin{enumerate}
        \item 参与者可无成本地进入或退出证券市场；
        \item 没有交易成本；
        \item 对参与者的证券持有没有限制（允许卖空）；
        \item 单个参与者的交易不影响市场价格（价格接受者）；
        \item 没有税收。
    \end{enumerate}
\end{enumerate}
\end{tcolorbox}

\begin{tcolorbox}[knowledge,title=2.4 基本经济模型]
现在我们已经对将要分析的经济的基本模型有了完整的描述。在作进一步的讨论和分析以前，先作一个总结。

\begin{tcolorbox}[knowledge]
\textbf{定义 2.5}\quad 一个\textbf{经济}定义如下：
\begin{enumerate}
    \item 有两个时期 $0$ 和 $1$；在 1 期有 $\Omega$ 个可能状态，对这些状态有概率测度 $P$。经济中只有一种不可储存的商品。
    \item 经济中有 $K$ 个参与者（$k=1,\dots,K$）：
    \begin{enumerate}[label=(\alph*)]
        \item 每一参与者对未来状态发生的可能性都有\textbf{相同的信息}，由 $P$ 描述；
        \item 每一参与者具有\textbf{禀赋} $e_k\in\mathrm{R}^{1+\Omega}_{+}$；
        \item 每一参与者有定义于 $C=\mathrm{R}^{1+\Omega}_{+}$ 上、且满足公理 1$\sim$公理 3 的偏好（不满足性、连续性、凸性）。
    \end{enumerate}
    \item 有一个市场结构为 $X$ 的\textbf{无摩擦证券市场}。
\end{enumerate}
\end{tcolorbox}

\begin{tcolorbox}[knowledge]
\textbf{定义 2.6}\quad 在上面定义的经济中，如果所有参与者的 1 期禀赋都可表示为其初始证券组合的支付，即对每个 $k$ 有 
\[
e_{k,1}=X\bar{\theta}_k
\]
则我们称之为\textbf{证券市场经济}（security-market economy）。
\end{tcolorbox}

为简化起见，在以后的讨论中有时会局限于证券市场经济的情形。
\end{tcolorbox}

\begin{tcolorbox}[knowledge,title=2.5 市场均衡]
在证券市场中，资源配置通过以下两方面的相互作用来实现：
\begin{enumerate}
    \item 给定证券及其未来支付与当前价格，每个参与者选择最优证券持有量；
    \item 所有参与者的总需求影响证券价格；当需求等于供给时达到\textbf{均衡}（equilibrium），此时各参与者都选择了最优证券需求且市场出清。
\end{enumerate}

\tcbline

\begin{enumerate}
    \item \textbf{参与者各自优化：}记 $\theta=\begin{bmatrix}\theta_1 \\ \vdots \\ \theta_N\end{bmatrix}$ 为\textbf{净购买}的证券组合（相对初始持有量的交易量），$\theta_n>0$ 为\textbf{多头}，$\theta_n<0$ 为\textbf{空头}。该组合在 0 期的市场价值（成本）为 $S^{\mathsf T}\theta$。对禀赋为 $e_k$ 的参与者，净购买 $\theta_k$ 后消费计划 $c_k=\begin{bmatrix}c_{k,0} \\ c_{k,1}\end{bmatrix}$ 为
    \[
        c_{k,0}=\underbrace{e_{k,0}-S^{\mathsf T}\theta_k}_{\text{第零期禀赋减去购买证券的支出}},\qquad
        c_{k,1}=\underbrace{e_{k,1}+X\theta_k}_{\text{第 1 期禀赋加上证券在第一期的支付}}.
    \]
    其中 $-S^{\mathsf T}\theta_k$ 为 0 期储蓄（为正）或借贷（为负）。所有可达消费计划的集合称为\textbf{预算集}
    \[
        B(e_k,\{X,S\})=\Bigl\{c\ge 0\;\Big|\;
        c_{k,0}=e_{k,0}-S^{\mathsf T}\theta_k,\;
        c_{k,1}=e_{k,1}+X\theta_k,\;
        \theta_k\in\mathrm{R}^{N}\Bigr\}.
    \]
    参与者选择 $\theta_k$ 以最大化效用：
    \[
    \begin{aligned}
        &\max_{\theta_k} U_k(c_k) \\
        &\text{s.t.}\quad
        c_{k,0}=e_{k,0}-S^{\mathsf T}\theta_k,\;
        c_{k,1}=e_{k,1}+X\theta_k,\;
        c_{k,0},\,c_{k,1}\ge 0
    \end{aligned}
    \]
    等价地，
    \[
        \max_{c_k\in B(e_k,\{X,S\})} U_k(c_k).
    \]
    该问题的解给出参与者 $k$ 对证券的\textbf{需求函数} $\theta_k(e_k,S)$，它依赖于价格向量 $S$ 与禀赋 $e_k$。

    \item \textbf{市场出清：}在价格向量 $S$ 下，记参与者 $k$ 的证券需求为 $\theta_k(e_k,S)$（即上面的净购买量）。已知，参与者的禀赋是 $(e_0,e_1)$ 和初始组合 $\bar{\theta}$，那么可以定义等价的总禀赋（即，将初始证券的形式，从第零期的持有 $\bar{\theta}_k$，转换为第 1 期的支付 $X\bar{\theta}_k$）
    \[
        e_k^T = 
        \begin{bmatrix} e_{k,0} \\ e_{k,1} + X\bar{\theta}_k\end{bmatrix}
    \]
    如此转化后：
    \begin{itemize}
        \item 对于参与者的禀赋，我们把$e^T$当做他的总禀赋，初始证券持有量 $\bar\theta_k=0$，他选择的任意证券需求就变成了初始组合之外的额外持有量 $\theta_k$
        \item 对于后续优化问题里的$\theta_k$，就表示在第零期市场上净买卖的量。
    \end{itemize}


    
    如此一来，\textbf{证券市场出清}要求有人净买 $\theta_k>0$，就要有人净卖 $\theta_k<0$。对所有参与者加总：净买入之和必须为 0，否则多出来的份额没人卖、或少掉的份额没人接，即净需求为零：
    \[
        \underbrace{\sum_{k=1}^{K}\theta_k(e_k,S)=0}_{\text{证券市场出清条件}}.
    \]
    满足该条件的 $S$ 即为\textbf{均衡价格}。由 $c_{k,0}=e_{k,0}-S^{\mathsf T}\theta_k(e_k,S)$，对 $k$ 求和得
    \[
        \underbrace{\underbrace{\sum_{k=1}^{K}c_{k,0}}_{\text{0 期商品总需求}}
        =\sum_{k=1}^{K}e_{k,0}-S^{\mathsf T}\Bigl(\sum_{k=1}^{K}\theta_k(e_k,S)\Bigr)
        =\underbrace{\sum_{k=1}^{K}e_{k,0}}_{\text{0 期总供给}}}_{\text{第0期商品市场的出清}}
    \]
    左边是第0 期商品的总需求，右边是第0期的总供给，表示的是第0期商品市场的出清。同理，对每个状态 $\omega=1,\dots,\Omega$ 有
    \[
        \sum_{k=1}^{K}c_{k,\omega}=\sum_{k=1}^{K}e_{k,\omega},
    \]
    即\textbf{商品市场出清}。合起来得到 \textbf{Walras 法则}（这表明，实物资源没有凭空产生或消失，只是通过交易被重新分配）：
    \[
        \sum_{k=1}^{K}c_k=\sum_{k=1}^{K}e_k.
    \]
    

    求均衡时，\textbf{首先}对任意价格 $S$ 解出各参与者的最优证券组合 $\theta_k(e_k,S)$，\textbf{第二}用市场出清条件解出均衡价格 $S$。形式上，
    \[
        S=S\bigl(P;\,\{U_k,\,e_k,\,k=1,\dots,K\};\,X\bigr),
    \]
    即证券价格取决于经济面临的风险 $P$、参与者的偏好 $U_k$ 与禀赋 $e_k$ 以及市场结构 $X$ 等\textbf{基本面}（fundamentals）或\textbf{本源}（primitives）：经济面临的风险，经济中参与者的偏好和禀赋以及证券市场的结构。
\end{enumerate}
\end{tcolorbox}

\begin{tcolorbox}[knowledge,title=2.6 最优性]
通过证券市场进行的资源配置是否已是\textbf{最好的可能配置}？参与者各自选择最优消费计划，并不意味着社会资源配置一定最优；若市场不完备，结果可能是\textbf{次优的}。

\begin{tcolorbox}[knowledge]
\textbf{定义 2.7}\quad 配置 $\{c_k,\,\forall k\}$ \textbf{Pareto 占优于}配置 $\{c_k',\,\forall k\}$，如果
\[
    \forall k,\quad U_k(c_k)\ge U_k(c_k'),
\]
且至少对一个参与者严格不等式成立。

这就是说，对比两套分配结果，如果新配置帕累托占优于旧配置，说明：新配置对所有参与者的效用都至少不比原来差，且至少对一个参与者的效用严格提高，且这种“好”不以“损人”为前提，这是一种\textbf{更具效率}的体现。
\end{tcolorbox}



\begin{tcolorbox}[knowledge]
\textbf{定义 2.8}\quad 给定经济的总供给 $\{e_k,\,\forall k\}$，配置 $\{c_k,\,\forall k\}$ 是\textbf{可行的}（feasible），如果
\[
    \underbrace{\sum_{k=1}^{K}c_k=\sum_{k=1}^{K}e_k}_{\text{说明所有状态下所有人的消费之和等于总禀赋之和，没有闲置}}.
\]
\end{tcolorbox}

\begin{tcolorbox}[knowledge]
\textbf{定义 2.9}\quad 若配置 $\{c_k,\,\forall k\}$ 可行，且不存在其他可行配置 Pareto 占优于它，则该配置是\textbf{Pareto 最优的}（Pareto optimal）。Pareto 最优即在不牺牲其他参与者福利、也不使用更多资源的前提下，无法改进任一参与者的福利（某人福利的提高，必定以损害他人的利益为前提）。给定偏好后，通常只有可行配置的一个子集是最优的。
\end{tcolorbox}

Pareto 最优配置也叫做\textbf{有效}（efficient）配置。若证券市场使参与者达到 Pareto 最优配置，则称该市场为\textbf{有效市场}（efficient market）。
\end{tcolorbox}


\begin{tcolorbox}[example,title=例题：帕累托最优]
    \begin{question}
        \textbf{例 2.5}\quad 考虑一个两期、1 期有两种状态 $a$、$b$ 的经济。经济总禀赋：
        \begin{center}
        \WJEventTree[left={$100$},stem]{400,25}
        \end{center}
        有两个参与者，效用函数均为 $u(c_0,c_{1a},c_{1b})=\sqrt{c_0}+\tfrac12(\sqrt{c_{1a}}+\sqrt{c_{1b}})$（$c=\begin{bmatrix}c_0 \\ c_{1a} \\ c_{1b}\end{bmatrix}$）。树中每个节点括号内为 $(c_{1,\cdot},\,c_{2,\cdot})$：该节点上参与者 1、2 的消费量（不是整段消费计划列向量）。考虑三种配置：
        \begin{center}
        \begin{tabular}{@{}c@{\quad}c@{\quad}c@{}}
            \textbf{配置 A} & \textbf{配置 B} & \textbf{配置 C} \\[0.35em]
            \WJEventTree[left={$(100,0)$},stem]{ {(0,400)}, {(0,25)} } &
            \WJEventTree[left={$(50,50)$},stem]{ {(200,200)}, {(\tfrac{25}{2},\tfrac{25}{2})} } &
            \WJEventTree[left={$(100,0)$},stem]{ {(400,0)}, {(25,0)} }
        \end{tabular}
        \end{center}
        配置 A 中参与者 1 在 0 期消费全部、参与者 2 在 1 期消费全部；B 为两参与者在各期、各状态均分；C 中参与者 1 始终消费全部。哪种配置是 Pareto 最优的？（完整消费计划 $c_1$、$c_2$ 见解答。）
    \end{question}
    \begin{solution}
        与前文及习题 2.2 相同，参与者 $i$ 的\textbf{消费计划}是 $\mathrm{R}^{3}_{+}$ 中的列向量（两状态记为 $a,b$）：
        \[
            c_i=\begin{bmatrix}c_{i,0} \\ c_{i,1a} \\ c_{i,1b}\end{bmatrix},
            \qquad U_i(c_i)=\sqrt{c_{i,0}}+\tfrac12\bigl(\sqrt{c_{i,1a}}+\sqrt{c_{i,1b}}\bigr).
        \]
        题中事件树每个节点上的 $(\cdot,\cdot)$ 只是该节点上\textbf{参与者 1、2 的消费量}；整段计划仍用上面的 $c_1$、$c_2$ 表示。

        \textbf{可行性:}三个配置合起来在每一日期—状态下满足 $c_{1,\cdot}+c_{2,\cdot}$ 等于总禀赋（0 期 $100$，状态 $a$ 为 $400$，状态 $b$ 为 $25$），故 A、B、C 均为可行配置。

        \textbf{各配置的效用：}
        \begin{itemize}
            \item \textbf{配置 A：}
            \[
                c_1=\begin{bmatrix}100 \\ 0 \\ 0\end{bmatrix},\qquad
                c_2=\begin{bmatrix}0 \\ 400 \\ 25\end{bmatrix}.
            \]
            \[
                \begin{aligned}
                    U_1(c_1) &= \sqrt{100}+\tfrac12(0+0)=10,\qquad 
                    U_2(c_2) &= 0+\tfrac12(\sqrt{400}+\sqrt{25})
                    =\tfrac12(20+5)=12.5.
                \end{aligned}
            \]
            效用对为 $(U_1,U_2)=(10,\,12.5)$。

            \item \textbf{配置 B}（均分）：$c_1=c_2=
            \begin{bmatrix}50 \\ 200 \\ \tfrac{25}{2}\end{bmatrix}$。
            \[
                U_1(c_1)=U_2(c_2)=\sqrt{50}+\tfrac12\Bigl(\sqrt{200}+\sqrt{\tfrac{25}{2}}\Bigr)
                \approx 7.071+8.839=15.91.
            \]
            效用对为 $(15.91,\,15.91)$。

            \item \textbf{配置 C：}
            \[
                c_1=\begin{bmatrix}100 \\ 400 \\ 25\end{bmatrix},\qquad
                c_2=\begin{bmatrix}0 \\ 0 \\ 0\end{bmatrix}.
            \]
            \[
                \begin{aligned}
                    U_1(c_1) = \sqrt{100}+\tfrac12(\sqrt{400}+\sqrt{25})=10+12.5=22.5, \quad
                    U_2(c_2) = 0.
                \end{aligned}
            \]
            效用对为 $(22.5,\,0)$。
        \end{itemize}

        \textbf{Pareto 占优与 Pareto 最优。}
        \begin{itemize}
            \item 比较 A 与 B：$U_1(B)=15.91>10=U_1(A)$ 且 $U_2(B)=15.91>12.5=U_2(A)$，两人都严格更好，故\textbf{配置 B Pareto 占优于配置 A}。因此 A \textbf{不是} Pareto 最优。

            \item 比较 B 与 C：$U_1(C)=22.5>U_1(B)$，但 $U_2(C)=0<U_2(B)$，有人变差，故 C 不 Pareto 占优 B，B 也不 Pareto 占优 C。

            \item \textbf{配置 B 是 Pareto 最优的}：不存在其他可行配置在不让参与者 2 变差的前提下再提高参与者 1 的效用（反之亦然）；B 与 C 只是在两人之间做了不同的福利权衡。

            \item \textbf{配置 C 是 Pareto 最优的}：若要提高参与者 2 的效用（必须给 2 正的消费），只能从参与者 1 已占用的 $100$、$400$、$25$ 中让出，从而必使 $U_1$ 下降，无法在不损害 1 的情况下改进 2。
        \end{itemize}

        
    \end{solution}

    \begin{tcolorbox}[warning,title=注意]
        \textbf{结论：}A 非 Pareto 最优；\textbf{B 与 C 都是} Pareto 最优配置（同一总禀赋下不同的有效分配）。可以看出：
        \begin{itemize}
            \item 在同一可行集（同一总禀赋、满足 $\sum_k c_k=\sum_k e_k$）里，Pareto 最优配置往往不止一个。
            \item 更能看到，帕累托最优仅考虑了“效率”，但没考虑“公平”，例如配置 C 中参与者 2 的效用为0，但仍属于帕累托最优。
        \end{itemize}
        \end{tcolorbox}
\end{tcolorbox}


\begin{tcolorbox}[knowledge,title=2.7 本章小结]
本章完成了在市场经济中通过资源配置分析所使用的基本框架的构建。资源配置是通过证券市场完成的：交易证券的价格引导参与者通过交易达到他们希望的配置，而交易又最终决定了证券的市场价格。证券价格在引导个体参与者和平衡整个经济方面都起着关键作用。因此，我们关注证券价格如何决定，这也称为\textbf{资产定价}（asset pricing）。

\begin{enumerate}
    \item 从金融的角度看，一个经济的结构包括三个方面：（1）它所处的\textbf{自然环境}；（2）\textbf{经济的参与者}；（3）\textbf{金融市场}。

    \item \textbf{外部环境}包括两个要素：\textbf{时间}和\textbf{风险}。为刻画时间，假设有 0 和 1 两个时期；为刻画风险，假设 1 期有 $\Omega$ 个可能状态，用概率测度 $P$ 描述。经济中只有一种不可储存的商品。

    \item 经济中有 $K$ 个\textbf{参与者}。每个参与者的经济特征包括其经济资源和经济需求。具体假设为：
    \begin{enumerate}[label=(\alph*)]
        \item 每个参与者 $k$ 具有禀赋 $e_k\in\mathrm{R}^{1+\Omega}_{+}$；
        \item 每个参与者对未来状态发生的可能性拥有相同的信息，完全由概率测度 $P$ 描述；
        \item 尽管参与者给定其禀赋和信息后可能拥有生产技术，这里先假设他们不拥有任何生产技术；
        \item 每个参与者在 $C=\mathrm{R}^{1+\Omega}_{+}$ 上具有偏好，且满足不满足性、连续性和凸性公理；
        \item 给定这些偏好关系，可用效用函数来描述。
    \end{enumerate}

    \item \textbf{金融市场}由一组证券构成：
    \begin{enumerate}[label=(\alph*)]
        \item 它们的支付只依赖于未来实际的经济状态；
        \item 它们的支付是外生给定的，且不受经济参与者行为的影响；
        \item 参与者知道且一致认同证券支付的概率分布；
        \item 市场结构为 $X$，且市场中无摩擦。
    \end{enumerate}

    \item \textbf{市场均衡}必须满足两个条件：
    \begin{enumerate}[label=(\alph*)]
        \item 每个参与者各自最优化其证券需求量；
        \item \textbf{市场出清}，即交易证券的总需求等于总供给，此时商品市场也是出清的。
    \end{enumerate}

    \item 如果配置 $\{c_k,\,\forall k\}$ 是可行的，且不存在另外的可行配置 Pareto 占优于它，则该配置是 \textbf{Pareto 最优的}。如果证券市场允许参与者达到 Pareto 最优配置，就称它为\textbf{有效市场}。
\end{enumerate}
\end{tcolorbox}


\begin{problem}
\problemno{2.1} $U(c)$ 和 $V(c)$ 是两个效用函数，$c\in\mathrm{R}_{+}^{n}$，且 $V(x)=f(U(x))$，其中 $f(\cdot)$ 是一正单调函数。证明这两个效用函数表示了相同的偏好。
\end{problem}
\begin{solution}
假设 $U(c)$ 表示的偏好关系为 $\succeq$，那么 $\forall\, c_{1},\,c_{2}\in\mathrm{R}_{+}^{N}$ 有
\[
    U(c_{1})\ge U(c_{2})\Leftrightarrow c_{1}\succeq c_{2}
\]
而 $f(\cdot)$ 是正单调函数，因而
\[
    V(c_{1})=f(U(c_{1}))\ge f(U(c_{2}))=V(c_{2})\Leftrightarrow U(c_{1})\ge U(c_{2})
\]
因此 $V(c_{1})\ge V(c_{2})\Leftrightarrow c_{1}\succeq c_{2}$，即 $V(c)$ 表示的偏好也是 $\succeq$。
\end{solution}

\begin{problem}
\problemno{2.2*} 在 1 期，经济有两个可能状态 $a$ 和 $b$，它们的发生概率相等：
\begin{center}
\WJEventTree{a,b}
\end{center}
考虑定义在消费计划 $c=\begin{bmatrix}c_{0} \\ c_{1a} \\ c_{1b}\end{bmatrix}$ 上的效用函数：
\[
    \begin{aligned}
        U(c) &= \log c_{0}+\tfrac12(\log c_{1a}+\log c_{1b})\\
        U(c) &= \tfrac{1}{1-\gamma}c_{0}^{1-\gamma}+\tfrac12\Bigl(\tfrac{1}{1-\gamma}c_{1a}^{1-\gamma}+\tfrac{1}{1-\gamma}c_{1b}^{1-\gamma}\Bigr)\\
        U(c) &= -e^{-ac_{0}}-\tfrac12\bigl(e^{-ac_{1a}}+e^{-ac_{1b}}\bigr)
    \end{aligned}
\]
证明它们满足：不满足性、连续性和凸性。
\end{problem}
\begin{solution}
在这里只证明第一个效用函数，可以类似地证明第二、第三个效用函数的性质。

\begin{enumerate}[label=(\alph*)]
    \item 先证明不满足性。假设 $c\ge c'$，那么有
    \[
        c_{0}\ge c_{0}',\quad c_{1a}\ge c_{1a}',\quad c_{1b}\ge c_{1b}'
    \]
    而 $\log(\cdot)$ 是单调增函数，因此有
    \[
        \log(c_{0})\ge\log(c_{0}'),\quad \log(c_{1a})\ge\log(c_{1a}'),\quad \log(c_{1b})\ge\log(c_{1b}')
    \]
    因而 $U(c)\ge U(c')$，即 $c\succeq c'$。

    \item 现在证明连续性。令 $\{c^{(n)}\}_{1}^{\infty}$ 为 $\mathrm{R}^{3}$ 中一个序列，且 $\lim_{n\to\infty}c^{(n)}=c$。对于 $\forall\,\varepsilon>0$，$\exists\,\delta$，当 $|c_{i}'-c_{i}|\le\delta$ 时我们有 $|\log(c_{i}')-\log(c_{i})|<\tfrac{\varepsilon}{3}$，$i=0,1a,1b$；对于 $\delta$，$\exists N$ 使得当 $n\ge N$ 时，
    \[
        \lVert c^{(n)}-c\rVert=\sqrt{(c_{0}-c_{0}^{(n)})^{2}+(c_{1a}-c_{1a}^{(n)})^{2}+(c_{1b}-c_{1b}^{(n)})^{2}}<\delta
    \]
    因而 $|U(c^{(n)})-U(c)|<\varepsilon$，故 $\lim_{c^{(n)}\to c}U(c^{(n)})=U(c)$。

    \item 最后证明凸性。假设 $U(C)>U(C')$，那么
    \[
        \log(c_{0})+\tfrac12\log(c_{1a})+\tfrac12\log(c_{1b})>\log(c_{0}')+\tfrac12\log(c_{1a}')+\tfrac12\log(c_{1b}')
    \]
    对于 $\forall\,\alpha\in(0,1)$，
    \[
        \begin{aligned}
            U(\alpha C'+(1-\alpha)C)
            &=\log(\alpha c_{0}'+(1-\alpha)c_{0})+\tfrac12\log(\alpha c_{1a}'+(1-\alpha)c_{1a})+\tfrac12\log(\alpha c_{1b}'+(1-\alpha)c_{1b})\\
            &\ge\alpha\bigl(\log(c_{0}')+\tfrac12\log(c_{1a}')+\tfrac12\log(c_{1b}')\bigr)+(1-\alpha)\bigl(\log(c_{0})+\tfrac12\log(c_{1a})+\tfrac12\log(c_{1b})\bigr)\\
            &=\alpha U(C')+(1-\alpha)U(C)>U(C')
        \end{aligned}
    \]
    故凸性成立。
\end{enumerate}
\end{solution}

\begin{problem}
\problemno{2.3} $U(c)=c-\tfrac12 a c^{2}$ 是一可能的效用函数，其中 $c\in\mathrm{R}_{+}$，$a$ 是非负的系数。$U(c)$ 具有不满足性吗？如果不是，那么 $a$ 取什么值和/或 $c$ 在什么范围内时 $U(c)$ 具有不满足性？
\end{problem}
\begin{solution}
不一定。比如当 $a=1$ 时，$U(\tfrac12)=\tfrac38<U(1)=\tfrac12>U(3)=-1.5$。$U(c)$ 不具有不满足性。当 $a=0$ 时，$U(c)=c$ 具有不满足性；当 $a>0$ 时，当 $c\in[0,\tfrac{1}{a}]$ 时 $U(c)$ 具有不满足性。
\end{solution}

\begin{problem}
\problemno{2.4} 考虑一个经济，它在 1 期有三个可能状态：$a$，$b$ 和 $c$：
\begin{center}
\WJEventTree{a,b,c}
\end{center}
证券市场包括证券 1 和 2，它们具有如下的支付向量：$X_{1}=\begin{bmatrix}1 \\ 1 \\ 1\end{bmatrix}$ 以及 $X_{2}=\begin{bmatrix}1 \\ 2 \\ 3\end{bmatrix}$。它们的价格分别为 $S_{1}$ 和 $S_{2}$。
\begin{enumerate}[label=(\alph*)]
    \item 描述这个经济的支付空间。
    \item 写出这个经济的市场结构矩阵 $X$。
    \item 考虑含有 $\theta_{1}$ 单位的证券 1 和 $\theta_{2}$ 单位的证券 2 的组合。写出这个组合的支付向量。这个组合的价格是多少？
    \item 假设这个市场中总共有 $K$ 个参与者。每个参与者的禀赋是 1 单位的证券 1 和 2 单位的证券 2。这时的市场组合是什么？市场组合的支付向量是什么？市场组合的总价值是多少？
    \item 写出市场化支付的集合。
    \item 如果市场不允许卖空。市场化支付的集合是什么？
    \item 现在引入新的证券 3，它的支付向量为 $X_{3}=\begin{bmatrix}0 \\ 0 \\ 1\end{bmatrix}$。写出新的市场结构矩阵。在这个市场结构下，市场化支付集合是什么？
\end{enumerate}
\end{problem}
\begin{solution}
三个状态 $a,b,c$ 对应 $\Omega=3$。记
\[
    X_1=\begin{bmatrix}1 \\ 1 \\ 1\end{bmatrix},\quad
    X_2=\begin{bmatrix}1 \\ 2 \\ 3\end{bmatrix},\quad
    \theta=\begin{bmatrix}\theta_1 \\ \theta_2\end{bmatrix}.
\]
组合在 1 期的支付向量为 $X_1\theta_1+X_2\theta_2$，也可写成矩阵乘积 $X\theta$，其中 $X=[X_1,\,X_2]$。

\begin{enumerate}[label=(\alph*)]
    \item \textbf{支付空间}：所有可能在三个状态下实现的支付向量的集合，即 $\mathrm{R}^{\Omega}=\mathrm{R}^{3}$（每个分量对应状态 $a,b,c$ 下的支付）。

    \item \textbf{市场结构矩阵}：两列分别为两只证券的支付向量，
    \[
        X=\begin{bmatrix}1&1\\1&2\\1&3\end{bmatrix}=[X_1,\,X_2].
    \]

    \item \textbf{组合支付与价格：}持有 $(\theta_1,\theta_2)$ 时，
    \[
        X\theta=\theta_1 X_1+\theta_2 X_2=
        \begin{bmatrix}\theta_1+\theta_2 \\ \theta_1+2\theta_2 \\ \theta_1+3\theta_2\end{bmatrix}.
    \]
    0 期购买成本（组合价格）为 $\theta_1 S_1+\theta_2 S_2$。

    \item \textbf{市场组合：}$K$ 个参与者每人禀赋为 $1$ 单位证券 1、$2$ 单位证券 2，加总得市场组合
    \[
        \theta_M=\begin{bmatrix}K \\ 2K\end{bmatrix}.
    \]
    其支付向量为
    \[
        X\theta_M=K X_1+2K X_2=
        \begin{bmatrix}3K \\ 5K \\ 7K\end{bmatrix}.
    \]
    总价值为 $K S_1+2K S_2$。

    \item \textbf{市场化支付（允许卖空）：}能由证券组合复制的 1 期支付集合为
    \[
        M=\bigl\{Y\in\mathrm{R}^{3}: Y=X\theta,\;\theta_1,\theta_2\in\mathrm{R}\bigr\},
    \]
    即 $X_1,X_2$ 张成的二维子空间（只有两只证券，维度 $\le 2$）。

    \item \textbf{禁止卖空：}仅允许 $\theta_1,\theta_2\ge 0$，为 $X_1,X_2$ 的\textbf{非负}线性组合：
    \[
        M_{+}=\bigl\{Y\in\mathrm{R}^{3}: Y=\theta_1 X_1+\theta_2 X_2,\;\theta_1,\theta_2\in\mathrm{R}_{+}\bigr\}.
    \]

    \item \textbf{加入证券 3：}$X_3=\begin{bmatrix}0 \\ 0 \\ 1\end{bmatrix}$ 仅在状态 $c$ 支付 $1$，市场结构变为
    \[
        X=\begin{bmatrix}1&1&0\\1&2&0\\1&3&1\end{bmatrix}=[X_1,\,X_2,\,X_3].
    \]
    三列线性无关，可表示任意 $Y\in\mathrm{R}^{3}$，故（允许卖空时）
    \[
        M=\bigl\{Y\in\mathrm{R}^{3}: Y=\theta_1 X_1+\theta_2 X_2+\theta_3 X_3,\;\theta_1,\theta_2,\theta_3\in\mathrm{R}\bigr\}=\mathrm{R}^{3}.
    \]
\end{enumerate}
\end{solution}

\begin{problem}
\problemno{2.5} 在练习 2.4 中定义的只存在证券 1 和 2 的经济中。考虑一个禀赋为 $\theta_{1}$ 单位的证券 1 和 $\theta_{2}$ 单位的证券 2 的参与者。写出他的预算集。
\end{problem}
\begin{solution}
参与者的预算集是 $\{C\in\mathrm{R}_{+}^{3}:C=\alpha_{1}X_{1}+\alpha_{2}X_{2}$，其中 $\alpha_{1}S_{1}+\alpha_{2}S_{2}\le\theta_{1}S_{1}+\theta_{2}S_{2}\}$。
\end{solution}

\begin{problem}
\problemno{2.6} 在上面的练习中引入练习 2.4 中定义的证券 3，它的价格为 $S_{3}$。这时，参与者的预算集是什么（他在证券 3 上的禀赋为 0）？证明由证券 1、2、3 构成的预算集包含仅由证券 1、2 构成的预算集。
\end{problem}
\begin{solution}
此时参与者的预算集变成了 $\{C\in\mathrm{R}_{+}^{3}:C=\alpha_{1}X_{1}+\alpha_{2}X_{2}+\alpha_{3}X_{3}$，其中 $\alpha_{1}S_{1}+\alpha_{2}S_{2}+\alpha_{3}S_{3}\le\theta_{1}S_{1}+\theta_{2}S_{2}+\theta_{3}S_{3}\}$。
\end{solution}

\begin{problem}
\problemno{2.7*} 考虑一个在 1 期只有一个可能状态的经济。（在这种情况下不存在不确定性。）参与者 1 的 0 期禀赋为 100 而 1 期禀赋为 1，即他的禀赋向量为 $\begin{bmatrix}100 \\ 1\end{bmatrix}$。他的偏好可以表示成如下形式：
\[
    U(c_{0},c_{1})=\log c_{0}+\rho\log c_{1}.
\]
系数 $\rho$ 为反映参与者在当前消费和未来消费之间相对偏好的参数。有一只证券，它的 0 期价格为 1、1 期支付为 $1+r_{F}$。这里，$r_{F}$ 是利率。
\begin{enumerate}[label=(\alph*)]
    \item 如果这个参与者不能在市场上进行交易，那么他的消费计划以及相应的效用 $U^{a}$ 是什么？
    \item 现在假设他可以在市场上进行交易。
    \begin{itemize}
        \item 他的预算集是什么？以当前消费为单位，他的总财富 $w$ 是多少？
        \item 写出参与者的优化问题。令 $c_{0}$ 为参与者的当前（即 0 期）最优消费，$s$ 为最优储蓄以及 $U^{b}$ 为在最优策略下得到的效用。求解他的最优消费/储蓄选择以及相应的效用。把 $U^{b}$ 表示成财富 $w$、利率 $r_{F}$ 和偏好系数 $\rho$ 的函数。
        \item 讨论参与者的最优选择如何依赖于利率 $r_{F}$ 和偏好系数 $\rho$。给出解释。
    \end{itemize}
    \item 证明 $U^{b}\ge U^{a}$。
    \item 令 $g$ 为参与者由于能够在证券市场上交易而获得的益处。它的定义为
    \[
        U^{b}(w-g)=U^{a}.
    \]
    计算 $g$。讨论 $g$ 如何依赖于 $\rho$？$g$ 如何依赖于 $r_{F}$？给出解释。
\end{enumerate}
\end{problem}
\begin{solution}
禀赋 $e_0=100$、$e_1=1$。无风险证券：0 期价格 $1$，1 期支付 $1+r_F$。记 $S$ 为 0 期\textbf{储蓄}（买入该证券的金额；$S<0$ 表示借贷）。

\begin{enumerate}[label=(\alph*)]
    \item \textbf{不能交易（自给自足）：}只能消费禀赋，
    \[
        c_0=100,\qquad c_1=1,\qquad
        U^{a}=\log(100)+\rho\log(1)=\log 100.
    \]

    \item \textbf{可以交易。}
    \begin{itemize}
        \item \textbf{预算集与财富 $w$。}买卖证券等价于在两期之间转移消费：0 期少消费 $S$ 用于储蓄，1 期增加 $S(1+r_F)$。故
        \[
            c_0=100-S,\qquad c_1=1+S(1+r_F),\qquad S\in\mathrm{R}.
        \]
        以 0 期消费为计价单位，把 1 期禀赋折现到 0 期（用市场利率 $r_F$）：
        \[
            w=100+\frac{1}{1+r_F}.
        \]
        这是参与者按市价出售全部禀赋后可在 0 期支配的\textbf{总财富}（现值）。任意满足上述两式的 $(c_0,c_1)$ 都在预算集内。

        \item \textbf{优化问题。}在预算集上最大化对数效用：
        \[
            \max_{S\in\mathrm{R}}\ \log(100-S)+\rho\log\bigl(1+S(1+r_F)\bigr).
        \]
        一阶条件（边际效用替代率 = 跨期相对价格）：
        \[
            \underbrace{\frac{1}{100-S}}_{\text{0 期边际效用}}
            =\rho\underbrace{\frac{1+r_F}{1+S(1+r_F)}}_{\text{1 期边际效用}\times\text{价格比}}.
        \]
        整理得 $\bigl[1+S(1+r_F)\bigr]=\rho(1+r_F)(100-S)$，即
        \[
            1+S(1+r_F)(1+\rho)=100\rho(1+r_F),
            \qquad
            S=\frac{100\rho(1+r_F)-1}{(1+\rho)(1+r_F)}.
        \]
        代回预算约束，并用 $w=100+1/(1+r_F)=\dfrac{100(1+r_F)+1}{1+r_F}$，得
        \[
            c_0=\frac{100(1+r_F)+1}{(1+\rho)(1+r_F)}=\frac{w}{1+\rho},\qquad
            c_1=\frac{\rho\bigl[100(1+r_F)+1\bigr]}{1+\rho}=\frac{\rho(1+r_F)}{1+\rho}\,w.
        \]
        即现期消费占财富 $\dfrac{1}{1+\rho}$，1 期消费占 $\dfrac{\rho(1+r_F)}{1+\rho}$（对数效用下的固定支出份额）。最优效用
        \[
            U^{b}=\log\Bigl(\frac{w}{1+\rho}\Bigr)+\rho\log\Bigl(\frac{\rho(1+r_F)}{1+\rho}\,w\Bigr).
        \]

        \item \textbf{比较静态。}$c_0=\dfrac{w}{1+\rho}$ 随 $r_F$、$\rho$ 上升而下降（$w$ 对 $r_F$ 略降，$1+\rho$ 对 $\rho$ 升）；$c_1=\dfrac{\rho(1+r_F)}{1+\rho}w$ 随 $r_F$、$\rho$ 上升而上升。直觉：$r_F$ 升高提高储蓄回报，$\rho$ 升高加重未来消费权重，都会促使\textbf{少消费于 0 期、多留于 1 期}。
    \end{itemize}

    \item \textbf{$U^{b}\ge U^{a}$。}不能交易时等价于强制 $S=0$，得到 $U^{a}$；能交易时在更大可行集上求最大值，最优值 $U^{b}$ 至少不低于 $S=0$，故 $U^{b}\ge U^{a}$。

    \item \textbf{交易益处 $g$。}定义 $U^{b}(w-g)=U^{a}$：$w-g$ 是「与交易后最优效用相同、但用一条 0 期财富数字表示」的\textbf{确定性等价财富}；$g=w-(w-g)$ 即为因能交易而多出来的等价 0 期资源。

    将 $U^{a}=\log 100$ 与上式 $U^{b}$ 联立，记 $w-g=\hat w$，则
    \[
        (1+\rho)\log\hat w+\rho\log\bigl(\rho(1+r_F)\bigr)-(1+\rho)\log(1+\rho)=\log 100.
    \]
    取指数并整理：
    \[
        \hat w^{1+\rho}\,\frac{\bigl[\rho(1+r_F)\bigr]^{\rho}}{(1+\rho)^{1+\rho}}=100
        \quad\Rightarrow\quad
        \hat w=\rho^{-\frac{\rho}{1+\rho}}(1+r_F)^{-\frac{\rho}{1+\rho}}
        100^{\frac{1}{1+\rho}}(1+\rho).
    \]
    故
    \[
        g=w-\hat w
        =w-\rho^{-\frac{\rho}{1+\rho}}(1+r_F)^{-\frac{\rho}{1+\rho}}
        100^{\frac{1}{1+\rho}}(1+\rho).
    \]
    $g$ 随 $\rho$、$r_F$ 增加而增加：二者都提高「把消费从 0 期转到 1 期」的动机，能在市场上做跨期配置带来的等价收益更大。
\end{enumerate}
\end{solution}

% !TeX root = ../../金融经济学习题讲义.tex
% 源 PDF：references/金融经济学/史上最全最完整的金融经济学答案(王江版).pdf，页 8–12（印刷页 7–11）

\zhchapter{Arrow-Debreu 经济}


本章运用第 2 章的框架分析\textcolor{red}{\textbf{参与者如何选择投资组合、证券价格如何影响这些选择以及市场如何决定这些价格}}。这一简单的市场结构也是理解资源配置效率的基准。

\begin{tcolorbox}[knowledge,title=3.1 Arrow-Debreu 证券市场]
下面引入一类最简单的证券：它们只在 1 期支付，且支付是已实现状态的函数。

\begin{tcolorbox}[knowledge]
\textbf{定义 3.1}\quad 状态 $\omega$ \textbf{或有要求权}（state-$\omega$ contingent claim）就是当状态为 $\omega$ 时支付为 $1$、而在其他状态下支付为 $0$ 的证券。
\end{tcolorbox}

因此，状态 $\omega$ 或有要求权具有如下支付分布（$x_{\omega'}$ 为状态 $\omega'$ 下的支付）：
\[
    x_{\omega'}=
    \begin{cases}
        1, & \omega'=\omega,\\
        0, & \omega'\neq\omega,
    \end{cases}
    \qquad\text{即}\qquad
    x(\omega')=\mathbf{1}_{\omega}(\omega'),
\]
其中 $\mathbf{1}_{\omega}(\cdot)$ 为示性函数：当 $\omega'=\omega$ 时取 $1$，否则取 $0$。

\begin{tcolorbox}[example]
\textbf{例 3.1}\quad 在例 2.1 的 Lucas 树经济中，1 期有两个可能状态，因而有两个或有要求权，其支付向量分别为
\[
    \text{状态 $a$ 或有要求权：}\begin{bmatrix}1 \\ 0\end{bmatrix},\qquad
    \text{状态 $b$ 或有要求权：}\begin{bmatrix}0 \\ 1\end{bmatrix}.
\]
\end{tcolorbox}

对每个状态都可以定义一个相应的或有要求权，共有 $\Omega$ 个。这些或有要求权称为\textbf{Arrow-Debreu 证券}；由它们的完全集合构成的证券市场称为\textbf{Arrow-Debreu 证券市场}。若按状态排列证券，支付矩阵（市场结构）为单位阵：
\[
    X^{A\text{-}D}=
    \begin{bmatrix}
        1 & \cdots & 0 & \cdots & 0 \\
        \vdots & \ddots & \vdots & \ddots & \vdots \\
        0 & \cdots & 1 & \cdots & 0 \\
        \vdots & \ddots & \vdots & \ddots & \vdots \\
        0 & \cdots & 0 & \cdots & 1
    \end{bmatrix}
    =I.
\]
\end{tcolorbox}

\begin{tcolorbox}[knowledge,title=3.2 状态价格]
记 $\phi_{\omega}$ 为 0 期购买、在状态 $\omega$ 才支付 $1$ 单位消费品的状态 $\omega$ 或有要求权的价格（以 0 期消费品计量）。\textcolor{red}{\textbf{因为状态 $\omega$ 或有要求权在 1 期只有当状态为 $\omega$ 时才支付 $1$ 单位的消费品}}，$\phi_{\omega}$ 也称为\textbf{状态价格}（state price）。所有状态价格组成向量
\[
    \phi=\begin{bmatrix}\phi_1 \\ \vdots \\ \phi_{\omega} \\ \vdots \\ \phi_{\Omega}\end{bmatrix}.
\]
在无套利（没有免费的午餐）下，状态价格严格为正：
\[
    \phi\gg 0.
\]
\end{tcolorbox}

\begin{tcolorbox}[knowledge,title=3.3 市场的完全性]
在 Arrow-Debreu 市场中，1 期的消费计划可写成 $x=\begin{bmatrix}c_{11} \\ \vdots \\ c_{1\omega} \\ \vdots \\ c_{1\Omega}\end{bmatrix}$。持有状态证券组合 $\theta=\begin{bmatrix}c_{11} \\ \vdots \\ c_{1\omega} \\ \vdots \\ c_{1\Omega}\end{bmatrix}$ 即可复制该消费计划：
\[
    \begin{aligned}
        \underbrace{X\theta}_{\substack{\text{组合在 1 期各状态下的总支付}\\\text{（一般：任意证券组合的收益）}}}
        =\underbrace{X^{A\text{-}D}\begin{bmatrix}c_{11} \\ \vdots \\ c_{1\Omega}\end{bmatrix}}_{\substack{\theta\text{：状态 }\omega\text{ 阿罗证券持有 }c_{1\omega}\text{ 份}\\\text{研究阿罗市场，此处 }X=X^{A\text{-}D}}}
        =\underbrace{\begin{bmatrix}c_{11} \\ \vdots \\ c_{1\Omega}\end{bmatrix}}_{X^{A\text{-}D}=I\text{，故状态 }\omega\text{ 下收入 }c_{1\omega}}
        =\underbrace{x}_{\text{欲实现的 1 期消费计划}}.
    \end{aligned}
\]
该组合在 0 期的成本为
\[
    \underbrace{\phi^{\mathsf T}\theta}_{\text{复制该计划在 0 期支付的总成本}}
    =\underbrace{\sum_{\omega\in\Omega}\phi_{\omega}c_{1\omega}}_{\text{各状态消费 }\times\text{ 对应状态价格}}.
\]

是否\textbf{完全}，取决于市场能否以\textbf{有限成本}为任意\textbf{有限消费计划}\textbf{融资}。

\begin{tcolorbox}[knowledge]
\textbf{定义 3.2}\quad 如果市场中的任一有限消费计划都可以通过有限成本的可交易证券的组合来融资，那么我们就称这个证券市场是\textbf{完全的}（complete）。即，任给消费计划 $c=\begin{bmatrix}c_{11} \\ \vdots \\ c_{1\Omega}\end{bmatrix}$，存在证券组合 $\theta=\begin{bmatrix}\theta_1 \\ \vdots \\ \theta_{\Omega}\end{bmatrix}$ 使得 $X\theta=c$ 有解。
\end{tcolorbox}



\begin{tcolorbox}[example]
\textbf{例 3.2}\quad 在 Lucas 树经济中，若市场上只有一只无风险债券，其支付为 $\begin{bmatrix}1 \\ 1\end{bmatrix}$。参与者若想要消费计划 $\begin{bmatrix}2 \\ 1\end{bmatrix}$（在状态 $a$ 消费 $2$、在状态 $b$ 消费 $1$），无法仅靠该债券实现，市场\textbf{不完备}。
\end{tcolorbox}


\begin{enumerate}
    \item 在完全市场中，任意消费计划 $x=[c_{11},\dots,c_{1\Omega}]$ 都是状态证券的线性组合；投资者可通过买卖状态证券在不同状态间调整消费。
    \item 完全市场中，消费者有最高的灵活度，可以实现最优的风险配置。
    \item 所有的完备市场都是等价的，对一个完备市场的结论可以应用到其他完备市场。
    \item 完备市场处理起来更轻松。
    \item \textbf{金融创新}的价值之一，正是通过增加可交易证券使市场更完全，从而更好地实现风险分担与效率。
\end{enumerate}

\end{tcolorbox}

\begin{tcolorbox}[knowledge,title=3.4 参与者的优化]
考虑单个参与者的效用优化问题。在 Arrow-Debreu 市场中，设参与者禀赋为 $e$、效用函数为 $U$。若 1 期禀赋以状态证券组合 $\bar{\theta}=\begin{bmatrix}e_{11} \\ \vdots \\ e_{1\Omega}\end{bmatrix}$ 持有（初始持有），则 1 期禀赋的支付为
\[
    X^{A\text{-}D}\bar{\theta}=\bar{\theta}=e_1=
    \begin{bmatrix}e_{11} \\ \vdots \\ e_{1\Omega}\end{bmatrix}.
\]

上式$\bar{\theta}$也称\textbf{复制组合}（replication portfolio）。复制消费计划 $c_1$ 的组合 $\bar{\theta}$ 在 0 期的市场价值为 $\phi^{\mathsf T}\bar{\theta}=\begin{bmatrix}
    \phi_1 & \cdots & \phi_{\omega} & \cdots & \phi_{\Omega}
\end{bmatrix}\begin{bmatrix}e_{11} \\ \vdots \\ e_{1\Omega}\end{bmatrix}$。因此，参与者的\textbf{金融财富}（简称\textbf{财富}）$w$ 是其禀赋在 0 期的市场价值（初始禀赋+持有状态证券在第零期的价值）：
\[
    w=e_0+\phi^{\mathsf T}e_1.
    \tag{3.2}
\]

参与者可以把自己的初始禀赋换为等值的现金，从而通过消费（现在的消费$c_0$+购买证券$\theta$获得未来的消费 $c_1=\bar{\theta}$），实现自己效用的最优化。优化过程要满足\textbf{预算约束}：消费总成本不超过财富，
\[
    c_0+\phi^{\mathsf T}c_1\le w=e_0+\phi^{\mathsf T}e_1.
    \tag{3.3}
\]

记扩展价格向量 $\hat{\phi}=\begin{bmatrix}1 \\ \phi\end{bmatrix}$（0 期消费品价格为 $1$，即计价物），则预算约束可写为 $\hat{\phi}^{\mathsf T}c\le\hat{\phi}^{\mathsf T}e$，或 $\hat{\phi}^{\mathsf T}(e-c)\ge 0$。

参与者的优化问题为（在预算内最大化效用，且消费非负）
\begin{align}
    \max_{c\in C}\ &U(c) \tag{3.4}\\
    \mathrm{s.t.}\quad &\hat{\phi}^{\mathsf T}(e-c)\ge 0, \notag\\
    &c\ge 0. \notag
\end{align}

\begin{tcolorbox}[knowledge]
\textbf{定理 3.1}\quad 设 $C=\mathrm{R}^{1+\Omega}_{+}$，$U(\cdot)$ 在 $C$ 上连续，则问题 (3.4) 有解。
\end{tcolorbox}

\textbf{不满足公理}意味着边际效用恒为正。记 $\partial_0 U$ 为 0 期消费的边际效用，$\partial_{\omega}U$ 为状态 $\omega$ 下 1 期消费的边际效用，则
\[
    \partial_0 U>0,\qquad \partial_{\omega}U>0,\quad\forall\omega\in\Omega,
    \tag{3.5}
\]
边际效用向量为 $DU=\begin{bmatrix}\partial_0 U \\ \partial_1 U \\ \vdots \\ \partial_{\Omega}U\end{bmatrix}$，不满足性即 $DU\gg 0$。

\begin{tcolorbox}[knowledge]
\textbf{定理 3.2}\quad 设 $DU\gg 0$。则 (3.4) （参与者的效用优化问题，在预算内最大化效用，且消费非负）的解满足：存在 $\lambda>0$、$\mu_i\ge 0$（$i=0,1,\dots,\Omega$）使得
\begin{align}
    \partial_i U&=\lambda\phi_i+\mu_i,\quad i=0,1,\dots,\Omega\quad(\phi_0=1),\tag{3.6a}\\
    \hat{\phi}^{\mathsf T}(e-c)&=0,\tag{3.6b}\\
    \mu_i c_i&=0,\quad i=0,1,\dots,\Omega.\tag{3.6c}
\end{align}


定理 3.2 说的是：在问题 (3.4)（预算约束 $+$ 消费非负）下，若不满足性成立（$DU\gg 0$），则最优消费满足 KKT 条件 (3.6a)–(3.6c)（不等式约束下的一阶必要条件；偏好与可行集凸时亦为充分条件）。含义简述如下：
\begin{itemize}
    \item \textbf{(3.6a)} \textbf{边际效用定价：}$\partial_i U=\lambda\phi_i+\mu_i$。内点（$c_i>0$）时 $\mu_i=0$，边际效用与扩展价格成比例；若某期/某状态消费为 $0$，则 $\mu_i\ge 0$ 表示非负约束在起作用。
    \item \textbf{(3.6b)} \textbf{预算紧绑：}按 $\hat{\phi}$ 计价时花光财富，$\hat{\phi}^{\mathsf T}c=\hat{\phi}^{\mathsf T}e$（在 $DU\gg 0$ 下不会故意少花）。
    \item \textbf{(3.6c)} \textbf{互补松弛：}$\mu_i c_i=0$。要么该分量消费为正、对应乘子为零；要么消费为零、允许 $\mu_i>0$（角点解）。
\end{itemize}
\end{tcolorbox}

对问题 (3.4) 构造拉格朗日函数（$\lambda$ 为预算约束 $\hat{\phi}^{\mathsf T}(e-c)\ge 0$ 的乘子，$\mu_i$ 为 $c_i\ge 0$ 的乘子）：
\[
    \mathcal{L}
    =U(c)+\lambda\,\hat{\phi}^{\mathsf T}(e-c)-\sum_{i=0}^{\Omega}\mu_i c_i.
\]
对 $c_i$ 求偏导并令其为零，即得 (3.6a)；与 (3.6b)、(3.6c) 联立即为定理 3.2。对参与者 $k$ 将 $U,c,e$ 换为 $U_k,c_k,e_k$ 即可。

对参与者 $k$，\textbf{内点解}（$c_k\gg 0$，故 $\mu_i=0$）时由 (3.6a) 可写成 $DU_k(c_k)=\lambda_k\hat{\phi}$，即
\begin{align}
    \underbrace{\frac{\partial U_k(c_k)}{\partial c_k(0)}=\lambda_k}_{\substack{\text{对第 0 期消费求偏导}\\(\phi_0=1)}},
    &\quad
    \underbrace{\frac{\partial U_k(c_k)}{\partial c_k(\omega)}=\lambda_k\phi_{\omega}}_{\substack{\text{对第 1 期各状态消费求偏导}}},\quad
    \forall\omega\in\Omega. \tag{3.7}
\end{align}

等价地，\textbf{跨期边际替代率}等于状态价格：
\[
\begin{aligned}
&\underbrace{\frac{\partial_{\omega}U_k(c_k)}{\partial_0 U_k(c_k)}}_{\text{跨期边际替代率}}=\phi_{\omega}\\ 
\Rightarrow \quad &\underbrace{\partial_{\omega}U_k(c_k)}_{\text{每一单位第0期消费的边际效用}}=\underbrace{\frac{1}{\phi_{\omega}}\partial_0 U_k(c_k)}_{\text{用第 0 期1单位效用，能换来 $1/\phi_{\omega}$ 单位第1期的边际效用}}.
\end{aligned}
\tag{3.8}
\]

这说明，你愿意用多少 0 期消费去换 1 期状态 $\omega$ 下的 1 单位消费，在最优时正好等于市场价格 $\phi_{\omega}$。对于参与者来说，在达到最优的时候，在不同时期和状态间转移消费是无差异的。之，如果这种转移是有差异的，那就说明参与者还可以通过这种转移进一步增加取得的效用，说明原本并不是最优。

\end{tcolorbox}

\begin{tcolorbox}[example]
    \begin{problem}
        \problemno{例 3.3} 考虑“Lucas 树经济”，1 期有两个概率相等的状态 $a$ 和 $b$。假设 1 期的两个可能状态的状态价格分别为 $\phi_{a}$ 和 $\phi_{b}$。考虑一个参与者，他的禀赋为 $(e_{0};e_{1a};e_{1b})$。他具有如下对数形式的效用函数：
        \[
            U(c_{0},c_{1a},c_{1b})=\log c_{0}+\tfrac12(\log c_{1a}+\log c_{1b})
        \]
        他的最优消费/组合选择是什么？
        \end{problem}
        \begin{solution}
        给定状态价格和他的禀赋，他的总财富是 $w=e_{0}+\phi_{a}e_{1a}+\phi_{b}e_{1b}$。他的最优化问题是
        \[
            \begin{aligned}
                \max_{\{c_{0},c_{1a},c_{1b}\}}\ &\log c_{0}+\tfrac12(\log c_{1a}+\log c_{1b})\\
                \mathrm{s.t.}\quad &w-(c_{0}+\phi_{a}c_{1a}+\phi_{b}c_{1b})=0\\
                &c_{0},c_{1a},c_{1b}\ge 0
            \end{aligned}
        \]
        拉格朗日函数为
        \[
            \mathcal{L}=\log c_{0}+\tfrac12(\log c_{1a}+\log c_{1b})+\lambda\bigl(w-(c_{0}+\phi_{a}c_{1a}+\phi_{b}c_{1b})\bigr)-\mu_{0}c_{0}-\mu_{a}c_{1a}-\mu_{b}c_{1b}
        \]
        其一阶条件为
        \[
            \begin{aligned}
                \frac{1}{c_{0}}&=\lambda+\mu_{0}\\
                \tfrac12\bigl(1/c_{1a}\bigr)&=\lambda\phi_{a}+\mu_{a}\\
                \tfrac12\bigl(1/c_{1b}\bigr)&=\lambda\phi_{b}+\mu_{b}\\
                c_{0}+\phi_{a}c_{1a}+\phi_{b}c_{1b}&=w\\
                \mu_{i}c_{i}&=0,\quad i=0,a,b
            \end{aligned}
        \]
        对数效用在 $c=0$ 处边际效用为无穷大，因而最优消费在每一时期每一状态都严格为正（当然，假设 $\phi\gg 0$，即所有状态价格严格为正）。也就是说，最优消费总是在 $\mathrm{R}^{1+\Omega}$ 的内部而不是边界上，因而非负性约束从不起限制作用。在这种情况下，我们可以在一阶条件中去掉这些约束（以及对应的乘子）而直接求解最优。因此，$\mu_{i}=0$（$i=0,a,b$）。对于 $c$ 我们立即得到如下解：
        \[
            c_{0}=\frac{1}{\lambda},\qquad
            c_{1a}=\frac{1}{\lambda}\frac{1}{2\phi_{a}},\qquad
            c_{1b}=\frac{1}{\lambda}\frac{1}{2\phi_{b}}
        \]
        把 $c$ 的解代入预算约束，我们可以得到 $\lambda$ 的解：
        \[
            \frac{1}{\lambda}+\frac{1}{2\lambda}+\frac{1}{2\lambda}=w\qquad\Rightarrow\qquad\lambda=\frac{2}{w}
        \]
        最后，我们有
        \[
            c_{0}=\frac12 w,\qquad
            c_{1a}=\frac14\frac{w}{\phi_{a}},\qquad
            c_{1b}=\frac14\frac{w}{\phi_{b}}
        \]
        可以看出，参与者把一半财富用作现在的消费，把另外一半财富作为未来的消费。某一状态下的消费与对应的状态价格负相关。这个结果很直观。状态价格高的状态下的消费更昂贵。结果，参与者在这些状态下选择较低的消费。
        \end{solution}
经济含义：一半财富用于 0 期消费，一半用于 1 期；状态价格越高，该状态下消费越低。\textbf{拉格朗日推导见后文「例 3.3」。}

\end{tcolorbox}

\begin{tcolorbox}[knowledge,title=3.5 市场均衡]
给定状态价格，Arrow-Debreu 经济中的\textbf{均衡}是一组状态价格 $\{\phi_{\omega},\,\omega\in\Omega\}$，使得：
\begin{enumerate}
    \item \textbf{给定价格，每个参与者最优化：}对参与者 $k$，
    \[
        c_k(e_k,\phi)=\operatorname*{Argmax}\,U_k(c_k)
    \]
    满足 $c_{k,0}+\phi^{\mathsf T}c_{k,1}=e_{k,0}+\phi^{\mathsf T}e_{k,1}$ 且 $c_k\ge 0$；
    \item \textbf{证券市场出清}（商品市场出清）：
    \[
        \sum_{k=1}^{K}c_{k,0}(e_k,\phi)=\sum_{k=1}^{K}e_{k,0},\qquad
        \sum_{k=1}^{K}c_{k,1}(e_k,\phi)=\sum_{k=1}^{K}e_{k,1}.
    \]
\end{enumerate}
由 Walras 法则，若 $m-1$ 个市场出清，则第 $m$ 个市场也出清。

\begin{tcolorbox}[knowledge,title=瓦尔拉斯均衡]
\textbf{瓦尔拉斯均衡}（Walrasian equilibrium，竞争均衡）指一组价格与配置，使得：
\begin{enumerate}
    \item 每个参与者在给定价格下最大化效用（或厂商最大化利润），且为\textbf{价格接受者}；
    \item 每个市场（每种商品）上\textbf{总需求等于总供给}（市场出清）。
\end{enumerate}
在本章中，商品是 $1+\Omega$ 种\textbf{或有商品}：0 期消费品（计价物，价格 $1$）与 1 期各状态 $\omega$ 下的消费品（价格 $\phi_\omega$，即状态 $\omega$ 阿罗证券的价格）。上文两条均衡条件正是瓦尔拉斯均衡在 Arrow-Debreu 证券市场经济中的写法，故也称\textbf{Arrow-Debreu 均衡}。第 2.5 节中的 Walras 法则说明：在标准预算下，$1+\Omega$ 个商品市场里若 $1+\Omega-1$ 个已出清，则剩余一个自动出清；实物总禀赋 $\sum_k e_k$ 仅被重新分配，不会凭空增减。
\end{tcolorbox}

\begin{tcolorbox}[knowledge]
\textbf{定理 3.3}\quad 对定义 2.5 （两期、$\Omega$ 状态、单一不可储存商品、$K$ 名同质信息参与者、禀赋与满足公理 1–3 的偏好（不满足性、连续性、凸性）、无摩擦证券市场 $X$）中的经济，Arrow-Debreu 证券市场存在均衡。
\end{tcolorbox}

均衡时所有人面对相同价格，因而在任一状态 $\omega$ 上\textcolor{red}{\textbf{相对边际效用都相等}}：
\[
    \frac{\partial_{\omega}U_k(c_k)}{\partial_0 U_k(c_k)}
    =\phi_{\omega}
    =\frac{\partial_{\omega}U_{k'}(c_{k'})}{\partial_0 U_{k'}(c_{k'})},
    \quad\forall k,k',\,\omega\in\Omega.
    \tag{3.9}
\]
$\phi_{\omega}$ 可理解为：为在状态 $\omega$ 多获得 1 单位 1 期消费，0 期须放弃的消费量——市场上关于跨期、跨状态资源配置的\textbf{共识价格}。

\end{tcolorbox}

\begin{tcolorbox}[example]
    \begin{problem}
        \problemno{例 3.4} 考虑一个经济，在 1 期有两个概率相等的状态 $a$ 和 $b$。经济中有参与者 1 和 2，他们具有的禀赋分别为：
        \begin{center}
        \WJEventTree[left={$e_{1}{:}100$}]{0,0}
        \qquad
        \WJEventTree[left={$e_{2}{:}0$}]{200,50}
        \end{center}
        两个参与者都具有如下形式的对数效用函数：
        \[
            U(c)=\log c_{0}+\tfrac12(\log c_{a}+\log c_{b})
        \]
        \end{problem}
        \begin{solution}
        在市场上存在一组完全的状态或有证券可以交易。因为有两个状态，因而只有两个状态或有证券。
        
        现在我们开始分析这个经济的均衡。从给定交易证券价格下参与者的最优化问题开始。记 $\phi=\begin{bmatrix}\phi_{a} \\ \phi_{b}\end{bmatrix}$ 为状态价格（向量），即两个状态或有证券的价格。我们可以定义每个参与者的财富为 $w=\hat{\phi}^{\top}e$，这里 $\hat{\phi}=\begin{bmatrix}1 \\ \phi\end{bmatrix}$，而 $e$ 是他的禀赋。这时，最优化问题变成了：
        \[
            \begin{aligned}
                \max_{c}\ &\log c_{0}+\tfrac12(\log c_{a}+\log c_{b})\\
                \mathrm{s.t.}\quad &c_{0}+\phi_{a}c_{a}+\phi_{b}c_{b}=w
            \end{aligned}
        \]
        由例 3.3，我们知道解为
        \[
            c_{k,0}=\tfrac12 w_{k},\qquad
            c_{k,a}=\tfrac14\frac{w_{k}}{\phi_{a}},\qquad
            c_{k,b}=\tfrac14\frac{w_{k}}{\phi_{b}},\qquad k=1,2
        \]
        在状态价格 $\phi_{a},\phi_{b}$ 下，参与者按禀赋市值计的财富为 $w_{k}=e_{k,0}+\phi_{a}e_{k,a}+\phi_{b}e_{k,b}$（0 期消费品价格为 $1$）。代入禀赋 $(100;0,0)$ 与 $(0;200,50)$：
        \[
            w_{1}=100+0\cdot\phi_{a}+0\cdot\phi_{b}=100,\qquad
            w_{2}=0+200\phi_{a}+50\phi_{b}=200\phi_{a}+50\phi_{b}.
        \]

        \textbf{0 期商品市场出清}（两人 0 期消费之和等于经济 0 期总禀赋 $100$）。由例 3.3，$c_{k,0}=\tfrac12 w_{k}$，故
        \[
            c_{1,0}+c_{2,0}=\tfrac12 w_{1}+\tfrac12 w_{2}=\tfrac12(w_{1}+w_{2})=100
            \quad\Longrightarrow\quad w_{1}+w_{2}=200.
        \]
        把 $w_{1},w_{2}$ 代入 $w_{1}+w_{2}=200$：
        \[
            100+(200\phi_{a}+50\phi_{b})=200.
        \]
        移项得 $200\phi_{a}+50\phi_{b}=100$；两边同除以 $50$，得到与状态价格有关的第一个方程
        \[
            4\phi_{a}+\phi_{b}=2.
        \]

        \textbf{1 期·状态 $a$ 商品市场出清}（状态 $a$ 下总消费 $200$）。由 $c_{k,a}=\tfrac14 w_{k}/\phi_{a}$，
        \[
            c_{1,a}+c_{2,a}=\tfrac{w_{1}}{4\phi_{a}}+\tfrac{w_{2}}{4\phi_{a}}
            =\frac{w_{1}+w_{2}}{4\phi_{a}}=200
            \quad\Longrightarrow\quad w_{1}+w_{2}=800\phi_{a}.
        \]
        与 $w_{1}+w_{2}=200$ 联立得 $\phi_{a}=\dfrac14$，从而 $\phi_{b}=1$。状态 $b$ 出清 $\dfrac{w_{1}+w_{2}}{4\phi_{b}}=50$ 自动成立。此时 $w_{1}=w_{2}=100$，最优消费为
        \[
            \begin{aligned}
                c_{1,0}=c_{2,0}&=50\\
                c_{1,a}=c_{2,a}&=100\\
                c_{1,b}=c_{2,b}&=25
            \end{aligned}
        \]
        \end{solution}
\end{tcolorbox}



\begin{tcolorbox}[knowledge,title=3.5 市场均衡（续）]
可行配置满足 $\sum_{k=1}^{K}c_k=\sum_{k=1}^{K}e_k$。对均衡配置，将各参与者预算约束在 $\hat{\phi}$ 下加总，得到
\[
    \sum_{k=1}^{K}\hat{\phi}^{\mathsf T}c_k=\sum_{k=1}^{K}\hat{\phi}^{\mathsf T}e_k.
    \tag{3.10}
\]

\begin{tcolorbox}[knowledge]
\textbf{定理 3.4}\quad 对存在 Arrow-Debreu 证券市场的经济（定义 2.5），均衡达到的资源配置是\textbf{Pareto 最优的}。
\end{tcolorbox}

定理 3.4 即\textbf{福利经济学第一定理}：完全、无摩擦的证券市场使参与者达到\textcolor{red}{\textbf{资源在不同时期和不同状态下的 Pareto 最优配置}}，为证券市场的资源配置效率提供了理论基础（亚当·斯密“看不见的手”的现代表述）。

\begin{tcolorbox}[knowledge,title=福利经济学两大定理]
\textbf{第一定理}（定理 3.4）：在完全、无摩擦的竞争市场中，任一均衡配置都是 Pareto 最优的。

\textbf{第二定理}（福利经济学第二定理）：在消费集凸、偏好凸（满足公理 3）且市场完全（如 Arrow-Debreu 市场）等条件下，\textbf{任一} Pareto 最优配置都可以成为某个竞争均衡的配置——只需在交易开始前对禀赋作适当的\textbf{一次性总付转移}（lump-sum redistribution），使参与者的初始禀赋变为 $e_k'$，再令他们在状态价格 $\phi$ 下自由交易即可达到该配置。

第二定理的含义：在一定条件下，\textbf{效率}问题与\textbf{分配}问题可以分开处理——先确定有效的总资源配置，再通过不扭曲价格信号的转移支付来实现合意的福利分配。（现实中转移支付往往受信息、政治与再分配成本限制，故第二定理更多是规范基准。）
\end{tcolorbox}


\end{tcolorbox}

\begin{tcolorbox}[knowledge,title=本章小结]
\begin{enumerate}
    \item 状态 $\omega$ 或有要求权在状态 $\omega$ 支付 $1$、其他状态支付 $0$；由全部状态或有证券构成的市场为 Arrow-Debreu 证券市场。
    \item 任一有限消费计划可用有限成本证券组合融资时，市场完全；Arrow-Debreu 市场是完全市场。
    \item 均衡存在，且各参与者最优化、市场出清；均衡时相对边际效用（跨期边际替代率）相等。
    \item \textbf{福利经济学第一定理：}Arrow-Debreu 均衡配置 Pareto 最优。
    \item \textbf{福利经济学第二定理：}在凸性与完全市场下，任一 Pareto 最优配置可通过一次性总付转移与市场交易实现。
\end{enumerate}
\end{tcolorbox}







\begin{problem}
\problemno{3.1*} 考虑如下经济，在 1 期有两个可能状态 $a$ 和 $b$：
\begin{center}
\WJEventTree{a,b}
\end{center}
\begin{enumerate}[label=(\alph*)]
    \item 描述所有 Arrow-Debreu 证券的支付向量。记这些证券的价格向量为 $\phi$。
    \item 考虑一个拥有如下禀赋的参与者：
    \begin{center}
    \WJEventTree[left={$0$}]{2,1}
    \end{center}
    把他的禀赋表示成 Arrow-Debreu 证券的组合。
    \item 计算他的金融财富。写出他的预算集。
    \item 假设参与者的效用函数如下：
    \[
        U(c_{0},c_{1a},c_{1b})=-e^{-c_{0}}-\tfrac12\bigl(e^{-c_{1a}}+e^{-c_{1b}}\bigr).
    \]
    不考虑消费的非负约束，写出他的优化问题。求解他的最优消费选择。
    \item 讨论他的消费如何依赖于 Arrow-Debreu 证券的价格向量 $\phi$。
    \item 证明在某些价格下，他（在某些时期/状态下）的消费可能是负的。
\end{enumerate}
\end{problem}
\begin{solution}
\begin{enumerate}[label=(\alph*)]
    \item Arrow-Debreu 证券的支付向量是 $X_{a}=\begin{bmatrix}1 \\ 0\end{bmatrix}$，$X_{b}=\begin{bmatrix}0 \\ 1\end{bmatrix}$；
    \item $2X_{a}+X_{b}$；
    \item 参与者的金融财富是 $w=2\phi_{a}+\phi_{b}$，他的预算集是 $\{c\in\mathrm{R}_{+}^{3}:c_{0}+\phi_{a}c_{1a}+\phi_{b}c_{1b}=w\}$；
    \item 由于不考虑非负约束，参与者的优化问题就变成了
    \[
        \begin{aligned}
            \max_{\{c_{0},c_{1a},c_{1b}\}}\ &\bigl\{-e^{-c_{0}}-\tfrac12\bigl(e^{-c_{1a}}+e^{-c_{1b}}\bigr)\bigr\}\\
            \mathrm{s.t.}\quad &c_{0}+\phi_{a}c_{1a}+\phi_{b}c_{1b}=w=2\phi_{a}+\phi_{b}
        \end{aligned}
    \]
    拉格朗日函数为
    \[
        \mathcal{L}=-e^{-c_{0}}-\tfrac12\bigl(e^{-c_{1a}}+e^{-c_{1b}}\bigr)+\lambda\bigl(w-c_{0}-\phi_{a}c_{1a}-\phi_{b}c_{1b}\bigr)
    \]
    对消费求导，一阶条件为
    \[
        e^{-c_{0}}=\lambda,\qquad
        \tfrac12 e^{-c_{1a}}=\lambda\phi_{a},\qquad
        \tfrac12 e^{-c_{1b}}=\lambda\phi_{b}
    \]
    从而
    \[
        c_{0}=-\log\lambda,\qquad
        c_{1a}=-\log(2\lambda\phi_{a}),\qquad
        c_{1b}=-\log(2\lambda\phi_{b})
    \]
    代入预算约束：
    \[
        -\log\lambda-\phi_{a}\log(2\lambda\phi_{a})-\phi_{b}\log(2\lambda\phi_{b})=w
    \]
    即 $-(1+\phi_{a}+\phi_{b})\log\lambda-\phi_{a}\log(2\phi_{a})-\phi_{b}\log(2\phi_{b})=w$。记 $S=1+\phi_{a}+\phi_{b}$，则
    \[
        \log\lambda=-\frac{w+\phi_{a}\log(2\phi_{a})+\phi_{b}\log(2\phi_{b})}{S}
    \]
    从而 $c_{0}^{*}=-\log\lambda$，且 $c_{1a}^{*}=-\log\lambda-\log(2\phi_{a})$、$c_{1b}^{*}=-\log\lambda-\log(2\phi_{b})$。代入 $w=2\phi_{a}+\phi_{b}$ 得到
    \[
        \begin{aligned}
            c_{0}^{*}&=\frac{2\phi_{a}+\phi_{b}+\phi_{a}\log(2\phi_{a})+\phi_{b}\log(2\phi_{b})}{1+\phi_{a}+\phi_{b}}\\
            c_{1a}^{*}&=\frac{2\phi_{a}+\phi_{b}+\phi_{b}\log(2\phi_{b})-\log(2\phi_{a})-\phi_{b}\log(2\phi_{a})}{1+\phi_{a}+\phi_{b}}\\
            c_{1b}^{*}&=\frac{2\phi_{a}+\phi_{b}+\phi_{a}\log(2\phi_{a})-\log(2\phi_{b})-\phi_{a}\log(2\phi_{b})}{1+\phi_{a}+\phi_{b}}
        \end{aligned}
    \]
    \item 状态价格变化对消费有财富效应和价格效应。当 $\phi_{a}$ 上升时，$w=2\phi_{a}+\phi_{b}$ 增加，对三种消费都有正的财富效应；同时 $a$ 更贵，对 $c_{1a}^{*}$ 有负的替代效应，对 $c_{0}^{*}$、$c_{1b}^{*}$ 有正的替代效应。因此 $c_{0}^{*}$、$c_{1b}^{*}$ 上升，而 $c_{1a}^{*}$ 的净效应不确定；
    \item 当 $\phi_{a}=\phi_{b}=0.1$ 时，$w=0.3$，$S=1.2$，$\log(0.2)<0$，故
    \[
        c_{0}^{*}=\frac{0.3+2\times 0.1\log(0.2)}{1.2}\approx -0.0182<0
    \]
\end{enumerate}
\end{solution}

\begin{problem}
\problemno{3.2} 考虑一个在 1 期只有一个可能状态的经济。（在这种情况下不存在不确定性。）参与者 1 的 0 期禀赋为 100 而 1 期禀赋为 1，即他的禀赋向量为 $\begin{bmatrix}100 \\ 1\end{bmatrix}$。他的偏好可以表示成如下形式：
\[
    U(c_{0},c_{1})=\log c_{0}+\rho\log c_{1}.
\]
有一只可交易证券，它的 0 期价格为 1、1 期支付为 $1+r_{F}$。这里，$r_{F}$ 是利率。
\begin{enumerate}[label=(\alph*)]
    \item 假设利率 $r_{F}$ 是给定的，导出参与者对证券的需求。
    \item 假设参与者 1 是经济中的唯一参与者。描述市场出清条件。
    \item 求解均衡利率。
    \item 均衡利率如何依赖于偏好参数 $\rho$？解释所得到的结论。
\end{enumerate}
\end{problem}
\begin{solution}
\begin{enumerate}[label=(\alph*)]
    \item 先不考虑消费的非负性，求解参与者的证券需求。此时参与者的优化问题变成了
    \[
        \max_{\theta}\bigl\{\log(100-\theta)+\rho\log\bigl(1+\theta(1+r_{F})\bigr)\bigr\}
    \]
    这与练习 2.7 是同一规划。一阶条件
    \[
        \frac{1}{100-\theta}=\rho\frac{1+r_{F}}{1+\theta(1+r_{F})}
    \]
    解得证券需求
    \[
        \theta^{*}=\frac{100\rho(1+r_{F})-1}{(1+\rho)(1+r_{F})}
    \]
    最优消费为
    \[
        c_{0}^{*}=\frac{100(1+r_{F})+1}{(1+\rho)(1+r_{F})},\quad
        c_{1}^{*}=\frac{\rho\bigl(100(1+r_{F})+1\bigr)}{1+\rho}
    \]
    由于最优消费都是非负的，因此上面的求解过程没有问题。
    \item 由于市场上只有一个参与者，因而市场出清条件是 $\theta^{*}=0$；
    \item 由上面的市场出清条件，我们可以求得均衡利率为 $r_{F}=\dfrac{1}{100\rho}-1$；
    \item $r_{F}$ 随着偏好系数的增加而减小。这是因为偏好系数增大，未来消费带来的边际效用相对增加，从而人们对能将财富在 0 期、1 期之间进行转换的证券需求增加，证券价格上升，利率下降。
\end{enumerate}
\end{solution}

\begin{problem}
\problemno{3.3*} 继续考虑练习 3.2 中定义的经济。现在假设还有另外一个参与者 2，他的效用函数和参与者 1 的一样但他的禀赋为 $\begin{bmatrix}1 \\ 100\end{bmatrix}$。
\begin{enumerate}[label=(\alph*)]
    \item 求解参与者 2 的最优消费/储蓄选择。
    \item 比较参与者 1 和 2 的最优消费/储蓄选择。
    \item 求解有两个参与者时的市场均衡。
    \item 计算均衡利率。它与前一情形（经济中只有参与者 1）下的均衡利率有何不同？
    \item 证明均衡配置是 Pareto 最优的。
\end{enumerate}
\end{problem}
\begin{solution}
\begin{enumerate}[label=(\alph*)]
    \item 先不考虑消费非负的约束，参与者 2 的优化问题就变成了
    \[
        \max_{\theta}\bigl\{\log(1-\theta)+\rho\log\bigl(100+\theta(1+r_{F})\bigr)\bigr\}
    \]
    一阶条件 $\dfrac{1}{1-\theta}=\rho\dfrac{1+r_{F}}{100+\theta(1+r_{F})}$ 给出证券需求
    \[
        \theta^{*}=\frac{\rho(1+r_{F})-100}{(1+\rho)(1+r_{F})}
    \]
    最优消费为
    \[
        c_{0}^{*}=\frac{(1+r_{F})+100}{(1+\rho)(1+r_{F})},\quad
        c_{1}^{*}=\frac{\rho\bigl((1+r_{F})+100\bigr)}{1+\rho}
    \]
    由于最优消费都是非负的，因此上面的求解过程没有问题。
    \item 两人的最优消费/储蓄选择形式上是完全一样的，差别在于两人的禀赋不一样；
    \item 此时的市场出清条件是 $\theta_{1}^{*}+\theta_{2}^{*}=0$，即
    \[
        \frac{100\rho(1+r_{F})-1}{(1+\rho)(1+r_{F})}+\frac{\rho(1+r_{F})-100}{(1+\rho)(1+r_{F})}=0
    \]
    分子之和为零：$101\rho(1+r_{F})=101$，故 $1+r_{F}=1/\rho$；
    \item 均衡利率是 $r_{F}=1/\rho-1$。只有参与者 1 时总禀赋为 $(100,1)$，均衡利率为 $1/(100\rho)-1$；现在总禀赋为 $(101,101)$，利率由跨期边际替代率 $\rho$ 单独决定，不再被极端的当期/未来禀赋扭曲；
    \item 边际效用之比为
    \[
        \frac{\partial U/\partial c_{1}}{\partial U/\partial c_{0}}=\rho\frac{c_{0}}{c_{1}}
    \]
    由两人的最优消费都有 $c_{0}^{*}/c_{1}^{*}=1/[\rho(1+r_{F})]$，因而两人的
    \[
        \rho\frac{c_{0}^{*}}{c_{1}^{*}}=\frac{1}{1+r_{F}}
    \]
    相等，均衡配置是 Pareto 最优的。
\end{enumerate}
\end{solution}

\begin{problem}
\problemno{3.4*} 在 1 期，经济有三个等可能的状态 $a$、$b$ 和 $c$。在证券市场上交易的 Arrow-Debreu 证券集合是完全的。经济中有两个参与者，参与者 1 和 2。他们的禀赋如下：
\begin{center}
\WJEventTree[left={$e_{1}{:}100$}]{0,0,0}
\qquad
\WJEventTree[left={$e_{2}{:}0$}]{200,100,50}
\end{center}
他们的偏好相同，形式如下：
\[
    U(c)=-\frac{1}{c_{0}}-\frac13\Bigl(\frac{1}{c_{1a}}+\frac{1}{c_{1b}}+\frac{1}{c_{1c}}\Bigr).
\]
\begin{enumerate}[label=(\alph*)]
    \item 求解每一参与者的最优消费/投资选择。
    \item 求解均衡状态价格。
    \item 证明对两个参与者的消费都有禀赋越高、消费越高的性质。
    \item 证明每一参与者的消费在总消费中所占的份额 $c_{k,1\omega}/e_{1\omega}$ 在所有状态下都是一样的，其中 $k=1,2$，而 $e_{1\omega}$ 是在状态 $\omega=a,b,c$ 下的 1 期总消费/禀赋。
\end{enumerate}
\end{problem}
\begin{solution}
\begin{enumerate}[label=(\alph*)]
    \item 求解最优消费选择时，先不考虑消费非负的约束。对于参与者 1，它的优化问题是：
    \[
        \begin{aligned}
            \max_{\{c_{0},c_{1}\}}\ &\Bigl\{-\frac{1}{c_{0}}-\frac13\Bigl(\frac{1}{c_{1a}}+\frac{1}{c_{1b}}+\frac{1}{c_{1c}}\Bigr)\Bigr\}\\
            \mathrm{s.t.}\quad &c_{0}+\phi_{a}c_{1a}+\phi_{b}c_{1b}+\phi_{c}c_{1c}=100
        \end{aligned}
    \]
    拉格朗日函数为
    \[
        \mathcal{L}=-\frac{1}{c_{0}}-\frac13\Bigl(\frac{1}{c_{1a}}+\frac{1}{c_{1b}}+\frac{1}{c_{1c}}\Bigr)+\lambda\bigl(100-c_{0}-\phi_{a}c_{1a}-\phi_{b}c_{1b}-\phi_{c}c_{1c}\bigr)
    \]
    一阶条件是边际效用与状态价格成比例：
    \[
        \frac{1}{c_{0}^{2}}=\lambda,\qquad
        \frac{1}{3c_{1\omega}^{2}}=\lambda\phi_{\omega},\quad\omega=a,b,c
    \]
    因而 $c_{1\omega}=c_{0}/\sqrt{3\phi_{\omega}}$。代入预算约束
    \[
        c_{0}\Bigl(1+\sqrt{\frac{\phi_{a}}{3}}+\sqrt{\frac{\phi_{b}}{3}}+\sqrt{\frac{\phi_{c}}{3}}\Bigr)=100
    \]
    得到参与者 1 的最优消费选择为
    \[
        \begin{aligned}
            c_{0}^{*}&=\frac{100}{1+\sqrt{\phi_{a}/3}+\sqrt{\phi_{b}/3}+\sqrt{\phi_{c}/3}}\\
            c_{1a}^{*}&=\frac{100}{\sqrt{3\phi_{a}}\bigl(1+\sqrt{\phi_{a}/3}+\sqrt{\phi_{b}/3}+\sqrt{\phi_{c}/3}\bigr)}\\
            c_{1b}^{*}&=\frac{100}{\sqrt{3\phi_{b}}\bigl(1+\sqrt{\phi_{a}/3}+\sqrt{\phi_{b}/3}+\sqrt{\phi_{c}/3}\bigr)}\\
            c_{1c}^{*}&=\frac{100}{\sqrt{3\phi_{c}}\bigl(1+\sqrt{\phi_{a}/3}+\sqrt{\phi_{b}/3}+\sqrt{\phi_{c}/3}\bigr)}
        \end{aligned}
    \]
    参与者 2 的财富为 $w_{2}=200\phi_{a}+100\phi_{b}+50\phi_{c}$，同一组一阶条件给出
    \[
        c_{0}^{*}=\frac{w_{2}}{1+\sqrt{\phi_{a}/3}+\sqrt{\phi_{b}/3}+\sqrt{\phi_{c}/3}},\qquad
        c_{1\omega}^{*}=\frac{c_{0}^{*}}{\sqrt{3\phi_{\omega}}}
    \]
    即
    \[
        \begin{aligned}
            c_{0}^{*}&=\frac{200\phi_{a}+100\phi_{b}+50\phi_{c}}{1+\sqrt{\phi_{a}/3}+\sqrt{\phi_{b}/3}+\sqrt{\phi_{c}/3}}\\
            c_{1a}^{*}&=\frac{200\phi_{a}+100\phi_{b}+50\phi_{c}}{\sqrt{3\phi_{a}}\bigl(1+\sqrt{\phi_{a}/3}+\sqrt{\phi_{b}/3}+\sqrt{\phi_{c}/3}\bigr)}\\
            c_{1b}^{*}&=\frac{200\phi_{a}+100\phi_{b}+50\phi_{c}}{\sqrt{3\phi_{b}}\bigl(1+\sqrt{\phi_{a}/3}+\sqrt{\phi_{b}/3}+\sqrt{\phi_{c}/3}\bigr)}\\
            c_{1c}^{*}&=\frac{200\phi_{a}+100\phi_{b}+50\phi_{c}}{\sqrt{3\phi_{c}}\bigl(1+\sqrt{\phi_{a}/3}+\sqrt{\phi_{b}/3}+\sqrt{\phi_{c}/3}\bigr)}
        \end{aligned}
    \]
    因为最优消费都是非负的，因而上面的求解过程是没有问题的；
    \item 市场出清条件是
    \[
        \begin{aligned}
            &\frac{100}{\sqrt{3\phi_{a}}\bigl(1+\sqrt{\phi_{a}/3}+\sqrt{\phi_{b}/3}+\sqrt{\phi_{c}/3}\bigr)}+\frac{200\phi_{a}+100\phi_{b}+50\phi_{c}}{\sqrt{3\phi_{a}}\bigl(1+\sqrt{\phi_{a}/3}+\sqrt{\phi_{b}/3}+\sqrt{\phi_{c}/3}\bigr)}=200\\
            &\frac{100}{\sqrt{3\phi_{b}}\bigl(1+\sqrt{\phi_{a}/3}+\sqrt{\phi_{b}/3}+\sqrt{\phi_{c}/3}\bigr)}+\frac{200\phi_{a}+100\phi_{b}+50\phi_{c}}{\sqrt{3\phi_{b}}\bigl(1+\sqrt{\phi_{a}/3}+\sqrt{\phi_{b}/3}+\sqrt{\phi_{c}/3}\bigr)}=100\\
            &\frac{100}{\sqrt{3\phi_{c}}\bigl(1+\sqrt{\phi_{a}/3}+\sqrt{\phi_{b}/3}+\sqrt{\phi_{c}/3}\bigr)}+\frac{200\phi_{a}+100\phi_{b}+50\phi_{c}}{\sqrt{3\phi_{c}}\bigl(1+\sqrt{\phi_{a}/3}+\sqrt{\phi_{b}/3}+\sqrt{\phi_{c}/3}\bigr)}=50
        \end{aligned}
    \]
    由 $c_{1\omega}=c_{0}/\sqrt{3\phi_{\omega}}$，0 期总禀赋为 100，故 $\sum_{k}c_{k,0}=100$。状态 $a$ 的出清化为
    \[
        \frac{\sum_{k}c_{k,0}}{\sqrt{3\phi_{a}}}=\frac{100}{\sqrt{3\phi_{a}}}=200
    \]
    即 $\sqrt{3\phi_{a}}=\dfrac12$，从而 $\phi_{a}=\dfrac{1}{12}$。同理 $\dfrac{100}{\sqrt{3\phi_{b}}}=100$ 给出 $\phi_{b}=\dfrac13$，$\dfrac{100}{\sqrt{3\phi_{c}}}=50$ 给出 $\phi_{c}=\dfrac43$。
    \item $\phi_{a}<\phi_{b}<\phi_{c}$，对于参与者 1、2 都有 $c_{1a}^{*}>c_{1b}^{*}>c_{1c}^{*}$，即总禀赋越高、消费越高；
    \item 因为对于消费者 1、2 都有
    \[
        \begin{aligned}
            \frac{c_{1a}^{*}}{c_{1b}^{*}}&=\sqrt{\frac{\phi_{b}}{\phi_{a}}}=2=\frac{200}{100}\\
            \frac{c_{1a}^{*}}{c_{1c}^{*}}&=\sqrt{\frac{\phi_{c}}{\phi_{a}}}=4=\frac{200}{50}
        \end{aligned}
    \]
    因此，每一参与者的消费在总消费中所占的份额在所有状态下都是相等的。
\end{enumerate}
\end{solution}

% !TeX root = ../../金融经济学习题讲义.tex
% 源 PDF：references/金融经济学/史上最全最完整的金融经济学答案(王江版).pdf，页 14–18（印刷页 13–17）

\zhchapter{套利和资产定价}

第 3 章在 Arrow--Debreu 市场结构下讨论均衡与效率；本章回到第 2 章的一般支付矩阵 $X$，讨论\textbf{证券价格的基本性质}、\textbf{定价原理}、\textbf{无套利}与\textbf{资产定价基本定理}。

\begin{tcolorbox}[knowledge,title=4.1 一般市场结构]
多数可交易证券在多个状态下都有正支付，称为\textbf{复合证券}（相对只在单一状态付 1 的阿罗证券而言）。设 1 期有 $\Omega$ 个状态，市场上有 $N$ 只证券（$n=1,\dots,N$），第 $n$ 只证券的支付向量为 $X_n=\begin{bmatrix}x_{1,n} \\ \vdots \\ x_{\Omega,n}\end{bmatrix}$，\textbf{市场结构矩阵}为（与第 2.3 节一致，行为状态、列为证券）：
\[
    X=[X_{\cdot,1},\,\cdots,\,X_{\cdot,N}]=
    \begin{bmatrix}
        x_{1,1}   & \cdots & x_{1,N}   \\
        \vdots    & \ddots & \vdots    \\
        x_{\omega,1} & \cdots & x_{\omega,N} \\
        \vdots    & \ddots & \vdots    \\
        x_{\Omega,1} & \cdots & x_{\Omega,N}
    \end{bmatrix}.
\]

\tcbline

\textbf{A.\ 冗余证券}\quad 若证券 $j$ 的支付列 $x_j$ 可由其余证券的\textbf{复制组合}（replicating portfolio）线性表出，即存在 $\theta_{\setminus j}$ 使得
\[
    x_j=X_{\setminus j}\theta_{\setminus j},
\]
则称 $j$ 为\textbf{冗余证券}。冗余证券不增加可达的 1 期支付集合；在无摩擦时，市场结构 $X$ 与删去第 $j$ 列后的 $X_{\setminus j}$ 在资源配置意义上\textbf{等价}，均衡与分配分析中常可忽略冗余列。\textcolor{red}{\textbf{更精确的说法：}}存在或去掉某证券\textbf{不改变均衡配置}时，该证券才是冗余的。存在\textbf{市场摩擦}时，用其他证券复制 $x_j$ 可能额外花费，表观冗余的证券在交易中仍可能不可替代。

\tcbline

\textbf{B.\ 证券市场的等价描述}\quad 无摩擦时，可用\textbf{线性无关}的 $N$ 只证券描述市场，$\mathrm{rank}(X)=N$（通常设 $N\le\Omega$）。任取 $N$ 个线性无关的组合向量 $\theta_1,\dots,\theta_N$，令 $H=[\theta_1,\cdots,\theta_N]$（$N\times N$ 可逆），则新支付矩阵 $X_\theta=XH$ 与原市场\textbf{等价}：原证券列空间不变，$X=X_\theta H^{-1}$，每只原证券都可由新组合复制。

\tcbline

\textbf{C.\ 生成（span）与完全市场}\quad 记 $\mathbf{1}_\omega$ 为第 $\omega$ 分量为 1、其余为 0 的列向量（状态 $\omega$ 阿罗证券的支付）。若 $\mathrm{rank}(X)=\Omega$ 且 $N\ge\Omega$，则对每一 $\omega$ 存在组合 $\theta_\omega$ 使 $X\theta_\omega=\mathbf{1}_\omega$。当 $X$ 为 $\Omega\times\Omega$ 可逆方阵时，解唯一：
\[
    \theta_\omega=X^{-1}\mathbf{1}_\omega.
\]
此时称不确定性由证券\textbf{生成}（span）；市场与第 3 章完全阿罗市场等价。

\begin{tcolorbox}[knowledge]
\textbf{定理 4.1}\quad 无摩擦时，证券市场\textbf{完全}当且仅当线性无关证券只数等于状态数（可取 $N=\Omega$ 且 $\mathrm{rank}(X)=\Omega$，此时 $X$ 可逆）。这与定义 3.2、\textcolor{red}{\textbf{行满秩}}判据一致。
\end{tcolorbox}
\end{tcolorbox}

\begin{tcolorbox}[knowledge,title=4.2 套利]
记交易证券的 0 期价格向量为 $S=\begin{bmatrix}S_1 \\ \vdots \\ S_N\end{bmatrix}$。组合 $\theta=\begin{bmatrix}\theta_1 \\ \vdots \\ \theta_N\end{bmatrix}$ 的 0 期成本为 $S^{\mathsf T}\theta$，1 期支付为 $X\theta$。从支付 $X$ 到价格 $S$ 的映射称为\textbf{资产定价关系}（资产定价模型）。若 $X\theta\ge 0$，称该组合在 1 期满足\textbf{有限责任}。

\begin{tcolorbox}[knowledge]
\textbf{定义 4.1}\quad 组合 $\theta$ 构成\textbf{套利机会}，若同时满足：
\begin{enumerate}
    \item $S^{\mathsf T}\theta\le 0$（0 期净支出非正）；
    \item $X\theta\ge 0$（1 期各状态支付非负）；
    \item 上述两式中\textbf{至少有一个严格不等式}成立（0 期严格收入或某状态严格正支付）。
\end{enumerate}
等价地，扩展支付向量 $\begin{bmatrix}-S^{\mathsf T}\theta \\ X\theta\end{bmatrix}$ 非负且至少一分量严格为正。
\end{tcolorbox}

常分为三类（记严格不等方向与教材一致）：
\begin{itemize}
    \item \textbf{第 1 类：}$S^{\mathsf T}\theta<0$ 且 $X\theta=0$——0 期白捡钱、未来无负债；
    \item \textbf{第 2 类：}$S^{\mathsf T}\theta=0$ 且 $X\theta>0$——\textbf{套利组合}（零成本、未来严格正支付）；
    \item \textbf{第 3 类：}$S^{\mathsf T}\theta<0$ 且 $X\theta>0$——0 期收入为正且未来仍为正。
\end{itemize}

\begin{tcolorbox}[example]
\textbf{例 4.1}\quad $S^{\mathsf T}=[1,\,1,\,2]$，
$X=\begin{bmatrix}1&0&2\\1&2&0\\1&2&0\end{bmatrix}$（三证券在各状态下的支付）。组合 $\theta_1=[2;\,-1;\,-1]$：$S^{\mathsf T}\theta_1=-1$，$X\theta_1=0$，为\textbf{第 1 类套利}。经济直觉：支付 $[2;2;2]$ 可由 2 单位证券 1 以成本 $2$ 复制，也可由 1 单位 $B$ 加 1 单位 $C$ 以成本 $3$ 复制——同支付不同价。
\end{tcolorbox}
\end{tcolorbox}

\begin{tcolorbox}[knowledge,title=4.3 无套利原理]
套利机会只依赖公开的支付 $X$ 与价格 $S$，不依赖私人信息或对状态概率的主观信念。

\begin{tcolorbox}[knowledge]
\textbf{定理 4.2}\quad 市场均衡中\textbf{不存在}套利机会。
\end{tcolorbox}

\textbf{直观：}若存在套利 $\theta$，在\textbf{不满足性}下参与者可无限放大该组合，得到严格更优的消费，与均衡最优矛盾。\textbf{前提：}至少一名参与者不满足；市场\textbf{无摩擦}（无成本、无限制、价格接受者等，见第 2.3 节）。均衡 $\Rightarrow$ 无套利是\textbf{必要}条件；无套利 $\not\Rightarrow$ 均衡（还需出清与偏好等）。

\begin{tcolorbox}[knowledge]
\textbf{定义 4.2}\quad 证券市场上不存在套利机会，称为满足\textbf{无套利原理}（no-arbitrage principle）。
\end{tcolorbox}
\end{tcolorbox}

\begin{tcolorbox}[knowledge,title=4.4 资产定价基本定理（开端）]
将证券（或组合）的 1 期支付 $x\in\mathrm{R}^{\Omega}$ 映射为 0 期价格 $S$，记\textbf{定价算子}（估值算子）$V$：
\[
    S=V(x).
    \tag{4.1}
\]
对列向量 $x$ 为某只证券支付时，$S_n=V(X_n)$。

\begin{tcolorbox}[knowledge]
\textbf{定理 4.3（一价定律）}\quad 若两只证券（或组合）未来支付相同，$x_1=x_2$，则 0 期价格必相同：$V(x_1)=V(x_2)$。
\end{tcolorbox}

无套利蕴含一价定律；反之一般不成立。进一步可证：\textbf{无套利}等价于存在\textbf{状态价格}（或正线性定价泛函）使 $S=X^{\mathsf T}\phi$（完全市场时 $\phi$ 唯一，不完全市场时 $\phi$ 可能不唯一甚至某些分量可为负，见练习 4.5）。鞅测度表述见练习 4.8。
\end{tcolorbox}


\begin{problem}
\problemno{4.1} 经济在 1 期有 4 个可能状态。在市场中有 5 只可交易证券，它们的支付矩阵 $X$ 如下：
\[
    X=\begin{bmatrix}
        1 & 1 & 0 & 0 & 0 \\
        1 & 2 & 1 & 0 & 0 \\
        1 & 3 & 2 & 1 & 0 \\
        1 & 4 & 3 & 2 & 1
    \end{bmatrix}.
\]
\begin{enumerate}[label=(\alph*)]
    \item 证明市场是完全的。
    \item 证明状态价格向量存在且唯一。并给出它的值。
    \item 证明市场上存在冗余证券。
    \item 选择一组足以保证市场完全的复合证券。把每一只 Arrow-Debreu 证券都表示成这些证券的组合。称这些组合为状态或有组合。
    \item 用状态或有组合把任一由五只交易证券组成的组合表示出来。
    \item 证明状态或有组合可生成状态空间。
\end{enumerate}
\end{problem}
\begin{solution}
记 5 只交易证券的支付向量分别为
\[
    X_{1}=\begin{bmatrix}1 \\ 1 \\ 1 \\ 1\end{bmatrix},\quad X_{2}=\begin{bmatrix}1 \\ 2 \\ 3 \\ 4\end{bmatrix},\quad X_{3}=\begin{bmatrix}0 \\ 1 \\ 2 \\ 3\end{bmatrix},\quad X_{4}=\begin{bmatrix}0 \\ 0 \\ 1 \\ 2\end{bmatrix},\quad X_{5}=\begin{bmatrix}0 \\ 0 \\ 0 \\ 1\end{bmatrix}
\]
\begin{enumerate}[label=(\alph*)]
    \item 因为 $X_{1}$、$X_{2}$、$X_{3}$、$X_{4}$ 的支付是互相独立的四只证券，而状态数等于 4，因而证券市场是完全的；
    \item 假设状态价格向量为 $\phi=\begin{bmatrix}\phi_{1} \\ \phi_{2} \\ \phi_{3} \\ \phi_{4}\end{bmatrix}$，那么有：
    \[
        \begin{aligned}
            \phi_{1}+\phi_{2}+\phi_{3}+\phi_{4}&=1\\
            \phi_{1}+2\phi_{2}+3\phi_{3}+4\phi_{4}&=2.5\\
            \phi_{2}+2\phi_{3}+3\phi_{4}&=1.5\\
            \phi_{3}+2\phi_{4}&=0.75\\
            \phi_{4}&=0.25
        \end{aligned}
    \]
    求解上述方程组可以得到唯一的解为 $\phi=\begin{bmatrix}\tfrac14 \\ \tfrac14 \\ \tfrac14 \\ \tfrac14\end{bmatrix}$；
    \item 买入 1 单位的证券 2 卖出 1 单位的证券 3，得到的支付是 $\begin{bmatrix}1 \\ 1 \\ 1 \\ 1\end{bmatrix}$，恰好就是证券 1 的支付，因而市场上存在冗余证券；
    \item 证券 2、3、4、5 可以保证市场完全，而且我们有
    \[
        \begin{aligned}
            X_{2}-2X_{3}+X_{4}&=\begin{bmatrix}1 \\ 0 \\ 0 \\ 0\end{bmatrix}\\
            X_{3}-2X_{4}+X_{5}&=\begin{bmatrix}0 \\ 1 \\ 0 \\ 0\end{bmatrix}\\
            X_{4}-2X_{5}&=\begin{bmatrix}0 \\ 0 \\ 1 \\ 0\end{bmatrix}\\
            X_{5}&=\begin{bmatrix}0 \\ 0 \\ 0 \\ 1\end{bmatrix}
        \end{aligned}
    \]
    \item 假设组合中 5 只交易证券的头寸分别为 $\theta_{1}$、$\theta_{2}$、$\theta_{3}$、$\theta_{4}$、$\theta_{5}$，那么状态或有组合的头寸分别为 $\theta_{1}+\theta_{2}$，$\theta_{1}+2\theta_{2}+\theta_{3}$，$\theta_{1}+3\theta_{2}+2\theta_{3}+\theta_{4}$，$\theta_{1}+4\theta_{2}+3\theta_{3}+2\theta_{4}+\theta_{5}$；
    \item 因为状态或有组合的支付是 $\mathrm{R}^{4}$ 空间的一组基，因而状态或有组合可生成支付空间。
\end{enumerate}
\end{solution}

\begin{problem}
\problemno{4.2} 银行 1 的存贷利率为 $r_{1}$ 而银行 2 的为 $r_{2}$，且 $r_{1}>r_{2}$。证明存在套利机会。
\end{problem}
\begin{solution}
从银行 2 借钱并全部存入银行 1，参与者就可以获得第 2 类套利。
\end{solution}

\begin{problem}
\problemno{4.3} 银行 $i$ 的贷款利率为 $r_{ia}$（即顾客借钱的利率）而存款利率为 $r_{ib}$，$i=1,2$。
\begin{enumerate}[label=(\alph*)]
    \item 证明无套利要求 $r_{ia}\ge r_{ib}$。
    \item 两个银行的存贷利率应该满足什么样的条件才能使市场上不存在套利机会。
\end{enumerate}
\end{problem}
\begin{solution}
\begin{enumerate}[label=(\alph*)]
    \item 反证法。如果 $r_{ia}<r_{ib}$，那么从银行 $i$ 借入 1 块钱并存入银行 $i$，那么在 0 期支付为 0 而 1 期支付为 $r_{ib}r_{ia}>0$，参与者可获得第 2 类套利。因此，$r_{ia}\ge r_{ib}$；
    \item $r_{1a}>r_{2b}$，$r_{1b}<r_{2a}$。
\end{enumerate}
\end{solution}

\begin{problem}
\problemno{4.4} 考虑练习 4.1 中定义的市场。假设 5 只交易证券的价格分别为 $1$，$2$，$3/2$，$1$，$1/2$。
\begin{enumerate}[label=(\alph*)]
    \item 在这样的价格和支付下存在套利机会吗？
    \item 你能构建第 1、2 以及 3 类套利机会吗？
\end{enumerate}
\end{problem}
\begin{solution}
\begin{enumerate}[label=(\alph*)]
    \item 买入 1 单位的证券 4，卖出 2 单位的证券 5，0 期支付为 0，而 1 期支付为 $\begin{bmatrix}0 \\ 0 \\ 1 \\ 0\end{bmatrix}>0$，因而存在套利机会。
    \item 上面构造了第 2 类套利机会。卖出 1 单位的证券 1 和 1 单位的证券 3，同时买入 1 单位的证券 2，那么 0 期支付为 $\tfrac12$，而 1 期支付为 0，这是第 1 类套利机会。同时进行上述两种交易就可以得到第 3 类套利机会。
\end{enumerate}
\end{solution}

\begin{problem}
\problemno{4.5} 在 1 期有三个可能状态：$\{\omega_{1},\omega_{2},\omega_{3}\}$。在经济中有三只交易证券，$a$，$b$，$c$，它们的支付向量如下：
\[
    X_{a}=\begin{bmatrix}1\\1\\0\end{bmatrix},\quad
    X_{b}=\begin{bmatrix}0\\1\\1\end{bmatrix},\quad
    X_{c}=\begin{bmatrix}1/2\\1\\1/2\end{bmatrix}.
\]
这些证券的 0 期均衡价格为：$S_{c}=S_{b}=S_{c}=1$。令 $\phi=\begin{bmatrix}\phi_{1} \\ \phi_{2} \\ \phi_{3}\end{bmatrix}$ 为状态价格向量。
\begin{enumerate}[label=(\alph*)]
    \item 证明 $\begin{bmatrix}2 \\ -1 \\ 2\end{bmatrix}$ 是一个可能的状态价格向量，因为它给出了三支现有证券的价格。
    \item “资产定价基本定理”指出状态价格必须是非负的。在一个两期经济里，对于证券市场这意味着什么？并给出其解释。
    \item 上面的例子是这个原理的反例吗？为什么？
    \item 证明 $\begin{bmatrix}1/2 \\ 1/2 \\ 1/2\end{bmatrix}$ 是另外一个可能的状态价格向量。在这种情况下，所有状态价格都是正的。
    \item 解释为什么存在多个可能的状态价格向量，它们对已有证券都给出相同的价格。
    \item 考虑一个标的资产为证券 $a$、执行价格为 $1/2$ 的欧式买权。这份期权的支付是多少？你能为它定价吗？
\end{enumerate}
\end{problem}
\begin{solution}
\begin{enumerate}[label=(\alph*)]
    \item 因为
    \[
        2\times 1+(-1)\times 1=1,\quad (-1)\times 1+2\times 1=1,\quad 2\times\tfrac12+(-1)\times 1+2\times\tfrac12=1
    \]
    因而 $\begin{bmatrix}2 \\ -\tfrac12 \\ 2\end{bmatrix}$ 是一个可能的状态价格向量；
    \item 状态价格向量的非负性意味着不存在套利机会；
    \item 不是。因而它只是可能的状态价格之一，并不一定是真实的状态价格；
    \item 因为
    \[
        \tfrac12\times 1+\tfrac12\times 1=1,\quad \tfrac12\times 1+\tfrac12\times 1=1,\quad \tfrac12\times\tfrac12+\tfrac12\times 1+\tfrac12\times\tfrac12=1
    \]
    因而 $\begin{bmatrix}\tfrac12 \\ \tfrac12 \\ \tfrac12\end{bmatrix}$ 是一个可能的状态价格向量；
    \item 因为市场是不完全的，因而可能存在多个状态价格，它们都给出已有证券的价格。
\end{enumerate}
\end{solution}

\begin{problem}
\problemno{4.6} 远期合约（Forward）是在未来（即 1 期）以当期（即 0 期）确定的价格交易某一证券或商品的合约，而在当期没有现金流产生。在当期约定的未来交易价格称作远期价格。比方说，考虑一份在 1 期以 $F_{0}$ 的价格购买 1 盎司黄金的远期合约。合约的买方同意在 1 期（合约的到期日）从卖方买入 1 盎司黄金。在当期没有现金交换。令 $S_{t}$ 为 $t=0,1$ 期黄金的价格，而 $r_{F}$ 为利率。考虑在 1 期得到 1 盎司黄金的两种方法：
\begin{itemize}
    \item 现在购买一份远期合约。在 1 期，支付远期价格 $F_{0}$ 并获得黄金。
    \item 现在支付 $S_{0}$ 买入 1 盎司黄金，并且持有它直到 1 期。假设从现在起持有黄金直到 1 期没有另外的成本或收益。
\end{itemize}
在这两种情况下，我们在 1 期同样得到 1 盎司黄金，而那时它的价值为 $S_{1}$。
\begin{enumerate}[label=(\alph*)]
    \item 由于这两种策略给予我们相同的支付 $S_{1}$，无套利原理意味着它们的成本必须一样。用当期的价格 $S_{1}$ 和 $r_{F}$ 给出远期价格。
    \item 假设持有 1 盎司黄金不仅得到它的 1 期出售收益，在 1 期还得到额外收益 $y$。$y$ 将如何影响远期价格？
\end{enumerate}
\end{problem}
\begin{solution}
\begin{enumerate}[label=(\alph*)]
    \item 第一种方法的成本是 $\dfrac{F_{0}}{1+r_{F}}$，而第二种方法的成本是 $S_{0}$，根据无套利原理，$\dfrac{F_{0}}{1+r_{F}}=S_{0}$，因此，$F_{0}=(1+r_{F})S_{0}$；
    \item 由于第二种方法还得到额外收益，由无套利原理原理有 $F_{0}<(1+r_{F})S_{0}$。（如果收益是确定性的，那么应该有 $F_{0}+y=(1+r_{F})S_{0}$）
\end{enumerate}
\end{solution}

\begin{problem}
\problemno{4.7} 英镑现时对美元的价格（即英镑/美元汇率）为 $S_{0}$（即一英镑的美元价格为本 $S_{0}$），美元利率为 $r_{A}$ 而英镑利率为 $r_{B}$。考虑一份在 1 期购买 1 英镑的远期合约。远期价格（即远期汇率）应该是多少？提示：考虑在 1 期得到 1 英镑的两种策略。
\begin{itemize}
    \item 购买一份远期合约。在 1 期，支付远期价格 $F_{0}$ 购入 1 英镑。
    \item 现在以当前汇率 $S_{0}$ 买入 $x$ 英镑并以 $r_{B}$ 的利率进行投资，在 1 期得到 1 英镑。
\end{itemize}
得到的远期价格、现时汇率和利率之间的关系就叫作抛补利率平价理论（Covered Interest Rate Parity）。
\end{problem}
\begin{solution}
第一种策略的 0 期支付是 $\dfrac{F_{0}}{1+r_{A}}$ 美元，而第二种策略的 0 期支付是 $\dfrac{S_{0}}{1+r_{B}}$ 美元，根据无套利原理，$\dfrac{F_{0}}{1+r_{A}}=\dfrac{S_{0}}{1+r_{B}}$，即 $F_{0}=\dfrac{1+r_{A}}{1+r_{B}}S_{0}$。
\end{solution}

\begin{problem}
\problemno{4.8} 考虑练习 4.1 定义的市场。假设证券 1，2，5 的价格是 1，2，$1/2$。
\begin{enumerate}[label=(\alph*)]
    \item 计算状态价格和利率。
    \item 构建等价鞅测度。
    \item 证明 5 只证券的价格都具有鞅性质。
\end{enumerate}
\end{problem}
\begin{solution}
\begin{enumerate}[label=(\alph*)]
    \item 假设状态价格向量为 $\phi=\begin{bmatrix}\phi_{1} \\ \phi_{2} \\ \phi_{3} \\ \phi_{4}\end{bmatrix}$，那么有：
    \[
        \begin{aligned}
            \phi_{1}+\phi_{2}+\phi_{3}+\phi_{4}&=0.8\\
            \phi_{1}+2\phi_{2}+3\phi_{3}+4\phi_{4}&=2\\
            \phi_{2}+2\phi_{3}+3\phi_{4}&=1.2\\
            \phi_{4}&=0.2
        \end{aligned}
    \]
    求解上述方程组可以得到唯一的解为 $\phi=\begin{bmatrix}0.2 \\ 0.2 \\ 0.2 \\ 0.2\end{bmatrix}$。因此，证券 4 的价格为 $0.6$，而利率为 $r_{F}=\dfrac{1}{\phi_{1}+\phi_{2}+\phi_{3}+\phi_{4}}-1=0.25$；
    \item 等价鞅测度为 $Q=\begin{bmatrix}0.25 \\ 0.25 \\ 0.25 \\ 0.25\end{bmatrix}$；
    \item 对于证券 1，$\dfrac{\mathrm{E}^{Q}[X_{1}]}{1+r_{F}}=\dfrac{1}{1+0.25}=0.8$，正好就是它的价格；类似可以证明其它证券的价格也有鞅性质。
\end{enumerate}

Notes：证券价格不完全，现在改为证券 1，2，3，5 的价格分别为 0.8，2，1.2，0.2。
\end{solution}

% !TeX root = ../../金融经济学习题讲义.tex
% 源 PDF：references/金融经济学/史上最全最完整的金融经济学答案(王江版).pdf，页 20–25（印刷页 19–24）

\zhchapter{期权：一个套利定价的例子}

\ifdefined\WJHandoutKnowledge
\fi

% !TeX root = ../../金融经济学习题讲义.tex
% 05_期权 课后习题（由章节文件拆分）

\begin{problem}
\problemno{5.1} 1 期有两个可能状态 $a$ 和 $b$，发生概率分别为 $\pi$ 和 $1-\pi$。存在一只收益率为 $r_{F}$ 的无风险证券。还存在一只风险证券，它的现时价格为 $S$，在 1 期支付红利 $\delta S$。它在 1 期的除权价（即红利后的价格）在状态 $a$ 下是 $(1+r_{F}+u)S$、在状态 $b$ 下是 $(1+r_{F}-d)S$。
\begin{enumerate}[label=(\alph*)]
    \item 为使市场中不存在套利机会，$\pi$，$r_{F}$，$u$，$d$ 和 $\delta$ 应该满足什么条件？
    \item 1 期到期、执行价为 $K$ 的欧式买权的 0 期价格是多少？
    \item 美式买权的价格是多少？比较它与相应的欧式买权价格。什么时候提前执行是最优的？
    \item 将买权换成卖权，回答 (b) 和 (c)。
\end{enumerate}
\end{problem}
\begin{solution}
\begin{enumerate}[label=(\alph*)]
    \item $u+\delta>1+r_{F}>d+\delta$；
    \item 令状态价格向量为 $\phi=\begin{bmatrix}\phi_{a} \\ \phi_{b}\end{bmatrix}$，那么有
    \[
        \begin{aligned}
            \phi_{a}+\phi_{b}&=\frac{1}{1+r_{F}}\\
            (u+\delta)S\phi_{a}+(d+\delta)S\phi_{b}&=S
        \end{aligned}
    \]
    求解方程组可得状态价格为
    \[
        \begin{aligned}
            \phi_{a}&=\frac{1}{1+r_{F}}\frac{(1+r_{F})-(d+\delta)}{u-d}\\
            \phi_{b}&=\frac{1}{1+r_{F}}\frac{(u+\delta)-(1+r_{F})}{u-d}
        \end{aligned}
    \]
    因而期权的价格为
    \[
        c=\phi_{a}[uS-K]_{+}+\phi_{b}[dS-K]_{+}=
        \begin{cases}
            0 & uS<K\\
            \phi_{a}(uS-K) & dS\le K\le uS\\
            \phi_{a}(uS-K)+\phi_{b}(dS-K) & K<dS
        \end{cases}
    \]
    \item 对于美式期权来说，如果在 0 期执行，那么获得的收益是 $[S-K]_{+}$，如果在 1 期执行，那么收益和欧式期权一样，因而 $C=\max\{[S-K]_{+},c\}$。当且仅当 $[S-K]_{+}>c$ 时提前执行是最优的。当 $\delta$ 很大时，股价 $S$ 中含有股利的成份越多，因而越有可能提前执行。
    \item 对于欧式看跌期权，价格为
    \[
        p=\phi_{a}[K-uS]_{+}+\phi_{b}[K-dS]_{+}=
        \begin{cases}
            0 & K<dS\\
            \phi_{b}(K-dS) & dS\le K\le uS\\
            \phi_{a}(K-uS)+\phi_{b}(K-dS) & K>uS
        \end{cases}
    \]
    而对于美式看跌期权，$P=\max\{[K-S]_{+},p\}$。当 $\delta$ 越大时，股价 $S$ 中含有股利的成份越多，因而越有可能不提前执行。
\end{enumerate}
\end{solution}

\begin{problem}
\problemno{5.2} 在 1 期有三个可能状态：$\{\omega_{1},\omega_{2},\omega_{3}\}$，在经济中有三只交易证券，它们的支付向量如下：
\[
    X_{a}=\begin{bmatrix}1\\1\\0\end{bmatrix},\quad
    X_{b}=\begin{bmatrix}0\\1\\1\end{bmatrix},\quad
    X_{c}=\begin{bmatrix}1/2\\1\\1/2\end{bmatrix}.
\]
这些证券的 0 期均衡价格是 $S_{a}=S_{b}=S_{c}=1$。令 $\phi=\begin{bmatrix}\phi_{1} \\ \phi_{2} \\ \phi_{3}\end{bmatrix}$ 为状态价格向量。[这与前一章练习 4.5 定义的经济一样。]考虑一份标的资产证券为 $a$、执行价格为 $1/2$ 的欧式买权。它的支付是什么？你能确定它的价格吗？
\end{problem}
\begin{solution}
它的支付是 $\begin{bmatrix}1/2\\1/2\\0\end{bmatrix}$，因为它的支付恰为证券 $a$ 的一半，因而它的价格也是一半即 $\tfrac12$。
\end{solution}

\begin{problem}
\problemno{5.3} 蝴蝶头寸是看涨期权的组合：（1）买入 1 份执行价格为 $K-\delta$ 的看涨期权，（2）卖出 2 份执行价格为 $K$ 的看涨期权，（3）买入 1 份执行价格为 $K+\delta$ 的看涨期权。因此，蝴蝶头寸可由两个参数 $K$ 和 $\delta$ 刻划。比较两个蝴蝶头寸，$K$ 相同但 $\delta$ 不同。证明 $\delta$ 较高的蝴蝶头寸的价值较高。
\end{problem}
\begin{solution}
假设 $\delta_{1}>\delta_{0}$，比较两只期权的支付。当 1 期价格 $S$ 大于 $K+\delta$ 或者小于 $K-\delta$ 时，两者的支付都为 0；当 $S\in[K-\delta_{1},K-\delta_{0}]$ 时 $X_{1}=S-(K-\delta_{1})>X_{0}=0$；当 $S\in[K-\delta_{0},K]$ 时 $X_{1}=S-(K-\delta_{1})>S-(K-\delta_{0})=X_{0}$；当 $S\in[K,K+\delta_{0}]$ 时 $X_{1}=-S+(K+\delta_{1})>-S+(K+\delta_{0})=X_{0}$；当 $S\in[K+\delta_{0},K+\delta_{1}]$ 时 $X_{1}=-S+(K+\delta_{1})>0=X_{0}$；根据无套利原理，$\delta$ 较高的期权价格较高。
\end{solution}

\begin{problem}
\problemno{5.4} 对于参数为 $K$ 和 $\delta$ 的蝴蝶头寸，给出它的价格的一个上界和下界。
\end{problem}
\begin{solution}
一个上界是 $[S-(K-\delta)]_{+}-[S-(K+\delta)]_{+}$，因为期权的支付是非负的，因而 0 是它的一个下界。
\end{solution}

\begin{problem}
\problemno{5.5} 股票价格具有如下形式的二叉树结构：
\begin{center}
\begin{tikzpicture}[x=1.35em,y=1.15em,baseline=(current bounding box.center)]
    \node[anchor=east] at (-0.15,0) {$10$};
    \draw[thick] (0,0) -- (1.1,0.7);
    \draw[thick] (0,0) -- (1.1,-0.7);
    \draw[thick] (1.1,0.35) -- (1.1,1.05);
    \draw[thick] (1.1,0.7) -- (1.85,0.7);
    \draw[thick] (1.1,-0.7) -- (1.85,-0.7);
    \node[anchor=west] at (2.0,0.7) {$12$};
    \node[anchor=west] at (2.0,-0.7) {$8$};
\end{tikzpicture}
\end{center}
两个可能结果的发生概率相等，无风险利率是 $10\%$。
\begin{enumerate}[label=(\alph*)]
    \item 为一个以此股票为标的资产、执行价格为 9 的看跌期权定价。
    \item 假设你得到一个新的利好信息：股票价格上升到 12 的概率大于 $1/2$，而股票价格仍然是 10。你愿意支付的看跌期权的价格是上涨还是下跌？解释原因。
\end{enumerate}
\end{problem}
\begin{solution}
\begin{enumerate}[label=(\alph*)]
    \item 由第 1 题中公式可以得到状态价格为 $\phi=[\tfrac{15}{22},\tfrac{5}{22}]$，因而看跌期权的价格为 $\dfrac{5}{22}\times[9-8]=\dfrac{5}{22}$；
    \item 由于股价不变，状态价格也不会变化，因而看跌期权的价格是不变的。
\end{enumerate}
\end{solution}

\begin{problem}
\problemno{5.6} 二叉树模型可以拓展到多于两个时期的情形。考虑拓展到三个时期 0，1，2 的情况。股票价格服从如下的二叉树过程：
\begin{center}
\begin{tikzpicture}[x=1.6em,y=1.15em,baseline=(current bounding box.center)]
    \node[anchor=east] at (0,0) {$10$};
    \draw[thick] (0.15,0) -- (1.5,1.6);
    \draw[thick] (0.15,0) -- (1.5,-1.6);
    \node[anchor=east] at (1.5,1.6) {$12$};
    \node[anchor=east] at (1.5,-1.6) {$8$};
    \draw[thick] (1.65,1.6) -- (3.2,2.7);
    \draw[thick] (1.65,1.6) -- (3.2,0.7);
    \draw[thick] (1.65,-1.6) -- (3.2,-0.7);
    \draw[thick] (1.65,-1.6) -- (3.2,-2.7);
    \draw[thick] (3.2,2.2) -- (3.2,3.2);
    \draw[thick] (3.2,2.7) -- (4.0,2.7);
    \draw[thick] (3.2,0.7) -- (4.0,0.7);
    \draw[thick] (3.2,-0.7) -- (4.0,-0.7);
    \draw[thick] (3.2,-2.7) -- (4.0,-2.7);
    \node[anchor=west] at (4.15,2.7) {$14.4$};
    \node[anchor=west] at (4.15,0.7) {$9.6$};
    \node[anchor=west] at (4.15,-0.7) {$9.6$};
    \node[anchor=west] at (4.15,-2.7) {$6.4$};
\end{tikzpicture}
\end{center}
每一期间（即，从 0 期到 1 期或从 1 期到 2 期）的利率都是 $10\%$。考虑一份在 2 期到期、执行价格为 9 的看涨期权。
\begin{enumerate}[label=(\alph*)]
    \item 描述这份看涨期权在 2 期的支付。
    \item 求解看涨期权在 1 期、股票价格分别为 12 或 8 时的价格。
    \item 求解看涨期权在 0 期的价格。
    \item 如果是美式看涨期权，它在 0 期的价格是多少？
\end{enumerate}
\end{problem}
\begin{solution}
\begin{enumerate}[label=(\alph*)]
    \item 支付是
    \[
        \begin{bmatrix}5.4\\0.6\\0.6\\0\end{bmatrix}
    \]
    \item 可以求得每个节点的状态价格都和第 5 题中的一样。当股价为 12 时，期权价格为 $\dfrac{15}{22}\times 5.4+\dfrac{5}{22}\times 0.6=3.82$；当股价为 8 时，期权价格为 $\dfrac{15}{22}\times 0.6=0.41$；
    \item 看涨期权的 0 期价格为 $\dfrac{15}{22}\times 3.82+\dfrac{5}{22}\times 0.41=2.70$；
    \item 因为没有股利，美式看涨期权不会提前执行，价格也是 2.70。
\end{enumerate}
\end{solution}

\begin{problem}
\problemno{5.7*} 在练习 5.6 描述的经济中，考虑一份以股票为标的资产、在 2 期到期、执行价格为 14 的看跌期权。
\begin{enumerate}[label=(\alph*)]
    \item 如果它是欧式期权，求解它的 0 期价格。
    \item 如果它是美式期权，求解它的 0 期价格。
\end{enumerate}
\end{problem}
\begin{solution}
\begin{enumerate}[label=(\alph*)]
    \item 期权的 2 期支付是 $\begin{bmatrix}0\\4.4\\4.4\\7.6\end{bmatrix}$；在 1 期，当股价为 12 时，期权价格为 $\dfrac{5}{22}\times 4.4=1$；当股价为 8 时，期权价格为 $\dfrac{15}{22}\times 4.4+\dfrac{5}{22}\times 7.6=4.73$；因而期权的 0 期价格为 $\dfrac{15}{22}\times 1+\dfrac{5}{22}\times 4.73=1.76$。
    \item 期权的 2 期支付是 $\begin{bmatrix}0\\4.4\\4.4\\7.6\end{bmatrix}$；在 1 期如果不提前执行，当股价为 12 时，期权价格为 $\dfrac{5}{22}\times 4.4=1$；当股价为 8 时，期权价格为 $\dfrac{15}{22}\times 4.4+\dfrac{5}{22}\times 7.6=4.73$；而如果在 1 期提前执行，那么期权的支付分别为 2 和 6，因而参与者会提前执行，期权在 1 期的支付为 $\begin{bmatrix}2 \\ 6\end{bmatrix}$；如果在 0 期不提前执行，期权的价格为 $\dfrac{15}{22}\times 2+\dfrac{5}{22}\times 6=2.73$，而如果在 0 期提前执行，参与者的收益是 4，因此参与者在 0 期就执行期权，期权在 0 期的价格是 4。
\end{enumerate}
\end{solution}

\begin{problem}
\problemno{5.8*} 这里的经济与练习 5.6 描述的经济一样，只是这里的参数更一般性：
\begin{center}
\begin{tikzpicture}[x=1.6em,y=1.15em,baseline=(current bounding box.center)]
    \node[anchor=east] at (0,0) {$S$};
    \draw[thick] (0.15,0) -- (1.5,1.6);
    \draw[thick] (0.15,0) -- (1.5,-1.6);
    \node[anchor=east] at (1.5,1.6) {$uS$};
    \node[anchor=east] at (1.5,-1.6) {$dS$};
    \draw[thick] (1.65,1.6) -- (3.2,2.7);
    \draw[thick] (1.65,1.6) -- (3.2,0.7);
    \draw[thick] (1.65,-1.6) -- (3.2,-0.7);
    \draw[thick] (1.65,-1.6) -- (3.2,-2.7);
    \draw[thick] (3.2,2.2) -- (3.2,3.2);
    \draw[thick] (3.2,2.7) -- (4.0,2.7);
    \draw[thick] (3.2,0.7) -- (4.0,0.7);
    \draw[thick] (3.2,-0.7) -- (4.0,-0.7);
    \draw[thick] (3.2,-2.7) -- (4.0,-2.7);
    \node[anchor=west] at (4.15,2.7) {$u^{2}S$};
    \node[anchor=west] at (4.15,0.7) {$udS$};
    \node[anchor=west] at (4.15,-0.7) {$duS$};
    \node[anchor=west] at (4.15,-2.7) {$d^{2}S$};
\end{tikzpicture}
\end{center}
每一期间的利率都是 $r_{F}$ 且 $0<d<1+r_{F}<u$。
\begin{enumerate}[label=(\alph*)]
    \item 求解在每一时期考虑下一期的两种可能情形（即股价上升或下降）的风险中性概率。
    \item 用 $S$，$u$，$d$ 和 $r_{F}$ 表示欧式看跌期权在 0 期的价格。
    \item 证明在 2 期到期的美式看涨期权的价格比具有相同执行价格、但在 1 期到期的美式看涨期权的价格要高。
\end{enumerate}
\end{problem}
\begin{solution}
\begin{enumerate}[label=(\alph*)]
    \item 状态价格在每期都是一样的，
    \[
        \begin{aligned}
            \phi_{u}&=\frac{1}{1+r_{F}}\frac{(1+r_{F})-d}{u-d}\\
            \phi_{d}&=\frac{1}{1+r_{F}}\frac{u-(1+r_{F})}{u-d}
        \end{aligned}
    \]
    因此，每一个节点处的风险中性概率也是一样的，
    \[
        q_{u}=\frac{(1+r_{F})-d}{u-d},\quad q_{d}=\frac{u-(1+r_{F})}{u-d}
    \]
    \item 欧式看跌期权在 0 期的价格为
    \[
        \begin{aligned}
            &\phi_{u}\bigl(\phi_{u}[K-u^{2}S]_{+}+\phi_{d}[K-udS]_{+}\bigr)+\phi_{d}\bigl(\phi_{u}[K-udS]_{+}+\phi_{d}[K-d^{2}S]_{+}\bigr)\\
            &=\phi_{u}^{2}[K-u^{2}S]_{+}+2\phi_{u}\phi_{d}[K-udS]_{+}+\phi_{d}^{2}[K-d^{2}S]_{+}
        \end{aligned}
    \]
    \item 因为美式期权可以提前执行，在 2 期到期的期权既可以选在 1 期执行也可以选择在 2 期执行，而在 1 期到期的美式看涨期权只会选择在 1 期执行（虽然都可以在 0 期执行，但是在没有股利的情况下，美式看涨期权是不会提前执行的），因而在 2 期到期的期权价格更高。
\end{enumerate}
\end{solution}

\begin{problem}
\problemno{5.9*} 考虑练习 5.8 描述的经济。假设一只证券在 2 期的支付是股票价格从 0 期到 2 期所达到的最大值。（这是一个关于回望期权的例子，它的支付依赖于标的证券的价格路径）。
\begin{enumerate}[label=(\alph*)]
    \item 写出它在 2 期的支付。
    \item 导出它的 0 期价格。
    \item 证明用在 2 期支付一美元的无风险债券的价格正规化以后，这只证券的价格过程在等价鞅测度下是鞅。
\end{enumerate}
\end{problem}
\begin{solution}
这里只考虑 $ud<1$ 的情形，其它情形可以类似考虑。
\begin{enumerate}[label=(\alph*)]
    \item 期权在 2 期的支付是
    \[
        \begin{bmatrix}u^{2}S\\uS\\S\end{bmatrix}
    \]
    \item 0 期价格为 $\phi_{u}^{2}u^{2}S+\phi_{u}\phi_{d}(u+1)S+\phi_{d}^{2}S$；
    \item 在 1 期，当股价为 $uS$ 时，
    \[
        \mathrm{E}^{Q}[\tilde{X}_{2}]=\frac{1}{\phi_{u}+\phi_{d}}(\phi_{u}u^{2}S+\phi_{d}uS)
    \]
    而正规化后的价格为
    \[
        (\phi_{u}u^{2}S+\phi_{d}uS)\Big/\frac{1}{1+r_{F}}=(1+r_{F})(\phi_{u}u^{2}S+\phi_{d}uS)=\frac{1}{\phi_{u}+\phi_{d}}(\phi_{u}u^{2}S+\phi_{d}uS)
    \]
    可以验证其它各节点均如此，因而期权的价格过程为鞅。
\end{enumerate}
\end{solution}

\begin{problem}
\problemno{5.10} 经济在 1 期有 3 个状态。状态一指数证券在三个状态下的支付分别为 1，2，3：
\begin{center}
\WJEventTree[stem]{1,2,3}
\end{center}
并有两只以指数证券为标的资产、执行价格分别为 1 和 2 的看涨期权。用指数证券和上述期权复制如下支付：
\begin{center}
\WJEventTree[stem]{1,1,1}
\qquad
\WJEventTree[stem]{3,2,1}
\qquad
\WJEventTree[stem]{1,0,0}
\end{center}
\end{problem}
\begin{solution}
记状态指数证券的支付向量为 $X_{0}=\begin{bmatrix}1\\2\\3\end{bmatrix}$，执行价格为 1 的指数期权的支付向量为 $X_{1}=\begin{bmatrix}0\\1\\2\end{bmatrix}$，执行价格为 2 的指数期权的支付向量为 $X_{2}=\begin{bmatrix}0\\0\\1\end{bmatrix}$ 那么我们可得：
\[
    X_{0}-X_{1}=\begin{bmatrix}1\\1\\1\end{bmatrix},\quad
    3X_{0}-4X_{1}=\begin{bmatrix}3\\2\\1\end{bmatrix},\quad
    X_{0}+X_{2}-2X_{1}=\begin{bmatrix}1\\0\\0\end{bmatrix}.
\]
\end{solution}



% !TeX root = ../../金融经济学习题讲义.tex
% 源 PDF：references/金融经济学/史上最全最完整的金融经济学答案(王江版).pdf，页 26–27（印刷页 25–26）

\zhchapter{期望效用函数}

\ifdefined\WJHandoutKnowledge
\fi

% !TeX root = ../../金融经济学习题讲义.tex
% 06_期望效用函数 课后习题（由章节文件拆分）

\begin{problem}
\problemno{6.1} 证明如下形式的期望效用函数：
\[
    U(c)=u_{0}(c_{0})+\sum_{\omega\in\Omega}\pi_{\omega}u_{1}(c_{1\omega})
\]
满足独立性公理。
\end{problem}
\begin{solution}
假设 $c$ 和 $c'$ 在状态 $\omega'$ 有相同的消费路径 $x=[c_{0},c_{x}]$，且 $c\succeq c'$，那么有
\[
    \begin{aligned}
        U(c)
        &=u_{0}(c_{0})+\sum_{\omega\in\Omega-\{\omega'\}}\pi_{\omega}u_{1}(c_{1\omega})+\pi_{\omega'}u_{1}(c_{x})\\
        &\ge u_{0}(c_{0})+\sum_{\omega\in\Omega-\{\omega'\}}\pi_{\omega}u_{1}(c_{1\omega}')+\pi_{\omega'}u_{1}(c_{x})=U(c')
    \end{aligned}
\]
\[
    \Leftrightarrow\sum_{\omega\in\Omega-\{\omega'\}}\pi_{\omega}u_{1}(c_{1\omega})\ge\sum_{\omega\in\Omega-\{\omega'\}}\pi_{\omega}u_{1}(c_{1\omega}')
\]
现在把消费路径换成 $y=[c_{0},c_{y}]$，那么我们有
\[
    \begin{aligned}
        U(c)
        &=u_{0}(c_{0})+\sum_{\omega\in\Omega-\{\omega'\}}\pi_{\omega}u_{1}(c_{1\omega})+\pi_{\omega'}u_{1}(c_{y})\\
        &\ge u_{0}(c_{0})+\sum_{\omega\in\Omega-\{\omega'\}}\pi_{\omega}u_{1}(c_{1\omega}')+\pi_{\omega'}u_{1}(c_{y})=U(c')
    \end{aligned}
\]
即 $c\succeq c'$ 仍然成立，故 $U(c)$ 满足独立性公理。
\end{solution}

\begin{problem}
\problemno{6.2*} 对于由练习 6.1 给出的效用函数表示的偏好，证明连续性意味着 $u_{0}(\cdot)$ 和 $u_{1}(\cdot)$ 都是连续的。
\end{problem}
\begin{solution}
我们只证明 $u_{0}(\cdot)$ 是连续的，$u_{1}(\cdot)$ 的连续性可以类似证明。

假设序列 $\{x_{n}\}_{n=1}^{\infty}$ 的极限是 $x$，且 $u_{0}(x_{n})\ge(\le)\,u_{0}(c_{0})$，那么我们只须证明 $u_{0}(x)\ge(\le)\,u_{0}(c_{0})$。令 $c$ 为 0 期消费为 $c_{0}$、1 期各状态消费均为 $y>0$ 的消费计划，而 $c_{n}'$ 为 0 期消费为 $x_{n}$、1 期各状态消费均为 $y$ 的消费计划，$c'$ 为 0 期消费为 $x$、1 期各状态消费均为 $y$ 的消费计划。由假设我们有 $c_{n}'\in\{a\in\mathrm{C}:a\succeq(\preceq)c\}$，由效用函数的连续性我们有 $c'$ 也是 $c_{n}'$ 的极限，因而 $c'\in\{a\in\mathrm{C}:a\succeq(\preceq)c\}$，因而 $u_{0}(x)\ge(\le)\,u_{0}(c_{0})$，$u_{0}(\cdot)$ 是连续的。
\end{solution}

\begin{problem}
\problemno{6.3**} 如果 $a\succeq b$ 意味着对于 $\forall\,\alpha\in[0,1]$ 都有 $\alpha a+(1-\alpha)b\succeq b$，那么称偏好关系 $\succ$ 是弱凸的（weakly convex）。这与第 2 章中的强凸性有所不同。假设偏好关系是连续的并且有练习 6.1 给出的期望效用函数表示（因而 $u_{0}(\cdot)$ 和 $u_{1}(\cdot)$ 都是连续函数）。证明此时强凸性包含弱凸性。
\end{problem}
\begin{solution}
当 $\alpha=0$ 或 1 时，$\alpha a+(1-\alpha)b\succeq b$ 显然成立；当 $a\succ b$，$\alpha\in(0,1)$ 时，由强凸性有 $\alpha a+(1-\alpha)b\succ b$；现在证明如果 $\alpha\in(0,1)$，$a\sim b$，那么 $\alpha a+(1-\alpha)b\succeq b$。不妨假设效用函数是递增的，另外假设 $\epsilon>0$ 为任意正实数，令 $a'=a+\begin{bmatrix}\epsilon \\ 0 \\ 0 \\ \ldots \\ 0\end{bmatrix}\succ b$，由强凸性有 $\alpha a'+(1-\alpha)b\succ b$，即
\[
    u_{0}(\alpha c_{0}^{a}+(1-\alpha)c_{0}^{b}+\alpha\epsilon)+\sum\pi_{\omega}u_{1}(\alpha c_{1\omega}^{a}+(1-\alpha)c_{1\omega}^{b})>u_{0}(c_{0}^{b})+\sum\pi_{\omega}u_{1}(c_{1\omega}^{b})
\]
由于 $\epsilon$ 是任意实数，并且 $u_{0}$ 是连续函数，因此
\[
    u_{0}(\alpha c_{0}^{a}+(1-\alpha)c_{0}^{b})+\sum\pi_{\omega}u_{1}(\alpha c_{1\omega}^{a}+(1-\alpha)c_{1\omega}^{b})\ge u_{0}(c_{0}^{b})+\sum\pi_{\omega}u_{1}(c_{1\omega}^{b})
\]
即 $\alpha a+(1-\alpha)b\succeq b$ 成立，命题得证。
\end{solution}

\begin{problem}
\problemno{6.4**} 假设定义在随机支付上的效用函数具有练习 6.1 中的形式，其中 $\Omega$ 是任意的状态空间而 $P$ 是相应的概率测度。如果它还满足弱凸性，那么 $u_{0}(\cdot)$ 和 $u_{1}(\cdot)$ 都是凹的。
\end{problem}
\begin{solution}
只能证明 $\{c\in\mathrm{C}:c\succeq b\}$ 是凸的，即 $u$ 是拟凹的，不能证明是凹的。
\end{solution}



% !TeX root = ../../金融经济学习题讲义.tex
% 源 PDF：references/金融经济学/史上最全最完整的金融经济学答案(王江版).pdf，页 28–31（印刷页 27–30）

\zhchapter{风险厌恶}

\ifdefined\WJHandoutKnowledge
\fi

% !TeX root = ../../金融经济学习题讲义.tex
% 07_风险厌恶 课后习题（由章节文件拆分）

\begin{problem}
\problemno{7.1**} 令 $R(w)\equiv -wu''(w)/u'(w)$ 而
\[
    a(w)=\operatorname{ArgMax}E\bigl[u(w+a\tilde{r})\bigr]
\]
为相对于效用函数 $u(\cdot)$ 的最优投资选择。假设 $w>0$、$u'(\cdot)>0$、$u''(\cdot)\le 0$ 且 $E[\tilde{r}]>0$。证明以下命题：
\begin{enumerate}[label=(\alph*)]
    \item 如果 $\lambda>1$，给定效用函数 $u(\cdot)$ 且 $R'(w)\ge 0$。定义
    \[
        u_{1}(w)\equiv u(\lambda w),\qquad u_{2}\equiv u(w).
    \]
    那么，$R_{1}(w)\ge R_{2}(w)$。
    \item 给定以上结论，我们有 $a_{1}(w)\le a_{2}(w)$，其中 $a_{1}(w)$ 和 $a_{2}(w)$ 分别是相对于 $u_{1}$ 和 $u_{2}$ 的最优投资选择。
    \item 定义 $\widehat{a}(w)\equiv a(w)/w$。那么，对于 $w_{1}\le w_{2}$ 有 $\widehat{a}(w_{1})\ge\widehat{a}(w_{2})$。（提示：令 $\lambda=w_{2}/w_{1}$。）
\end{enumerate}
最后一个结论说的是，对于具有单调递增的相对风险厌恶的投资者，他在风险资产上的投资额占他财富的比例随着财富的增加而减少。
\end{problem}
\begin{solution}
\begin{enumerate}[label=(\alph*)]
    \item
    \[
        u_{1}'(w)=\lambda u'(\lambda w),\qquad u_{1}''(w)=\lambda^{2}u''(\lambda w)
    \]
    由此以及 $R'(w)\ge 0$ 可得
    \[
        R_{1}(w)=\frac{wu_{1}''(w)}{u_{1}'(w)}=\frac{\lambda^{2}wu''(\lambda w)}{\lambda u'(\lambda w)}=\frac{\lambda wu''(\lambda w)}{u'(\lambda w)}=R(\lambda w)\ge R(w)=R_{2}(w);
    \]

    \item 参与者优化的一阶条件是 $E\bigl[U'\bigl(\lambda(W+a\tilde{r})\bigr)\lambda\tilde{r}\bigr]=0$，因为 $E[\tilde{r}]>0$，从而有 $a>0$。一阶条件的两边对 $\lambda$ 求导得：
    \[
        E\bigl[u''(w+\lambda a\tilde{r})(w+a\tilde{r}+\lambda a'\tilde{r})\tilde{r}\bigr]=0
    \]
    \[
        \begin{aligned}
            &\Rightarrow\lambda^{2}a'E\bigl[-u''(\tilde{w})\tilde{r}^{2}\bigr]=E\bigl[u''(\tilde{w})\tilde{w}\tilde{r}\bigr]=E\bigl[-u'(\tilde{w})R(\tilde{w})\tilde{r}\bigr]\\
            &\Rightarrow\operatorname{sign}(a')=-\operatorname{sign}\bigl(E\bigl[u'(\tilde{w})R(\tilde{w})\tilde{r}\bigr]\bigr)
        \end{aligned}
    \]
    因为
    \[
        E\bigl[u'(\tilde{w})R(\tilde{w})\tilde{r}\bigr]\ge E\bigl[u'(\tilde{w})R(w)\tilde{r}\bigr]=R(w)E\bigl[u'(\tilde{w})\tilde{r}\bigr]=0
    \]
    最后一个等式是因为一阶条件；因此 $a'\le 0$；
    \[
        \begin{aligned}
            &\Rightarrow a(\lambda w)\le a(w),\ \lambda>1\\
            &\Rightarrow a_{1}(w)\le a_{2}(w)
        \end{aligned}
    \]

    \item 令 $\lambda=\dfrac{w_{2}}{w_{1}}$，令 $u_{1}(w)=u(\lambda w)$，$u_{2}(w)=u(w)$，那么由 (a) 和 (b) 有 $\dfrac{a(w_{2})}{\lambda}\le a(w_{1})$，即 $\widehat{a}(w_{1})\ge\widehat{a}(w_{2})$。
\end{enumerate}
\end{solution}

\begin{problem}
\problemno{7.2} 参与者的初始财富为 $w$ 且他的绝对风险厌恶系数为常数 $a$。现在他必须承担风险 $\tilde{g}$，$\tilde{g}$ 是一公平赌博（即，他的财富变成 $w+\tilde{g}$）。当 $\tilde{g}$ 具有如下分布时计算确定性等价：
\begin{enumerate}[label=(\alph*)]
    \item 取值为 $-b$ 和 $b$ 的二项分布；
    \item 区间 $[-c,c]$ 上的均匀分布；
    \item 均值为 0、标准差为 $\sigma$ 的正态分布。
    \item 讨论在以上情形下，确定性等价如何依赖于初始财富 $w$。从中你能得到什么结论？并解释你的结论。
\end{enumerate}
\end{problem}
\begin{solution}
为简便起见，假设效用函数为 $U(w)=-e^{-aw}$，
\begin{enumerate}[label=(\alph*)]
    \item 由定义有
    \[
        \begin{aligned}
            U(\pi_{1}+w)
            &=-e^{-a(w+\pi_{1})}=E\bigl[U(w+\tilde{g})\bigr]=\tfrac12 U(w+b)+\tfrac12 U(w-b)\\
            &=-\tfrac12 e^{-a(w+b)}-\tfrac12 e^{-a(w-b)}\\
            \Rightarrow\pi_{1}&=-\frac{1}{a}\log\bigl(\tfrac12 e^{-ab}+\tfrac12 e^{ab}\bigr)
        \end{aligned}
    \]

    \item 由定义有
    \[
        \begin{aligned}
            U(w+\pi_{2})
            &=-e^{-a(w+\pi_{2})}=E\bigl[U(w+\tilde{g})\bigr]=\frac{1}{2c}\int_{-c}^{c}-e^{-a(w+x)}\,dx\\
            &=-\frac{1}{2ac}e^{-aw}(e^{ac}-e^{-ac})\\
            \Rightarrow\pi_{2}&=-\frac{1}{a}\log\Bigl(\frac{e^{ac}-e^{-ac}}{2ac}\Bigr)
        \end{aligned}
    \]

    \item 由定义有
    \[
        \begin{aligned}
            U(w+\pi_{3})
            &=-e^{-a(w+\pi_{3})}=E\bigl[U(w+\tilde{g})\bigr]\\
            &=\int_{-\infty}^{\infty}-e^{-a(w+x)}\frac{1}{\sqrt{2\pi\sigma}}\exp\Bigl(-\frac{x^{2}}{2\sigma^{2}}\Bigr)dx\\
            &=-e^{-aw+\frac{a^{2}\sigma^{2}}{2}}\\
            \Rightarrow\pi_{3}&=-\frac{a\sigma^{2}}{2}
        \end{aligned}
    \]

    \item 在以上各种情形下，确定性等价都不依赖于初始财富 $w$，也就是说，参与者对待公平赌博的态度与财富无关。
\end{enumerate}
\end{solution}

\begin{problem}
\problemno{7.3} 参与者的初始财富为 $w$ 且他的绝对风险厌恶系数为 1。现在他必须承担与他的财富成比例的风险 $w\tilde{g}$，$w\tilde{g}$ 是一公平赌博。当 $\tilde{g}$ 具有如下分布时，计算确定性等价占初始财富的比例——即相对确定性等价：
\begin{enumerate}[label=(\alph*)]
    \item 取值为 $-b$ 和 $b$ 的二项分布，其中 $0<b<1$；
    \item 区间 $[-c,c]$ 上的均匀分布，其中 $0<c<1$；
    \item 讨论在以上情形下，相对确定性等价如何依赖于初始财富 $w$。从中你能得到什么结论？并解释你的结论。
\end{enumerate}
\end{problem}
\begin{solution}
为简便起见，假设效用函数为 $U(w)=\log(w)$，
\begin{enumerate}[label=(\alph*)]
    \item 由定义有
    \[
        \begin{aligned}
            U(w\pi_{1}+w)&=\log(w+w\pi_{1})=E\bigl[U(w+w\tilde{g})\bigr]\\
            &=\tfrac12 U(w+b)+\tfrac12 U(w-b)=\log w+\tfrac12\log(1-b^{2})\\
            \Rightarrow\pi_{1}&=\sqrt{1-b^{2}}-1
        \end{aligned}
    \]

    \item 由定义有
    \[
        \begin{aligned}
            U(w+w\pi_{2})
            &=\log(w+w\pi_{2})=E\bigl[U(w+w\tilde{g})\bigr]=\frac{1}{2c}\int_{-c}^{c}\log(w+wx)\,dx\\
            &=\log w+\frac{1}{2c}\int_{-c}^{c}\log(1+x)\,dx\\
            \Rightarrow\pi_{2}&=\exp\Bigl(\frac{1}{2c}\int_{-c}^{c}\log(1+x)\,dx\Bigr)-1
        \end{aligned}
    \]

    \item 在以上各种情形下，确定性等价都不依赖于初始财富 $w$，也就是说，参与者对待与初始财富成比例的公平赌博的态度与初始财富无关。
\end{enumerate}
\end{solution}

\begin{problem}
\problemno{7.4} 参与者的初始财富为 $w$ 且相对风险厌恶系数为 2。现在他要承担风险 $\tilde{g}$，$\tilde{g}$ 是有两个可能取值 $-b$ 和 $b$ 的公平赌博，$0<b<w$。我们能用两种方式定义风险溢价：
\begin{itemize}
    \item $E[u(w+\tilde{g})]=u(w-\pi)$；
    \item $E[u(w+\tilde{g}+\pi)]=u(w)$。
\end{itemize}
用这两种定义计算风险溢价。它们相等吗？解释所得到的结论。
\end{problem}
\begin{solution}
不妨设 $U(w)=w^{\frac12}$，

第一种定义：
\[
    \begin{aligned}
        \sqrt{w-\pi_{1}}&=\tfrac12\sqrt{w+b}+\tfrac12\sqrt{w-b}\\
        \Rightarrow\pi_{1}&=\tfrac12\bigl(w-\sqrt{w^{2}-b^{2}}\bigr)
    \end{aligned}
\]

第二种定义：
\[
    \begin{aligned}
        \sqrt{w+\pi_{2}+\tilde{g}}&=\tfrac12\sqrt{w+b+\pi_{2}}+\tfrac12\sqrt{w-b+\pi_{2}}\\
        \Rightarrow\pi_{2}&=\frac{b^{2}}{4w}
    \end{aligned}
\]

两种定义所得的结果不一样，$\pi_{1}>\pi_{2}$，这是因为定义中的初始财富不一样，而效用函数的绝对风险厌恶系数又不是常数，因此得到的风险溢价不一样。
\end{solution}

\begin{problem}
\problemno{7.5*} 参与者的初始财富为 $w$ 且是严格风险厌恶的。$\tilde{g}_{1}$ 和 $\tilde{g}_{2}$ 是两个独立的公平赌博，具有相同的、定义于 $\{-b,b\}$ 上的二项分布，其中 $0<b<w/2$。
\begin{enumerate}[label=(\alph*)]
    \item 假设他必须承担风险 $\tilde{g}_{1}$，即他的财富变成 $w+\tilde{g}_{1}$。在这种情况下，他的期望效用记为 $V_{1}=E[u(w+\tilde{g}_{1})]$。证明风险使得他的情况恶化，即 $V_{1}<u(w)$。
    \item 现在假设他必须承担的是分散化的风险 $\tilde{g}=\frac12(\tilde{g}_{1}+\tilde{g}_{2})$。记这种情形下的期望效用为 $V_{2}=E[u(w+\tilde{g})]$。$V_{2}$ 是否总是高于 $V_{1}$？
\end{enumerate}
\end{problem}
\begin{solution}
\begin{enumerate}[label=(\alph*)]
    \item 因为参与者是严格风险厌恶的，因而效用函数 $u(\cdot)$ 是严格凹的，
    \[
        V_{1}=E[u(w+\tilde{g}_{1})]=\tfrac12 u(w+b)+\tfrac12 u(w-b)<u(w)
    \]

    \item
    \[
        \begin{aligned}
            V_{2}
            &=E[u(w+\tilde{g})]=\tfrac14 u(w+b)+\tfrac12 u(w)+\tfrac14 u(w-b)\\
            &>\tfrac14 u(w+b)+\tfrac12\bigl(\tfrac12 u(w+b)+\tfrac12 u(w-b)\bigr)+\tfrac14 u(w-b)=V_{1}
        \end{aligned}
    \]
\end{enumerate}
\end{solution}



% !TeX root = ../../金融经济学习题讲义.tex
% 源 PDF：references/金融经济学/史上最全最完整的金融经济学答案(王江版).pdf，页 32–36（印刷页 31–35）

\zhchapter{组合选择（一）}

\ifdefined\WJHandoutKnowledge
\fi

% !TeX root = ../../金融经济学习题讲义.tex
% 08_组合选择一 课后习题（由章节文件拆分）

\begin{problem}
\problemno{8.1} 假设有 $N$ 项资产，收益率为 $\tilde{r}_{i}$，$i=1,\cdots,N$，它们是独立同分布的。对于一个具有不满足性和凹性效用函数的投资者，计算他的最优投资组合。
\end{problem}
\begin{solution}
设组合权重 $\theta=(\theta_{1},\ldots,\theta_{N})$，$\sum_{i}\theta_{i}=1$，期末财富 $\tilde{w}=w\sum_{i}\theta_{i}\tilde{r}_{i}$。最大化 $E\bigl[u(\tilde{w})\bigr]$ 的一阶条件是：对任意 $j,k$，
\[
    E\bigl[u'(\tilde{w})(\tilde{r}_{j}-\tilde{r}_{k})\bigr]=0
\]
由于 $\{\tilde{r}_{i}\}$ 独立同分布，目标函数对 $\theta$ 对称；严格凹效用下最优解唯一，故必有 $\theta_{i}=1/N$。即最优组合是所有风险资产的等权重组合。
\end{solution}

\begin{problem}
\problemno{8.2*} 假设你必须将 $w$ 投资于两只证券：一只无风险债券和一只风险证券。无风险债券的收益率确定为 $r_{F}$。风险证券的收益率为 $\tilde{r}$，它的均值为 $\bar{r}$ 而方差为 $\sigma^{2}$。假设你的目标就是最大化你的二次效用函数的期望值：
\[
    E\Bigl[\tilde{w}-\tfrac12 a\tilde{w}^{2}\Bigr].
\]
\begin{enumerate}[label=(\alph*)]
    \item 求解最优组合。
    \item 讨论最优组合如何依赖于股票和债券的期望收益率 $r_{F}$ 和 $\bar{r}$，股票收益率的波动率 $\sigma$，以及偏好系数 $a$。为这些依赖关系提供经济解释。
\end{enumerate}
\end{problem}
\begin{solution}
\begin{enumerate}[label=(\alph*)]
    \item 假设投资于股票的资产比例为 $x$，那么有
    \[
        \tilde{w}=w\bigl(1+r_{F}+x(\tilde{r}-r_{F})\bigr)
    \]
    参与者的优化问题是
    \[
        \max_{x}E\Bigl[\tilde{w}-\tfrac12 a\tilde{w}^{2}\Bigr]
    \]
    记 $\mu=\bar{r}-r_{F}$。则 $E[\tilde{w}]=w(1+r_{F}+x\mu)$，$E[\tilde{w}^{2}]=w^{2}\bigl[(1+r_{F}+x\mu)^{2}+x^{2}\sigma^{2}\bigr]$。对 $x$ 求导，一阶条件为
    \[
        w\mu-aw^{2}\bigl[(1+r_{F}+x\mu)\mu+x\sigma^{2}\bigr]=0
    \]
    即 $\mu\bigl[1-aw(1+r_{F})\bigr]=aw\bigl(\mu^{2}+\sigma^{2}\bigr)x$，解得
    \[
        x=\frac{(\bar{r}-r_{F})\bigl(1-aw(1+r_{F})\bigr)}{aw\bigl(\sigma^{2}+(\bar{r}-r_{F})^{2}\bigr)}
    \]

    \item $x$ 随着 $\bar{r}$ 的增加而增加，随着 $r_{F}$、$\sigma^{2}$、$a$ 的增加而减小。股票的期望收益率增加使得股票的需求增加，而无风险收益率增加使得债券的需求增加；$\sigma^{2}$ 增大，股票的风险变大，对股票的需求变小；最后，$a$ 增大表明参与者更加厌恶风险，因而对股票的需求减小。
\end{enumerate}
\end{solution}

\begin{problem}
\problemno{8.3*} 考虑与练习 8.2 一样的问题。不过现在假设效用函数具有常数绝对风险厌恶系数：
\[
    E\bigl[-e^{-a\tilde{w}}\bigr].
\]
另外，假设股票收益率 $\tilde{r}$ 是正态分布的。
\begin{enumerate}[label=(\alph*)]
    \item 求解最优投资组合。
    \item 讨论最优组合如何依赖于股票和债券的期望收益率 $r_{F}$ 和 $\bar{r}$，股票收益率的波动率 $\sigma$，以及偏好系数 $a$。为这些依赖性提供经济解释。[如果需要的话，你可以添加其它条件。]
\end{enumerate}
\end{problem}
\begin{solution}
\begin{enumerate}[label=(\alph*)]
    \item 假设投资于股票的资产比例为 $x$，那么有
    \[
        \tilde{w}=w\bigl(1+r_{F}+x(\tilde{r}-r_{F})\bigr)
    \]
    参与者的优化问题是
    \[
        \max_{x}E\bigl[-e^{-a\tilde{w}}\bigr]
    \]
    因 $\tilde{r}$ 正态，$\tilde{w}$ 亦正态。对正态随机变量有
    \[
        E\bigl[-e^{-a\tilde{w}}\bigr]=-\exp\Bigl(-aE[\tilde{w}]+\tfrac12 a^{2}\operatorname{Var}(\tilde{w})\Bigr)
    \]
    其中 $E[\tilde{w}]=w\bigl(1+r_{F}+x(\bar{r}-r_{F})\bigr)$，$\operatorname{Var}(\tilde{w})=w^{2}x^{2}\sigma^{2}$。最大化上式等价于最大化确定性等价 $E[\tilde{w}]-\dfrac{a}{2}\operatorname{Var}(\tilde{w})$。对 $x$ 的一阶条件为
    \[
        w(\bar{r}-r_{F})-a\,w^{2}x\sigma^{2}=0
    \]
    因而最优权重（CARA 下风险资产的金额 $wx$ 与财富无关）为
    \[
        x=\frac{\bar{r}-r_{F}}{aw\sigma^{2}}
    \]
    \item $x$ 随着 $\bar{r}$ 的增加而增加，随着 $r_{F}$、$\sigma^{2}$、$a$ 以及财富 $w$ 的增加而减小。超额收益上升提高股票需求；$\sigma^{2}$ 或绝对风险厌恶 $a$ 增大则降低股票需求。与二次效用不同：CARA 下风险资产金额 $wx=(\bar{r}-r_{F})/(a\sigma^{2})$ 不随 $w$ 变化，故权重 $x$ 随财富下降。
\end{enumerate}
\end{solution}

\begin{problem}
\problemno{8.4**} 赋予投资者初始金融财富 $w$ 以及 $\lambda$ 份不可交易收入 $\tilde{y}$，其中 $\lambda\ge 0$ 且 $E[\tilde{y}]=0$。他可以交易一只无风险证券和一只风险证券。无风险证券的收益率为 0 而风险证券的收益率 $\tilde{r}$ 的均值为 $\bar{r}$。他具有不满足性且是风险厌恶的。他的绝对风险厌恶系数随着他的财富增加而减少。令 $a$ 为风险证券上的投资额。他的期望效用函数为：
\[
    \max E\bigl[u(w+a\tilde{r}+\lambda\tilde{y})\bigr].
\]
\begin{enumerate}[label=(\alph*)]
    \item 当 $\lambda=0$ 时，风险资产的最优头寸 $a_{0}^{*}$ 如何随着 $w$ 变化而变化？
    \item 假设 $\tilde{y}$ 和 $\tilde{r}$ 是独立的。当 $\lambda>0$ 时，记风险资产的最优头寸为 $a_{\lambda}^{*}$。$a_{\lambda}^{*}$ 是大于还是小于 $a_{0}^{*}$？
    \item 假设 $\tilde{y}=c(\tilde{r}-\bar{r})+\tilde{\varepsilon}$ 以及 $E[\tilde{\varepsilon}\mid\tilde{r}]=0$。因此，如果 $c\neq 0$，$\tilde{y}$ 和 $\tilde{r}$ 是相关的。当 $\lambda$ 极小时，他持有的股票头寸如何随着 $w$ 变化而变化？解释所得到的结论。
\end{enumerate}
\end{problem}
\begin{solution}
\begin{enumerate}[label=(\alph*)]
    \item 当 $\lambda=0$ 时，优化的一阶条件是
    \[
        E\bigl[u'(w+a_{0}^{*}\tilde{r})\tilde{r}\bigr]=0
    \]
    两边对 $w$ 求导可得
    \[
        \frac{da_{0}^{*}}{dw}=\frac{E[u''(\tilde{w})\tilde{r}]}{-E[u''(\tilde{w})\tilde{r}^{2}]}
    \]
    又
    \[
        \begin{aligned}
            E[u''(\tilde{w})\tilde{r}]
            &=-E[u'(\tilde{w})R(\tilde{w})\tilde{r}]\\
            &\ge -E[u'(\tilde{w})R(w)\tilde{r}]\\
            &=-R(w)E[u'(\tilde{w})\tilde{r}]\\
            &=0
        \end{aligned}
    \]
    故 $\dfrac{da_{0}^{*}}{dw}\ge 0$，即最优头寸随着 $w$ 的增加而增加；

    \item 参与者优化的一阶条件是
    \[
        E\bigl[u'(w+a_{\lambda}^{*}\tilde{r}+\lambda\tilde{y})\tilde{r}\bigr]=0
    \]
    两边对 $\lambda$ 求导可得
    \[
        \frac{da_{\lambda}^{*}}{d\lambda}=\frac{E[u''(\tilde{w})\tilde{r}\tilde{y}]}{-E[u''(\tilde{w})\tilde{r}^{2}]}
    \]
    记 $K(w+a_{\lambda}^{*}\tilde{r})=E[-R(\tilde{w})\tilde{y}\mid\tilde{r}]$，可以证明 $K(\cdot)>0$，$K'(\cdot)>0$，因此
    \[
        E[u''(\tilde{w})\tilde{r}\tilde{y}]=E[-R(\tilde{w})u'(\tilde{w})\tilde{r}\tilde{y}]=E\bigl[K(w+a_{\lambda}^{*}\tilde{r})u'(\tilde{w})\tilde{r}\bigr]\ge K(w)E[u'(\tilde{w})\tilde{r}]=0
    \]
    最后一个等式是因为一阶条件，因此，$\dfrac{da_{\lambda}^{*}}{d\lambda}\ge 0$，故 $a_{\lambda}^{*}\ge a_{0}^{*}$；

    \item 由 (a) 有
    \[
        \frac{da_{0}^{*}}{dw}=\frac{E[u''(\tilde{w})\tilde{r}]}{-E[u''(\tilde{w})\tilde{r}^{2}]}=\frac{-E[u'(\tilde{w})R(\tilde{w})\tilde{r}]}{-E[u''(\tilde{w})\tilde{r}^{2}]}
    \]
    而
    \[
        E[u'(\tilde{w})R(\tilde{w})\tilde{r}]=E\bigl[R(w+a\tilde{r})u'(\tilde{w})\tilde{r}\bigr]+\lambda E\bigl[R'(w+a\tilde{r})u'(\tilde{w})\tilde{r}\tilde{y}\bigr]+o(\lambda^{2})
    \]
    由 (a) 知道 $E\bigl[R(w+a\tilde{r})u'(\tilde{w})\tilde{r}\bigr]\le 0$，另外 $\lambda$ 又极小，因此
    \[
        E[u'(\tilde{w})R(\tilde{w})\tilde{r}]\le 0,\qquad \frac{da^{*}}{dw}\ge 0
    \]
    即最优头寸随着初始财富 $w$ 的增加而增加。这是因为 $\lambda$ 极小，因而不可交易收入不会改变 (a) 中的结论。
\end{enumerate}
\end{solution}

\begin{problem}
\problemno{8.5**} 假设你是对冲基金中交易政府债的交易员。你有 $w$ 的现金额度。
\begin{itemize}
    \item 在 0 期，当美联储拍卖新债券时，你可以提交报价而买入政府债。在 1 期，你可以在二级市场上以价格 $\tilde{v}$ 再把它们卖掉，其中 $\tilde{v}=\bar{v}+\tilde{\varepsilon}$，$\bar{v}>0$ 且 $E[\tilde{\varepsilon}]=0$。假设无风险收益率为零。
    \item 你购买债券的报价是 $p$。你购买到的数量 $x(p)$ 是 $p$ 的递增函数且 $x(0)=0$。选择 $p$ 以最大化你的期望效用函数：
    \[
        V(p)=E[u(\tilde{w})],\qquad \tilde{w}=w+x(p)(\bar{v}+\tilde{\varepsilon}-p).
    \]
    假设你是不满足的且严格风险厌恶的。为简单起见，假设 $V(p)$ 是严格凹的。
\end{itemize}
\begin{enumerate}[label=(\alph*)]
    \item 写出关于最优报价 $p^{*}$ 的一阶条件。
    \item 证明 $0<p^{*}<\bar{v}$。
    \item 对于两个绝对风险厌恶系数分别为 $A_{1}(w)$ 和 $A_{2}(w)$ 的交易员，如果 $\forall\,w$，$A_{1}(w)\ge A_{2}(w)$，那么 $p_{1}^{*}\le p_{2}^{*}$，其中 $p_{i}^{*}$ 是交易员 $i$ 的最优报价，$i=1,2$。
    \item 假设 $u(\cdot)$ 是绝对风险厌恶递减的。证明初始财富 $w$ 的增加使得 $p^{*}$ 也增加。
\end{enumerate}
\end{problem}
\begin{solution}
\begin{enumerate}[label=(\alph*)]
    \item 一阶条件是
    \[
        V'(p^{*})=E\bigl[u'(\tilde{w})\bigl(x'(p^{*})(\bar{v}+\tilde{\varepsilon}-p^{*})-x(p^{*})\bigr)\bigr]=0
    \]

    \item
    \[
        \begin{aligned}
            V'(0)&=E\bigl[u'(w)\bigl(x'(0)(\bar{v}+\tilde{\varepsilon})\bigr)\bigr]=u'(w)x'(0)\bar{v}>0\\
            V(\bar{v})&=E\bigl[u(w+x(\bar{v})\tilde{\varepsilon})\bigr]<u(w)=V(0)
        \end{aligned}
    \]
    又因为 $V(p)$ 是凹的，故 $0<p^{*}<\bar{v}$；

    \item 由定理 (7.4) 可知，存在 $f'>0$，$f''<0$ 使得 $u_{1}=f(u_{2})$。参与者的优化问题的一阶条件为：
    \[
        E\bigl[u_{2}'(\tilde{w}_{2})\bigl(x'(p_{2})(\bar{v}+\tilde{\varepsilon}-p_{2})-x(p_{2})\bigr)\bigr]=0
    \]
    考虑
    \[
        V_{1}'(p_{2})=E\bigl[f'(u_{2}(\tilde{w}_{2}))u_{2}'(\tilde{w}_{2})\bigl(x'(p_{2})(\bar{v}+\tilde{\varepsilon}-p_{2})-x(p_{2})\bigr)\bigr]
    \]
    令 $\tilde{r}=x'(p_{2})(\bar{v}+\tilde{\varepsilon}-p_{2})-x(p_{2})$，那么 $\tilde{w}_{2}$ 是 $\tilde{r}$ 的增函数，又 $f'(u_{2}(\tilde{w}_{2}))>0$ 且是 $\tilde{w}_{2}$ 的减函数，因此，$V_{1}'(p_{2})\le 0$，另外 $V(p)$ 是凹函数，因此 $p_{1}^{*}\le p_{2}^{*}$。

    \item 假设 $w_{1}>w_{0}$，当 $w=w_{0}$ 时优化问题的一阶条件为
    \[
        E\bigl[u'\bigl(w_{0}+x(p_{0})(\bar{v}+\tilde{\varepsilon}-p_{0})\bigr)\bigl(x'(p_{0})(\bar{v}+\tilde{\varepsilon}-p_{0})-x(p_{0})\bigr)\bigr]=0
    \]
    等式两边对 $w_{0}$ 求导可得
    \[
        \begin{aligned}
            &E\bigl[u''(\tilde{w}_{0})\bigl(x'(p_{0})(\bar{v}+\tilde{\varepsilon}-p_{0})-x(p_{0})\bigr)\bigr]\\
            &=E\bigl[-R(\tilde{w}_{0})u'(\tilde{w}_{0})\bigl(x'(p_{0})(\bar{v}+\tilde{\varepsilon}-p_{0})-x(p_{0})\bigr)\bigr]\\
            &\ge 0
        \end{aligned}
    \]
    因此，
    \[
        E\bigl[u'\bigl(w_{1}+x(p_{0})(\bar{v}+\tilde{\varepsilon}-p_{0})\bigr)\bigl(x'(p_{0})(\bar{v}+\tilde{\varepsilon}-p_{0})-x(p_{0})\bigr)\bigr]\ge 0
    \]
    而 $V(p)$ 又是凹的，因此 $p_{1}\ge p_{0}$，即初始财富增加报价也会增加。
\end{enumerate}
\end{solution}



% !TeX root = ../../金融经济学习题讲义.tex
% 源 PDF：references/金融经济学/史上最全最完整的金融经济学答案(王江版).pdf，页 38–42（印刷页 37–41）

\zhchapter{组合选择（二）}

\ifdefined\WJHandoutKnowledge
\fi

% !TeX root = ../../金融经济学习题讲义.tex
% 09_组合选择二 课后习题（由章节文件拆分）

\begin{problem}
\problemno{9.1} 假设有四个等概率的可能状态。状态空间是 $\Omega=\{\omega_{1},\cdots,\omega_{4}\}$。考虑风险资产 $A$ 和 $B$，它们的收益率如下：
\begin{center}
\begin{tabular}{c|cccc}
    & $\omega_{1}$ & $\omega_{2}$ & $\omega_{3}$ & $\omega_{4}$ \\
    \hline
    $\tilde{r}_{A}$ & 0.5 & 0.5 & 0.7 & 0.7 \\
    $\tilde{r}_{B}$ & 0.9 & 0.8 & 0.4 & 0.3
\end{tabular}
\end{center}
在随机占优的意义上，你能为它们排序吗？
\end{problem}
\begin{solution}
\[
    F_{A}(x)=\begin{cases}
        0 & 0\le x<0.5\\
        0.5 & 0.5\le x<0.7\\
        1 & 0.7\le x
    \end{cases}
\]
\[
    F_{B}(x)=\begin{cases}
        0 & 0\le x<0.3\\
        0.25 & 0.3\le x<0.4\\
        0.5 & 0.4\le x<0.8\\
        0.75 & 0.8\le x<0.9\\
        1 & 0.9\le x
    \end{cases}
\]
\[
    S(y)=\int_{0}^{y}\bigl(F_{A}(x)-F_{B}(x)\bigr)\,dx=\begin{cases}
        0 & 0\le y<0.3\\
        -0.25(y-0.3) & 0.3\le y<0.4\\
        -0.025-0.5(y-0.4) & 0.4\le y<0.5\\
        -0.075 & 0.5\le y<0.7\\
        -0.075+0.5(y-0.7) & 0.7\le y<0.8\\
        -0.025+0.25(y-0.8) & 0.8\le y<0.9\\
        0 & 0.9\le y
    \end{cases}
\]
可以验证 $S(y)\le 0$，$y\in[0,1]$；且 $E[x_{A}]=E[x_{B}]$，因此 $A\succeq_{\mathrm{SSD}}B$。
\end{solution}

\begin{problem}
\problemno{9.2*} 两只风险资产的收益率分别为 $\tilde{r}_{A}$ 和 $\tilde{r}_{B}$。考虑一个投资者，他的初始财富为 $w=1$、效用函数为 $u(\cdot)$，其中 $u'(\cdot)>0$ 且 $u''(\cdot)<0$。他最大化期望效用 $E[u(\tilde{w})]$。令他的最优组合为 $[a,1-a]$。
\begin{enumerate}[label=(\alph*)]
    \item 假设 $\tilde{r}_{A}=\tilde{x}$，$\tilde{r}_{B}=\tilde{x}+\tilde{e}$，其中 $\tilde{e}$ 和 $\tilde{x}$ 互相独立且 $E[\tilde{e}]=0$。
    \begin{itemize}
        \item $\tilde{r}_{A}$ 是否二阶随机占优于 $\tilde{r}_{B}$？
        \item 对于投资者的最优投资组合 $a$，你能得到什么样的结论？
    \end{itemize}
    \item 假设 $\tilde{r}_{A}=\tilde{x}$，$\tilde{r}_{B}=\tilde{y}+\tilde{e}$，$\tilde{x}$、$\tilde{y}$、$\tilde{e}$ 互相独立，$\tilde{x}\sim_{d}\tilde{y}$ 且 $E[\tilde{e}]=0$。
    \begin{itemize}
        \item $\tilde{r}_{A}$ 是否二阶随机占优于 $\tilde{r}_{B}$？
        \item 证明投资者在资产 $A$ 上的投资额要大于他资产 $B$ 上的投资额，即 $a>1/2$。
    \end{itemize}
    \item 解释上面两种情形下结论之间的差异。
\end{enumerate}
\end{problem}
\begin{solution}
\begin{enumerate}[label=(\alph*)]
    \item 根据定义可以判断出
    \[
        A\succeq_{\mathrm{SSD}}B
    \]
    \[
        \begin{aligned}
            E[u(\tilde{w})]
            &=E\bigl[u\bigl(\tilde{x}+(1-a)\tilde{e}\bigr)\bigr]\\
            &=E\bigl[E\bigl[u\bigl(\tilde{x}+(1-a)\tilde{e}\bigr)\bigm|\tilde{x}\bigr]\bigr]\\
            &\le E\bigl[u\bigl(E\bigl[\tilde{x}+(1-a)\tilde{e}\bigm|\tilde{x}\bigr]\bigr)\bigr]\\
            &=E\bigl[u(\tilde{x})\bigr]\\
            &=E\bigl[u(\tilde{r}_{A})\bigr]
        \end{aligned}
    \]
    当 $a=1$ 时等式成立，因此 $a=1$。

    \item 因为 $\tilde{x}$，$\tilde{y}$ 是独立同分布的，因而 $\tilde{r}_{B}\sim_{d}\tilde{r}_{A}+\tilde{e}$，且 $E[\tilde{e}\mid\tilde{r}_{A}]=0$，故 $A\succeq_{\mathrm{SSD}}B$。令 $V(a)=E\bigl[u\bigl(a\tilde{x}+(1-a)\tilde{y}+(1-a)\tilde{e}\bigr)\bigr]$，那么
    \[
        \begin{aligned}
            V'(a)&=E\bigl[u'\bigl(a\tilde{x}+(1-a)\tilde{y}+(1-a)\tilde{e}\bigr)(\tilde{x}-\tilde{y}-\tilde{e})\bigr]\\
            V''(a)&=E\bigl[u''\bigl(a\tilde{x}+(1-a)\tilde{y}+(1-a)\tilde{e}\bigr)(\tilde{x}-\tilde{y}-\tilde{e})^{2}\bigr]\\
            V'\bigl(\tfrac12\bigr)&=E\bigl[u'\bigl(\tfrac12\tilde{x}+\tfrac12\tilde{y}+\tfrac12\tilde{e}\bigr)(\tilde{x}-\tilde{y}-\tilde{e})\bigr]\\
            &=E\bigl[E\bigl[u'\bigl(\tfrac12\tilde{x}+\tfrac12\tilde{y}+\tfrac12\tilde{e}\bigr)(\tilde{x}-\tilde{y}-\tilde{e})\bigm|\tilde{x}+\tilde{y}\bigr]\bigr]\\
            &=E\bigl[u'\bigl(\tfrac12\tilde{x}+\tfrac12\tilde{y}+\tfrac12\tilde{e}\bigr)E[\tilde{x}-\tilde{y}\mid\tilde{x}+\tilde{y}]-E\bigl[u'\bigl(\tfrac12\tilde{x}+\tfrac12\tilde{y}+\tfrac12\tilde{e}\bigr)\tilde{e}\bigm|\tilde{x}+\tilde{y}\bigr]\bigr]\\
            &=-E\bigl[u'\bigl(\tfrac12\tilde{x}+\tfrac12\tilde{y}+\tfrac12\tilde{e}\bigr)\tilde{e}\bigr]\\
            &=-\operatorname{Cov}\bigl(u'\bigl(\tfrac12\tilde{x}+\tfrac12\tilde{y}+\tfrac12\tilde{e}\bigr),\tilde{e}\bigr)>0
        \end{aligned}
    \]
    故，$a^{*}>\frac12$，同样的方法可以验证 $a^{*}<1$。

    \item 因为在第一种情形下，$\tilde{r}_{B}$ 只是比 $\tilde{r}_{A}$ 多一个噪音项 $\tilde{e}$，引入 $\tilde{r}_{B}$ 的效果是增加了收益中的噪音，因此 $a=1$；而在第二种情形下，因为 $\tilde{x}$，$\tilde{y}$ 是独立同分布的，引入 $\tilde{r}_{B}$ 可以分散风险，但由于 $\tilde{r}_{B}$ 比 $\tilde{r}_{A}$ 多一个噪音项，因此 $1>a^{*}>\frac12$。
\end{enumerate}
\end{solution}

\begin{problem}
\problemno{9.3} 证券市场最初由 $N$ 只证券构成。证明：当把另外 $M$ 只证券添加到市场中去时，如果投资者的最优投资组合中包含新证券，那么他的福利（即他的期望效用）会增加（也就是说，更多选择总是好的）。
\end{problem}
\begin{solution}
因为在添加了新证券以后，投资者总是可以选择原来的最优组合，只有当他的效用增加时他才会选择新的证券去构成新的组合。
\end{solution}

\begin{problem}
\problemno{9.4} 假设有 $N$ 只证券，它们的供给都是无限的、收益率为 $\tilde{r}_{i}$，$i=1,\cdots,N$。经济中的参与者都是不满足的、风险中性的且具有正的禀赋。市场不允许卖空。在什么条件下 $K$ 基金分离成立（$K<N$）？
\end{problem}
\begin{solution}
当且仅当期望收益率最大的证券数等于 $K$ 时，$K$ 基金分离定理成立。
\end{solution}

\begin{problem}
\problemno{9.5} 假设在市场上有 $N$ 只证券，它们的收益率是独立同分布的。证明单基金分离成立。
\end{problem}
\begin{solution}
由练习 8.1：收益率独立同分布时，严格凹效用下最优组合必为等权重。所有投资者持有同一风险基金，故单基金分离成立。
\end{solution}

\begin{problem}
\problemno{9.6*} 证券市场由一只无风险债券和 $N$ 只风险证券组成。无风险证券的收益率为 $r_{F}$ 而风险证券的收益率向量为 $\tilde{r}=\begin{bmatrix}\tilde{r}_{1} \\ \ldots \\ \tilde{r}_{N}\end{bmatrix}$。另外，$E[\tilde{r}]=\bar{r}=\begin{bmatrix}\bar{r}_{1} \\ \ldots \\ \bar{r}_{N}\end{bmatrix}$ 且 $E[(\tilde{r}-\bar{r})(\tilde{r}-\bar{r})^{\top}]=\Sigma$（$\Sigma$ 是风险证券收益率的协方差阵）。考虑一个财富为 $w$ 的参与者的投资问题。假设他具有二次效用函数 $u(\tilde{w})=E[\tilde{w}]-\tfrac12 a E[\tilde{w}^{2}]$。
\begin{enumerate}[label=(\alph*)]
    \item 求解他的最优投资组合。
    \item 证明对于具有二次效用函数的参与者，两基金货币分离成立。
\end{enumerate}
\end{problem}
\begin{solution}
\begin{enumerate}[label=(\alph*)]
    \item 令风险资产权重为 $x$，并把期末财富写成与练习 8.2 相同的形式
    \[
        \tilde{w}=w\bigl(1+r_{F}+x^{\top}(\tilde{r}-r_{F}\iota)\bigr)
    \]
    记 $\mu=\bar{r}-r_{F}\iota$。则 $E[\tilde{w}]=w(1+r_{F}+x^{\top}\mu)$，且
    \[
        E[\tilde{w}^{2}]=w^{2}\bigl[(1+r_{F}+x^{\top}\mu)^{2}+x^{\top}\Sigma x\bigr]
    \]
    对 $x$ 求导，一阶条件为
    \[
        w\mu-aw^{2}\bigl[(1+r_{F}+x^{\top}\mu)\mu+\Sigma x\bigr]=0
    \]
    即 $\bigl[1-aw(1+r_{F})\bigr]\mu=aw(\mu\mu^{\top}+\Sigma)x$，解得
    \[
        x=\frac{1-aw(1+r_{F})}{aw}\bigl[(\bar{r}-r_{F}\iota)(\bar{r}-r_{F}\iota)^{\top}+\Sigma\bigr]^{-1}(\bar{r}-r_{F}\iota)
    \]
    其中，$\iota$ 是元素全为 1 的 $N$ 维列向量。风险组合中各资产的相对权重由 $(\mu\mu^{\top}+\Sigma)^{-1}\mu$ 决定，与财富、风险厌恶的水平只差一个共同标量。

    \item 因为投资者持有的风险组合的权重之比恒为常数，因而对于所有的投资者来说，风险组合是相同的，只是投资于风险组合的总金额不一样而已；另外，投资者还持有无风险资产，因此两基金分离成立。
\end{enumerate}
\end{solution}

\begin{problem}
\problemno{9.7*} 考虑练习 9.6 中描述的市场。另外，我们假设风险收益率是联合正态分布的且参与者的绝对风险厌恶系数恒为 $a$。
\begin{enumerate}[label=(\alph*)]
    \item 求解他的最优投资组合。
    \item 证明对于具有常数绝对风险厌恶的参与者，两基金货币分离成立。
\end{enumerate}
\end{problem}
\begin{solution}
\begin{enumerate}[label=(\alph*)]
    \item 仍令风险资产权重为 $x$，
    \[
        \tilde{w}=w\bigl(1+r_{F}+x^{\top}(\tilde{r}-r_{F}\iota)\bigr)
    \]
    $\tilde{r}$ 联合正态则 $\tilde{w}$ 正态。与练习 8.3 相同，
    \[
        E\bigl[-e^{-a\tilde{w}}\bigr]=-\exp\Bigl(-aE[\tilde{w}]+\tfrac12 a^{2}\operatorname{Var}(\tilde{w})\Bigr)
    \]
    其中 $E[\tilde{w}]=w(1+r_{F}+x^{\top}\mu)$，$\operatorname{Var}(\tilde{w})=w^{2}x^{\top}\Sigma x$，$\mu=\bar{r}-r_{F}\iota$。对 $x$ 的一阶条件为
    \[
        w\mu-aw^{2}\Sigma x=0
    \]
    因而
    \[
        x=\frac{1}{aw}\Sigma^{-1}(\bar{r}-r_{F}\iota)
    \]
    其中，$\iota$ 是元素全为 1 的 $N$ 维列向量。相对权重由 $\Sigma^{-1}\mu$ 决定，与 $a$、$w$ 无关。

    \item 因为投资者持有的风险组合的权重之比恒为常数，因而对于所有的投资者来说，风险组合是相同的，只是投资于风险组合的总金额不一样而已；另外，投资者还持有无风险资产，因此两基金分离成立。
\end{enumerate}
\end{solution}

\begin{problem}
\problemno{9.8*} 继续练习 9.7 中的问题。
\begin{enumerate}[label=(\alph*)]
    \item 计算在最优投资组合下参与者的期望效用。
    \item 假设现在有 $M$ 支新风险证券加入到证券市场中。（现有 $N$ 只证券的收益率不能生成新证券的收益率。）所有风险证券的收益率都是联合正态分布的。证明，如果参与者的最优投资组合中含有新的证券，无论 $M$ 只新证券的收益率分布如何，参与者的福利总是增加。
    \item 当引入新证券后，计算他的福利增加的确定性等价值。
\end{enumerate}
\end{problem}
\begin{solution}
\begin{enumerate}[label=(\alph*)]
    \item 由第七题我们知道最优组合的风险资产权重为 $x=\dfrac{1}{aw}\Sigma^{-1}(\bar{r}-r_{F}\iota)$，因而期望效用为
    \[
        \begin{aligned}
            E[u(\tilde{w})]
            &=-\exp\Bigl(-aE[\tilde{w}]+\tfrac12 a^{2}\operatorname{Var}(\tilde{w})\Bigr)\\
            &=-\exp\Bigl(-aw(1+r_{F})-aw\,x^{\top}(\bar{r}-r_{F}\iota)+\tfrac12 a^{2}w^{2}x^{\top}\Sigma x\Bigr)\\
            &=-\exp\Bigl(-aw(1+r_{F})-\tfrac12(\bar{r}-r_{F}\iota)^{\top}\Sigma^{-1}(\bar{r}-r_{F}\iota)\Bigr)
        \end{aligned}
    \]

    \item 由第三题可以得到此结论；

    \item 假设原来的风险资产的收益率和协方差矩阵分布为 $\bar{r}_{\mathrm{old}}$，$\Sigma_{\mathrm{old}}$，加入新证券的风险资产的收益率和协方差矩阵分布为 $\bar{r}_{\mathrm{new}}$，$\Sigma_{\mathrm{new}}$，那么福利增加的确定性等价值为
    \[
        \tfrac12(\bar{r}_{\mathrm{new}}-r_{F}\iota)^{\top}\Sigma_{\mathrm{new}}^{-1}(\bar{r}_{\mathrm{new}}-r_{F}\iota)-\tfrac12(\bar{r}_{\mathrm{old}}-r_{F}\iota)^{\top}\Sigma_{\mathrm{old}}^{-1}(\bar{r}_{\mathrm{old}}-r_{F}\iota)
    \]
\end{enumerate}
\end{solution}

\begin{problem}
\problemno{9.9*} 在练习 9.7 中，假设只有一只新证券（即 $M=1$），它的收益率是 $\tilde{r}_{N+1}$。令 $\bar{r}_{N+1}$ 为它的预期收益率、$\sigma_{N+1}$ 为收益率波动率而 $\sigma_{\cdot,N+1}$ 为它与 $N$ 只原有证券之间的协方差向量。另外，证券 $N+1$ 上的最优组合权重记作 $z_{N+1}$。
\begin{enumerate}[label=(\alph*)]
    \item 假设 $\bar{r}_{N+1}-r_{F}>(<)\,0$。我们能得到结论 $z_{N+1}>(<)\,0$ 吗？给出解释。
    \item 假设 $\bar{r}_{N+1}<\bar{r}_{n}$ 且 $\sigma_{N+1}>\sigma_{n}$，$\forall\,n=1,\ldots,N$。也就是说，其它所有证券在收益率和方差上都优于证券 $N+1$。这是否意味着 $z_{N+1}=0$ 即最优组合不应含有证券 $N+1$？给出解释。
\end{enumerate}
\end{problem}
\begin{solution}
\begin{enumerate}[label=(\alph*)]
    \item 由第七题的结论我们知道 $z_{N+1}=(\Sigma^{-1})_{N+1}(\bar{r}-r_{F}\iota)$。其中 $(\Sigma^{-1})_{N+1}$ 表示 $\Sigma^{-1}$ 的第 $N+1$ 行。在构造组合的时候，一方面我们需要考虑资产对组合收益的影响，另一方面我们还得考虑资产对组合风险的影响，如果 $\bar{r}_{N+1}-r_{F}>(<)\,0$，那么买入（卖出）证券 $N+1$ 会增加组合的收益，但同时可能会提高组合的风险，因而总体效应是不确定的，即不能得出 $z_{N+1}>(<)\,0$ 的结论。

    \item 单个资产的权重不但与自身的收益率、标准差有关，还与其它的资产的收益率以及它与其它资产的相关性有关，如果资产 $N+1$ 与其它资产是负（正）相关的，那么 $z_{N+1}$ 有可能大于（小于）0。
\end{enumerate}
\end{solution}



% !TeX root = ../../金融经济学习题讲义.tex
% 源 PDF：references/金融经济学/史上最全最完整的金融经济学答案(王江版).pdf，页 44–53（印刷页 43–52）

\chapter{完全市场中的资源配置与资产价格}

\begin{problem}
\problemno{10.1*} 考虑一个两期（$t=0,\,1$）、两个状态（$\omega=a,\,b$），存在 $K$ 个参与者的经济。状态 $a$ 发生的概率为 $\pi_{a}$ 而状态 $b$ 发生的概率为 $\pi_{b}$。所有参与者都是不满足的且严格风险厌恶的。他们期望效用有如下形式：
\[
    u_{k}(c_{k,0})+\pi_{a}u_{k}(c_{k,a})+\pi_{b}u_{k}(c_{k,b}),\quad k=1,\dots,K.
\]
参与者 $k$ 的禀赋是 $\begin{bmatrix}e_{k,0} \\ e_{k,a} \\ e_{k,b}\end{bmatrix}$。整个经济的总禀赋是 $\begin{bmatrix}C_{0} \\ C_{a} \\ C_{b}\end{bmatrix}$，其中 $C_{s}=\sum_{k}e_{k,s}$，$s=0,a,b$。假设 $C_{a}>C_{b}$。假设市场上交易的 Arrow--Debreu 状态或有要求权的集合是完全的。令 $\phi_{a}$ 和 $\phi_{b}$ 分别是状态 $a$ 和状态 $b$ 的状态价格。
\begin{enumerate}[label=(\alph*)]
    \item 定义这个经济的均衡。简要讨论均衡的存在性和唯一性。
    \item 令代表性参与者的（间接）期望效用函数为
    \[
        u(C_{0})+\pi_{a}u(C_{a})+\pi_{b}u(C_{b})
    \]
    且 $u'''(\cdot)<0$。当以下变量变化而其它变量保持不变时，状态价格（$\phi_{a}$ 和 $\phi_{b}$）、它们的比例（$\phi_{a}/\phi_{b}$）和无风险利率如何变化（上升或下降）？给出证明。
    \begin{itemize}
        \item $C_{0}$ 上升；
        \item $C_{a}$ 上升；
        \item $C_{b}$ 上升；
        \item $C_{a}$ 和 $C_{b}$ 同比例上升；
        \item $C_{a}-C_{b}$ 上升（禀赋的风险上升）。
    \end{itemize}
\end{enumerate}
\end{problem}
\begin{solution}
\begin{enumerate}[label=(\alph*)]
    \item 这个经济的均衡是存在状态价格 $\begin{bmatrix}\phi_{a} \\ \phi_{b}\end{bmatrix}$，使得
    \begin{enumerate}[label=（\arabic*）]
        \item 给定状态价格和自己的禀赋，每个参与者都最优化地选择自己的消费；
        \item 市场出清：$\sum c_{k,0}=C_{0}$，$\sum c_{k,a}=C_{a}$，$\sum c_{k,b}=C_{b}$。由于市场是完全的，而参与者的效用函数是严格凹的，因而市场均衡是存在且唯一的。
    \end{enumerate}
    \item 我们知道
    \[
        \begin{aligned}
            \phi_{a}&=\pi_{a}u'(C_{a})/u'(C_{0}),\quad
            \phi_{b}=\pi_{b}u'(C_{b})/u'(C_{0}),\quad
            \frac{\phi_{a}}{\phi_{b}}=\frac{\pi_{a}u'(C_{a})}{\pi_{b}u'(C_{b})}\\
            r_{F}&=\frac{1}{\phi_{a}+\phi_{b}}-1=\frac{u'(C_{0})}{\pi_{a}u'(C_{a})+\pi_{b}u'(C_{b})}-1\\
            u'&>0,\quad u''<0,\quad u'''<0
        \end{aligned}
    \]
    \begin{itemize}
        \item 因为 $u''<0$，故 $C_{0}$ 上升使得 $u'(C_{0})$ 减小；$\phi_{a}$，$\phi_{b}$ 上升，$\phi_{a}/\phi_{b}$ 不变，$r_{F}$ 下降；
        \item 因为 $C_{a}$ 上升使得 $u'(C_{a})$ 下降，因而 $\phi_{a}$，$\phi_{a}/\phi_{b}$ 下降，$\phi_{b}$ 不变，$r_{F}$ 上升；
        \item 类似地，因为 $C_{b}$ 上升使得 $u'(C_{b})$ 下降，因而 $\phi_{b}$ 下降，$\phi_{a}/\phi_{b}$ 上升，$\phi_{a}$ 不变，$r_{F}$ 上升；
        \item $\phi_{a}$，$\phi_{b}$ 下降，$r_{F}$ 上升。现在考虑 $\phi_{a}/\phi_{b}$：
        \[
            \frac{\phi_{a}'/\phi_{b}'}{\phi_{a}/\phi_{b}}=\frac{u'(kC_{a})/u'(kC_{b})}{u'(C_{a})/u'(C_{b})}=\frac{u'(kC_{a})/u'(C_{a})}{u'(kC_{b})/u'(C_{b})}
        \]
        令 $f(x)=\dfrac{u'(kx)}{u'(x)}$，那么 $f'(x)=\dfrac{u'(kx)}{xu'(x)}(-R(kx)+R(x))$，其中 $R(x)$ 是相对风险厌恶系数，可以证明当 $u'''<0$ 时，$R(x)$ 是递增的，因而 $f'(x)<0$。因此，如果 $C_{a}>C_{b}$，那么 $\phi_{a}/\phi_{b}$ 会下降，反之则会上升。
        \item $\phi_{a}$，$\phi_{a}/\phi_{b}$ 下降，$\phi_{b}$ 上升，现在考虑 $r_{F}$。不妨假设 $C_{a}>C_{b}$，现在 $C_{a}$ 上升，$C_{b}$ 下降使得 $\phi_{a}C_{a}+\phi_{b}C_{b}=\tilde{C}$ 保持不变。考虑 $\phi_{a}u'(C_{a})+\phi_{b}u'(C_{b})=\phi_{a}u'(C_{a})+\phi_{b}u'(\tilde{C}-\pi_{a}C_{a})$，对 $C_{a}$ 求导可得 $\pi_{a}[u''(C_{a})-u''(C_{b})]$，因为 $u'''<0$，故 $r_{F}$ 上升。
    \end{itemize}
\end{enumerate}
\end{solution}

\begin{problem}
\problemno{10.2**} 考虑一个含有参与者 $k=1,\dots,K$ 的 Arrow--Debreu 经济。参与者 $k$ 具有形式为 $u_{k,0}(c_{0})+\mathrm{E}[u_{k,1}(c_{1})]$ 的效用函数，对于所有 $k$，$u_{0,k}$ 和 $u_{k,1}$ 是严格递增和严格凹的。$f_{k}^{0}$ 和 $f_{k}$ 是在给定权重集合 $\lambda=\{\lambda_{1},\dots,\lambda_{K}\}$ 下对于参与者 $k$ 的 Pareto 最优分配规则。定义 $u_{0}$ 和 $u_{1}$ 如下
\[
    u_{0}'(x)=\lambda_{k}u_{k,0}(f_{k}^{0}(x)),\qquad u_{1}'(x)=\lambda_{k}u_{k,1}(f_{k}(x)).
\]
\begin{enumerate}[label=(\alph*)]
    \item 证明 $u_{0}$ 和 $u_{1}$ 是严格递增和严格凹的。
    \item 证明 $u_{0}(c_{0})+\mathrm{E}[u_{1}(c_{1})]$ 是代表性参与者的效用函数。
    \item 证明如果所有参与者都是绝对风险厌恶递减的，那么以上定义的代表性参与者的效用函数也是绝对风险厌恶递减的。
\end{enumerate}
\end{problem}
\begin{solution}
\begin{enumerate}[label=(\alph*)]
    \item $u_{1}'(x)=\lambda_{k}u_{k,1}'(f_{k}(x))>0$，$u_{1}''(x)=\lambda_{k}u_{k,1}''(f_{k}(x))f_{k}'(x)<0$，这是因为 $u_{k}''<0$ 而 $f_{k}'>0$（定理 10.6），因而 $u_{1}$ 是严格递增且凹的，同样可以证明 $u_{0}$ 是严格递增且凹的；
    \item 由 (10.14) 式我们知道，代表性参与者的效用函数满足
    \[
        u_{0}'(x)=\lambda_{k}u_{k,0}(f_{k}^{0}(x)),\qquad u_{1}'(x)=\lambda_{k}u_{k,1}(f_{k}(x)),\quad\forall k
    \]
    因此（除相差一个常数为）$u_{0}(c_{0})+\mathrm{E}[u_{1}(c_{1})]$ 就是代表性参与者的效用函数。
    \item 由于 $u_{0}$、$u_{1}$ 的证明方法完全一样，因而我们忽略下标 $0$、$1$。如果参与者是绝对风险厌恶递减的，那么有 $[u_{k}'']^{2}-u_{k}'u_{k}'''<0$。对于代表性参与者有
    \[
        A(x)=-\frac{u_{k}''(f_{k}(x))f_{k}'(x)}{u_{k}'(f_{k}(x))}
    \]
    那么
    \[
        A'(x)=\frac{[u_{k}'']^{2}[f_{k}']^{2}-u_{k}'u_{k}''f_{k}''-u_{k}'u_{k}'''[f_{k}']^{2}}{[u_{k}']^{2}}
    \]
    因此，
    \[
        f_{k}''=\frac{A'(x)[u_{k}']^{2}+u_{k}'u_{k}'''[f_{k}']^{2}-[u_{k}'']^{2}[f_{k}']^{2}}{-u_{k}'u_{k}''}
    \]
    对于分配规则，我们有
    \[
        \sum f_{k}(x)=x,\quad \sum f_{k}'(x)=1,\quad \sum f_{k}''(x)=0
    \]
    因此我们有
    \[
        A'(x)=\sum\frac{[f_{k}']^{2}\bigl([u_{k}'']^{2}-u_{k}'u_{k}'''\bigr)}{-u_{k}'u_{k}''}\Big/\sum\frac{-u_{k}'}{u_{k}''}<0
    \]
    因而代表性参与者的效用函数也是绝对厌恶递减的。
\end{enumerate}
\end{solution}

\begin{problem}
\problemno{10.3*} 在 1 期，经济有两个可能状态 $a$ 和 $b$。它们的发生的概率分别为 $\pi$ 和 $1-\pi$。市场有两只交易证券，证券 1 和 2，支付如下：
\begin{center}
\begin{tabular}{@{}c@{\hspace{2.5em}}c@{}}
证券 1：
\begin{tikzpicture}[x=1.15em,y=1.15em,baseline=(current bounding box.center)]
    \draw[thick] (-0.15,-1) -- (-0.15,1);
    \draw[thick] (-0.15,0.55) -- (0.85,0.55);
    \draw[thick] (-0.15,-0.55) -- (0.85,-0.55);
    \node[anchor=west] at (1.0,0.55) {$1$};
    \node[anchor=west] at (1.0,-0.55) {$1$};
\end{tikzpicture}
&
证券 2：
\begin{tikzpicture}[x=1.15em,y=1.15em,baseline=(current bounding box.center)]
    \draw[thick] (-0.15,-1) -- (-0.15,1);
    \draw[thick] (-0.15,0.55) -- (0.85,0.55);
    \draw[thick] (-0.15,-0.55) -- (0.85,-0.55);
    \node[anchor=west] at (1.0,0.55) {$uC$};
    \node[anchor=west] at (1.0,-0.55) {$dC$};
\end{tikzpicture}
\end{tabular}
\end{center}
其中 $C>0$ 且 $0<d<u$。因此，证券 1 是无风险债券而证券 2 是风险证券，我们可以称它为股票。经济中有参与者 $k=1,\dots,K$。参与者 $k$ 的禀赋是 $s_{k}C$ 单位的当期消费、0 单位的证券 1 和 $s_{k}$ 单位的证券 2，其中 $s_{k}>0$ 且 $\sum_{k}s_{k}=1$。所有参与者具有如下的期望效用函数：
\[
    U_{k}(c_{k})=\log c_{k,0}+\rho\bigl[\pi\log c_{k,1a}+(1-\pi)\log c_{k,1b}\bigr]
\]
其中 $\rho>0$。
\begin{enumerate}[label=(\alph*)]
    \item 证明证券市场是完全的。
    \item 求解每一参与者的最优消费/投资选择。
    \item 求解均衡证券价格。
    \item 将价格过程与期权定价中的二叉树模型相比较，你能得到什么结论？
\end{enumerate}
\end{problem}
\begin{solution}
\begin{enumerate}[label=(\alph*)]
    \item 市场上有两只支付独立的证券，而状态数目恰为 2，因而市场是完全的。
    \item 参与者 $k$ 面临的问题是
    \[
        \begin{aligned}
            &\max_{\{c_{k,0},c_{k,1a},c_{k,1b}\}}\ \log c_{k,0}+\rho\bigl[\pi\log c_{k,1a}+(1-\pi)\log c_{k,1b}\bigr]\\
            &\quad\text{s.t.}\quad c_{k,0}+\phi_{a}c_{k,1a}+\phi_{b}c_{k,1b}=S_{k}C(1+\phi_{a}u+\phi_{b}d)
        \end{aligned}
    \]
    可以求得
    \[
        \begin{aligned}
            c_{k,0}&=\frac{S_{k}C(1+\phi_{a}u+\phi_{b}d)}{1+\rho}\\
            c_{k,1a}&=\frac{\rho\pi S_{k}C(1+\phi_{a}u+\phi_{b}d)}{\phi_{a}(1+\rho)}\\
            c_{k,1b}&=\frac{\rho(1-\pi)S_{k}C(1+\phi_{a}u+\phi_{b}d)}{\phi_{b}(1+\rho)}
        \end{aligned}
    \]
    当市场达到均衡时：
    \[
        \begin{aligned}
            \sum c_{k,0}&=\sum S_{k}C=C\\
            \sum c_{k,1a}&=\sum S_{k}uC=uC\\
            \sum c_{k,1b}&=\sum S_{k}dC=dC
        \end{aligned}
    \]
    因此求得
    \[
        \begin{aligned}
            \phi_{a}&=\frac{\rho\pi}{u},\qquad\phi_{b}=\frac{\rho(1-\pi)}{d}\\
            c_{k,0}&=S_{k}C,\qquad c_{k,1a}=S_{k}uC,\qquad c_{k,1b}=S_{k}dC
        \end{aligned}
    \]
    \item 债券价格为 $\dfrac{\rho\pi}{u}+\dfrac{\rho(1-\pi)}{d}$，股票价格为 $\rho C$。
    \item 在二叉树模型中，债券、股票的价格都是给定的，因而状态价格只由债券、股票的支付给定，得到的是部分均衡的结果；在这个模型中，价格是通过市场均衡求得的，因而状态价格还与参与者的效用函数、时间偏好系数、状态发生的概率有关，得到的是一般均衡的结果。
\end{enumerate}
\end{solution}

\begin{problem}
\problemno{10.4**} 继续考虑练习 10.3 定义的经济，
\begin{enumerate}[label=(\alph*)]
    \item 构建代表性参与者。
    \item 计算在代表性参与者（即中央计划者）的效用函数中每个参与者的权重。它依赖于什么？为什么？
    \item 证明代表性参与者的偏好与禀赋在参与者之间的分布无关。
    \item 用基于消费的 CAPM 为两只证券定价。
    \item 计算无风险收益和股票的风险溢价。
    \item 利率和股票的风险溢价依赖于什么？解释它们对经济原生变量：当前消费水平、消费增长率的期望值 $\mathrm{E}[\tilde{C}_{1}/C_{0}]$、消费增长率的波动率 $\mathrm{Var}[\tilde{C}_{1}/C_{0}]$ 以及时间折现系数 $\rho$ 的依赖性？
\end{enumerate}
\end{problem}
\begin{solution}
\begin{enumerate}[label=(\alph*)]
    \item 由第二题我们知道，代表性参与者的效用函数满足
    \[
        u_{0}'(C)=\mu_{k}u_{k,0}'(c_{k})=\frac{\mu_{k}}{c_{k}}=\frac{\mu_{k}}{S_{k}C},\quad\forall k
    \]
    因此 $\mu_{k}=S_{k}$，$u_{0}'(C)=\dfrac{1}{C}$，$u_{0}(C)=\log C$，同样地，$u_{1}(C)=\log C$；因而代表性参与者的效用函数为
    \[
        \log C_{0}+\rho\bigl[\pi\log C_{1a}+(1-\pi)\log C_{1b}\bigr]
    \]
    \item 参与者 $k$ 的权重是 $S_{k}$，即他的禀赋占总禀赋的比例。如果参与者拥有的资源愈多，在均衡条件下，他的消费越多，权重也越大；
    \item 代表性参与者的偏好与 $S_{k}$ 无关，即与禀赋在参与者之间的分布无关；
    \item 由 (10.16) 式我们有：
    \[
        \begin{aligned}
            \text{债券价格为：}&\quad
            \pi\frac{\rho C}{uC}+(1-\pi)\frac{\rho C}{dC}=\rho\Bigl[\frac{\pi}{u}+\frac{1-\pi}{d}\Bigr]\\
            \text{股票价格为：}&\quad
            \pi\frac{\rho C}{uC}uC+(1-\pi)\frac{\rho C}{dC}dC=\rho C
        \end{aligned}
    \]
    \item 无风险收益率为
    \[
        r_{F}=1\big/\Bigl(\rho\Bigl[\frac{\pi}{u}+\frac{1-\pi}{d}\Bigr]\Bigr)-1
    \]
    风险溢价为
    \[
        p=\Bigl(\frac{\pi}{u}+\frac{1-\pi}{d}\Bigr)\big/\rho-1\big/\Bigl(\rho\Bigl[\frac{\pi}{u}+\frac{1-\pi}{d}\Bigr]\Bigr)
    \]
    \item 如果 $u$，$d$ 与 1 的距离很小，我们可以把风险溢价写成 $p=\dfrac{\mathrm{E}[\tilde{C}_{1}]/C_{0}}{\rho}-r_{F}$，其中 $r_{F}=1\big/\bigl\{\rho\bigl[\mathrm{Var}[\tilde{C}_{1}/C_{0}]-2\mathrm{E}[\tilde{C}_{1}/C_{0}]+3\bigr]\bigr\}$。无风险利率依赖于时间偏好系数 $\rho$、消费增长率的波动率 $\mathrm{Var}[\tilde{C}_{1}/C_{0}]$、以及消费增长率的期望值 $\mathrm{E}[\tilde{C}_{1}/C_{0}]$。当 $\rho$ 减小时，参与者偏好于当前消费，而为了吸引投资者储蓄以平衡当前消费和未来消费，利率 $r_{F}$ 应该上升；当未来消费的风险上升时，为规避未来消费的风险，参与者对债券的需求上升，其价格上升、利率下降；当 $\mathrm{E}[\tilde{C}_{1}/C_{0}]$ 上升时，人们倾向于在时间上平滑消费水平，即增加当前消费，为吸引投资，利率上升。至于风险溢价，当 $\rho$ 上升时，参与者倾向于未来消费，因而债券、股票的需求量都会上升，利率和股票收益率都减小，而且对两者的影响相对一样，因而风险溢价变小（风险溢价一般都是正的）。当 $\mathrm{E}[\tilde{C}_{1}/C_{0}]$ 变大时，参与者倾向于当前消费，因而对债券和股票的需求量都减小，两者的价格都下降、收益率都上升，对风险溢价的影响是不确定的，但当 $\mathrm{E}[\tilde{C}_{1}/C_{0}]$ 较小时，利率上升更快，风险溢价会变小。当 $\mathrm{Var}[\tilde{C}_{1}/C_{0}]$ 变大时，因为股票的风险和禀赋的风险完全一致，为了减少未来消费的波动，人们对债券的需求增加而对股票的需求减少，因而债券收益率下降而股票收益上升，风险溢价上升。
\end{enumerate}
\end{solution}

\begin{problem}
\problemno{10.5} 继续前面练习中的经济。
\begin{enumerate}[label=(\alph*)]
    \item 求解以证券 2 为标的资产、执行价格为 $K$ 的看涨期权的价格。
    \item 证明基于消费的 CAPM 也适用于期权。
\end{enumerate}
\end{problem}
\begin{solution}
\begin{enumerate}[label=(\alph*)]
    \item 不失一般性，假设 $dC\le K\le uC$，期权的支付为 $\begin{bmatrix}uC-K \\ 0\end{bmatrix}$，用复制的方法可以求得期权的价格为 $\dfrac{\rho\pi}{u}(uC-K)$；
    \item 而用 CCAPM 定价，期权的价格为
    \[
        \pi\frac{\rho C}{uC}[uC-K]_{+}+(1-\pi)\frac{\rho C}{dC}[dC-K]_{+}=\frac{\rho\pi}{u}(uC-K)
    \]
    因而 CCAPM 成立。
\end{enumerate}
\end{solution}

\begin{problem}
\problemno{10.6**} 在 1 期，经济有两个可能状态 $a$ 和 $b$。它们的发生概率分别为 $\pi$ 和 $1-\pi$。市场有两只交易证券，证券 1 和 2，支付如下：
\begin{center}
\begin{tabular}{@{}c@{\hspace{2.5em}}c@{}}
证券 1：
\begin{tikzpicture}[x=1.15em,y=1.15em,baseline=(current bounding box.center)]
    \draw[thick] (-0.15,-1) -- (-0.15,1);
    \draw[thick] (-0.15,0.55) -- (0.85,0.55);
    \draw[thick] (-0.15,-0.55) -- (0.85,-0.55);
    \node[anchor=west] at (1.0,0.55) {$1$};
    \node[anchor=west] at (1.0,-0.55) {$1$};
\end{tikzpicture}
&
证券 2：
\begin{tikzpicture}[x=1.15em,y=1.15em,baseline=(current bounding box.center)]
    \draw[thick] (-0.15,-1) -- (-0.15,1);
    \draw[thick] (-0.15,0.55) -- (0.85,0.55);
    \draw[thick] (-0.15,-0.55) -- (0.85,-0.55);
    \node[anchor=west] at (1.0,0.55) {$C_{1a}$};
    \node[anchor=west] at (1.0,-0.55) {$C_{1b}$};
\end{tikzpicture}
\end{tabular}
\end{center}
其中 $C>0$ 且 $0<d<u$。因此，证券 1 是无风险债券而证券 2 是风险证券，我们可以称它为股票。经济中有参与者 $k=1,\dots,K$。参与者 $k$ 的禀赋是 $s_{k}C_{0}$ 单位的当期消费、0 单位的证券 1 和 $s_{k}$ 单位的证券 2，其中 $s_{k}>0$ 且 $\sum_{k}s_{k}=1$。所有参与者具有如下的期望效用函数：
\[
    U_{k}(c_{k})=-e^{-\alpha c_{k,0}}-\rho\bigl[\pi e^{-\alpha c_{k,1a}}+(1-\pi)e^{-\alpha c_{k,1b}}\bigr]
\]
其中 $\rho>0$。
\begin{enumerate}[label=(\alph*)]
    \item 证明证券市场是完全的。
    \item 求解每一参与者的最优消费/投资选择。
    \item 求解均衡证券价格。
\end{enumerate}
\end{problem}
\begin{solution}
\begin{enumerate}[label=(\alph*)]
    \item 有 2 个可能状态，而支付独立的状态数目恰为 2，因此证券市场是完全的；
    \item 假设状态价格为 $\begin{bmatrix}\phi_{a} \\ \phi_{b}\end{bmatrix}$；考虑参与者 $k$，他的优化问题是
    \[
        \begin{aligned}
            &\max_{c_{k,0},\,c_{k,1a},\,c_{k,1b}}\ -e^{-\alpha c_{k,0}}-\rho\bigl[\pi e^{-\alpha c_{k,1a}}+(1-\pi)e^{-\alpha c_{k,1b}}\bigr]\\
            &\quad\text{s.t.}\quad c_{k,0}+\phi_{a}c_{k,1a}+\phi_{b}c_{k,1b}=s_{k}C_{0}+\phi_{a}s_{k}C_{1a}+\phi_{b}s_{k}C_{1b}
        \end{aligned}
    \]
    令 $W=C_{0}+\phi_{a}C_{1a}+\phi_{b}C_{1b}$ 为经济中的总财富。求解优化问题可得参与者的消费选择为
    \[
        \begin{aligned}
            c_{k,0}&=\frac{1}{\alpha}\log\alpha-\frac{\beta-\alpha s_{k}W}{\alpha(1+\phi_{a}+\phi_{b})}\\
            c_{k,1a}&=\frac{1}{\alpha}\log\frac{\rho\pi\alpha}{\phi_{a}}-\frac{\beta-\alpha s_{k}W}{\alpha(1+\phi_{a}+\phi_{b})}\\
            c_{k,1b}&=\frac{1}{\alpha}\log\frac{\rho(1-\pi)\alpha}{\phi_{b}}-\frac{\beta-\alpha s_{k}W}{\alpha(1+\phi_{a}+\phi_{b})}
        \end{aligned}
    \]
    其中，$\beta=\log\alpha+\phi_{a}\log\dfrac{\rho\pi\alpha}{\phi_{a}}+\phi_{b}\log\dfrac{\rho(1-\pi)\alpha}{\phi_{b}}$；
    \item 市场出清条件是
    \[
        \begin{aligned}
            \sum c_{k,0}&=C_{0}\\
            \sum c_{k,1a}&=C_{1a}\\
            \sum c_{k,1b}&=C_{1b}
        \end{aligned}
    \]
    由此可以求得均衡状态价格为
    \[
        \begin{aligned}
            \phi_{a}&=\rho\pi e^{-\frac{\alpha}{K}(C_{1a}-C_{0})}\\
            \phi_{b}&=\rho(1-\pi)e^{-\frac{\alpha}{K}(C_{1b}-C_{0})}
        \end{aligned}
    \]
    因而债券的均衡价格为 $\phi_{a}+\phi_{b}$，股票的均衡价格为 $\phi_{a}C_{1a}+\phi_{b}C_{1b}$；在均衡条件下，参与者 $k$ 的最优消费选择为
    \[
        \begin{aligned}
            c_{k,0}^{*}&=\frac{1}{K}C_{0}+\Bigl(s_{k}-\frac{1}{K}\Bigr)\frac{W}{1+\phi_{a}+\phi_{b}}\\
            c_{k,1a}^{*}&=\frac{1}{K}C_{1a}+\Bigl(s_{k}-\frac{1}{K}\Bigr)\frac{W}{1+\phi_{a}+\phi_{b}}\\
            c_{k,1b}^{*}&=\frac{1}{K}C_{1b}+\Bigl(s_{k}-\frac{1}{K}\Bigr)\frac{W}{1+\phi_{a}+\phi_{b}}
        \end{aligned}
    \]
    在 0 期他持有的无风险债券、股票的头寸分别为
    \[
        \begin{aligned}
            \theta_{B}&=\Bigl(s_{k}-\frac{1}{K}\Bigr)\frac{W}{1+\phi_{a}+\phi_{b}}\\
            \theta_{S}&=\frac{1}{K}-s_{k}.
        \end{aligned}
    \]
\end{enumerate}
\end{solution}

\begin{problem}
\problemno{10.7**} 继续考虑练习 10.6 定义的经济
\begin{enumerate}[label=(\alph*)]
    \item 构建代表性参与者。
    \item 计算在代表性参与者（即中央计划者）的效用函数中每个参与者的权重。它依赖于什么？为什么？
    \item 证明代表性参与者的偏好与禀赋在参与者之间的分布无关。
    \item 用基于消费的 CAPM 为两只证券定价。
    \item 计算无风险收益和股票的风险溢价，现在风险溢价的定义为 $\mathrm{E}[\tilde{S}_{1}]-(1+r_{F})S_{0}$。即一股股票的绝对超额收益。
    \item 利率和股票的风险溢价依赖于什么？解释它们对经济原生变量：当前消费水平、消费增长的期望值 $\mathrm{E}[\tilde{C}_{1}/C_{0}]$、消费增长的波动率 $\mathrm{Var}[\tilde{C}_{1}/C_{0}]$ 以及时间折现系数 $\rho$ 的依赖性？
    \item 求解以证券 2 为标的资产、执行价格为 $K$ 的看涨期权的价格。
\end{enumerate}
\end{problem}
\begin{solution}
\begin{enumerate}[label=(\alph*)]
    \item 因为每个参与者的 0、1 期效用函数只差时间偏好系数 $\rho$，因而代表性参与者的 0、1 期效用函数只差时间偏好系数 $\rho$；假设中央计划者问题中参与者 $k$ 的权重为 $\mu_{k}>0$，代表性参与者的 0 期效用函数为 $u(C_{0})$，那么
    \[
        \begin{aligned}
            u'(C_{0})
            &=\mu_{k}u_{k}'(c_{k,0})\\
            &=\mu_{k}\alpha e^{-\alpha\bigl(\frac{1}{K}C_{0}+\bigl(s_{k}-\frac{1}{K}\bigr)\frac{W}{1+\phi_{a}+\phi_{b}}\bigr)}\\
            &=e^{-\alpha\frac{C_{0}}{K}}\Bigl(\alpha\mu_{k}e^{-\alpha\bigl(s_{k}-\frac{1}{K}\bigr)\frac{W}{1+\phi_{a}+\phi_{b}}}\Bigr),\quad\forall k
        \end{aligned}
    \]
    因此，$\alpha\mu_{k}e^{-\alpha\bigl(s_{k}-\frac{1}{K}\bigr)\frac{W}{1+\phi_{a}+\phi_{b}}}$ 是与 $k$ 无关的常数，即 $\mu_{k}$ 与 $e^{\frac{\alpha}{1+\phi_{a}+\phi_{b}}s_{k}}$ 成正比，而且代表性参与者的效用函数具有负指数形式 $u(C_{0})=-e^{-\alpha\frac{C_{0}}{K}}$，因此，代表性参与者的效用函数为
    \[
        -e^{-\alpha\frac{C_{0}}{K}}-\rho\Bigl[\pi e^{-\alpha\frac{C_{1a}}{K}}+(1-\pi)e^{-\alpha\frac{C_{1b}}{K}}\Bigr]
    \]
    \item 权重 $\mu_{k}$ 与 $e^{\frac{\alpha}{1+\phi_{a}+\phi_{b}}s_{k}}$ 成正比，依赖于参与者的禀赋参数 $s_{k}$，即他的禀赋占总禀赋的比例。如果参与者拥有的资源愈多，在均衡条件下，他的消费越多，权重也越大；
    \item 代表性参与者的偏好与 $s_{k}$ 无关，即与禀赋在参与者之间的分布无关；
    \item
    \[
        \begin{aligned}
            \text{债券价格为：}&\quad
            \rho\pi\frac{\alpha e^{-\alpha\frac{C_{1a}}{K}}}{\alpha e^{-\alpha\frac{C_{0}}{K}}}+\rho(1-\pi)\frac{\alpha e^{-\alpha\frac{C_{1b}}{K}}}{\alpha e^{-\alpha\frac{C_{0}}{K}}}\\
            &\qquad=\rho\Bigl[\pi e^{-\alpha\frac{C_{1a}-C_{0}}{K}}+(1-\pi)e^{-\alpha\frac{C_{1b}-C_{0}}{K}}\Bigr]\\
            \text{股票价格为：}&\quad
            \rho\pi\frac{\alpha e^{-\alpha\frac{C_{1a}}{K}}}{\alpha e^{-\alpha\frac{C_{0}}{K}}}C_{1a}+\rho(1-\pi)\frac{\alpha e^{-\alpha\frac{C_{1b}}{K}}}{\alpha e^{-\alpha\frac{C_{0}}{K}}}C_{1b}\\
            &\qquad=\rho\Bigl[\pi e^{-\alpha\frac{C_{1a}-C_{0}}{K}}C_{1a}+(1-\pi)e^{-\alpha\frac{C_{1b}-C_{0}}{K}}C_{1b}\Bigr]
        \end{aligned}
    \]
    \item 无风险收益率为
    \[
        r_{F}=\frac{1}{\phi_{a}+\phi_{b}}-1=\frac{1}{\rho\Bigl[\pi e^{-\alpha\frac{C_{1a}-C_{0}}{K}}+(1-\pi)e^{-\alpha\frac{C_{1b}-C_{0}}{K}}\Bigr]}-1
    \]
    风险溢价为
    \[
        \begin{aligned}
            p&=\mathrm{E}[\tilde{S}_{1}]-(1+r_{F})S_{0}\\
            &=\pi C_{1a}+(1-\pi)C_{1b}-\frac{\phi_{a}C_{1a}+\phi_{b}C_{1b}}{\phi_{a}+\phi_{b}}\\
            &=\Bigl(\pi-\frac{\phi_{a}}{\phi_{a}+\phi_{b}}\Bigr)(C_{1a}-C_{1b})\\
            &=\frac{\pi(1-\pi)\bigl(e^{\alpha\frac{C_{1a}-C_{1b}}{K}}-1\bigr)}{\pi+(1-\pi)e^{\alpha\frac{C_{1a}-C_{1b}}{K}}}(C_{1a}-C_{1b})
        \end{aligned}
    \]
    \item 如果 $\tilde{g}=\tilde{C}_{1}/C_{0}-1$ 很小，而人均消费 $C_{0}/K$ 是很大，那么无风险利率可以写成
    \[
        r_{F}\doteq\frac{1}{\rho\Bigl[1-\alpha\frac{C_{0}}{K}\mathrm{E}[\tilde{g}]+\alpha^{2}\frac{C_{0}^{2}}{K^{2}}\bigl(\mathrm{Var}[\tilde{g}]+\mathrm{E}^{2}[\tilde{g}]\bigr)\Bigr]}
    \]
    可以看出，无风险利率依赖于时间偏好系数 $\rho$、消费增长率的波动率 $\mathrm{Var}[\tilde{C}_{1}/C_{0}]$、以及消费增长率的期望值 $\mathrm{E}[\tilde{C}_{1}/C_{0}]$。当 $\rho$ 减小时，参与者偏好于当前消费，而为了吸引投资者储蓄以平衡当前消费和未来消费，利率 $r_{F}$ 应该上升；当未来消费的风险上升时，为规避未来消费的风险，参与者对债券的需求上升，其价格上升、利率下降；当 $\mathrm{E}[\tilde{C}_{1}/C_{0}]$ 上升时，人们倾向于在时间上平滑消费水平，即增加当前消费，为吸引投资，利率上升。至于风险溢价，它只与 $C_{1a}-C_{1b}$ 有关，且是 $C_{1a}-C_{1b}$ 的增函数；某种意义上我们可以把 $C_{1a}-C_{1b}$ 作为未来禀赋风险的一种度量，也就是说，当未来禀赋的风险增加时，风险溢价增加，这是因为股票的风险和禀赋的风险完全一致，为了减少未来消费的波动，人们对债券的需求增加而对股票的需求减少，因而债券收益率下降而股票收益上升，风险溢价上升。
    \item 可以用 CCAPM 为期权定价
    \[
        \begin{aligned}
            c&=\rho\pi\frac{\alpha e^{-\alpha\frac{C_{1a}}{K}}}{\alpha e^{-\alpha\frac{C_{0}}{K}}}[C_{1a}-K]_{+}+\rho(1-\pi)\frac{\alpha e^{-\alpha\frac{C_{1b}}{K}}}{\alpha e^{-\alpha\frac{C_{0}}{K}}}[C_{1b}-K]_{+}\\
            &=\rho\Bigl[\pi e^{-\alpha\frac{C_{1a}-C_{0}}{K}}[C_{1a}-K]_{+}+(1-\pi)e^{-\alpha\frac{C_{1b}-C_{0}}{K}}[C_{1b}-K]_{+}\Bigr]
        \end{aligned}
    \]
\end{enumerate}
\end{solution}

% !TeX root = ../../金融经济学习题讲义.tex
% 源 PDF：references/金融经济学/史上最全最完整的金融经济学答案(王江版).pdf，页 54–61（印刷页 53–60）

\zhchapter{不完全市场中的资源配置和资产定价}

\ifdefined\WJHandoutKnowledge
\fi

% !TeX root = ../../金融经济学习题讲义.tex
% 11_不完全市场 课后习题（由章节文件拆分）

\begin{problem}
\problemno{11.1*} 经济有两期 $t=0,1$。有两个参与者，1 和 2，他们都有定义于消费路径上的对数效用函数 $c=(c_{0},c_{1})$：$u(c)=\log c_{0}+\log c_{1}$。在 $t=1$ 期有三个等概率的可能状态，$\omega_{1}$，$\omega_{2}$ 以及 $\omega_{3}$。在 $t=0$ 期的总禀赋是 1 而 $t=1$ 期的总禀赋是 $\begin{bmatrix}1 \\ 1 \\ 2\end{bmatrix}$。有两只交易证券，一只是无风险债券，在 $t=1$ 期的所有状态下都支付 1 单位的消费品。另外一只是股票，它的 $t=1$ 期支付和总禀赋完全一样，即它的支付向量是 $\begin{bmatrix}1 \\ 1 \\ 2\end{bmatrix}$。
\begin{enumerate}[label=(\alph*)]
    \item 假设参与者 1 和 2 的禀赋分别为 $\begin{bmatrix}0.4\\0.4\\0.4\\0.8\end{bmatrix}$ 和 $\begin{bmatrix}0.6\\0.6\\0.6\\1.2\end{bmatrix}$。求解均衡利率和均衡股价。均衡配置是 Pareto 最优的吗？解释原因。
    \item 假设参与者 $a$ 和 $b$ 的禀赋分别为 $\begin{bmatrix}0.5\\1\\0\\1\end{bmatrix}$ 和 $\begin{bmatrix}0.5\\0\\1\\1\end{bmatrix}$。导出求解均衡的一阶条件和市场出清条件。在不解出均衡的情况下，讨论均衡配置是否是 Pareto 最优的。解释原因。
\end{enumerate}
\end{problem}
\begin{solution}
\begin{enumerate}[label=(\alph*)]
    \item 假设债券股票的价格分别为 $B$，$S$。对于参与者 1 来说，假设他持有的债券和股票头寸分别为 $\theta_{11}$，$\theta_{12}$，那么他的优化问题是
    \[
        \begin{aligned}
            &\max_{\theta_{11},\,\theta_{12}}\ \log(c_{0})+\tfrac13\bigl[\log(c_{1,1})+\log(c_{1,2})+\log(c_{1,3})\bigr]\\
            &\quad\text{s.t.}\quad c_{0}=0.4-\theta_{11}B-\theta_{12}S\\
            &\qquad c_{1,1}=0.4+\theta_{11}+\theta_{12}\\
            &\qquad c_{1,2}=0.4+\theta_{11}+\theta_{12}\\
            &\qquad c_{1,3}=0.8+\theta_{11}+2\theta_{12}
        \end{aligned}
    \]
    可以求得
    \[
        \begin{aligned}
            \theta_{11}&=-\frac{1}{15}(5S^{2}+5S-6BS-6B)/(S^{2}-3BS+2B^{2})\\
            \theta_{12}&=-\frac{1}{15}(3S^{2}-3S-14BS+12B^{2}+4B)/(S^{2}-3BS+2B^{2})
        \end{aligned}
    \]
    同样地可以求得参与者 2 对债券、股票的需求量为
    \[
        \begin{aligned}
            \theta_{21}&=-\frac{1}{10}(5S^{2}+5S-6BS-6B)/(S^{2}-3BS+2B^{2})\\
            \theta_{22}&=-\frac{1}{10}(3S^{2}-3S-14BS+12B^{2}+4B)/(S^{2}-3BS+2B^{2})
        \end{aligned}
    \]
    市场出清条件为
    \[
        \begin{aligned}
            \theta_{11}+\theta_{21}&=0\\
            \theta_{12}+\theta_{22}&=0
        \end{aligned}
    \]
    因而可以求得债券、股票的价格分别为
    \[
        B=\frac{5}{6},\quad S=1
    \]
    均衡利率为 $r_{F}=1/B-1=20\%$；可以验证 $\dfrac{u_{1,1}'(\tilde{c}_{1,1})}{u_{1,0}'(\tilde{c}_{1,0})}=\dfrac{u_{2,1}'(\tilde{c}_{2,1})}{u_{2,0}'(\tilde{c}_{2,0})}$，即边际效用之比在各状态下均相等，因此均衡配置是 Pareto 最优的。
    \item 对于参与者 1 来说，最优化的一阶条件是
    \[
        \begin{aligned}
            \frac{-B}{0.5-\theta_{11}B-\theta_{12}S}+\frac13\Bigl[\frac{1}{1+\theta_{11}+\theta_{12}}+\frac{1}{\theta_{11}+\theta_{12}}+\frac{1}{1+\theta_{11}+2\theta_{12}}\Bigr]&=0\\
            \frac{-S}{0.5-\theta_{11}B-\theta_{12}S}+\frac13\Bigl[\frac{1}{1+\theta_{11}+\theta_{12}}+\frac{1}{\theta_{11}+\theta_{12}}+\frac{2}{1+\theta_{11}+2\theta_{12}}\Bigr]&=0
        \end{aligned}
    \]
    对于参与者 2 来说，最优化的一阶条件是
    \[
        \begin{aligned}
            \frac{-B}{0.5-\theta_{21}B-\theta_{22}S}+\frac13\Bigl[\frac{1}{\theta_{21}+\theta_{22}}+\frac{1}{1+\theta_{21}+\theta_{22}}+\frac{1}{1+\theta_{21}+2\theta_{22}}\Bigr]&=0\\
            \frac{-S}{0.5-\theta_{21}B-\theta_{22}S}+\frac13\Bigl[\frac{1}{\theta_{21}+\theta_{22}}+\frac{1}{1+\theta_{21}+\theta_{22}}+\frac{2}{1+\theta_{21}+2\theta_{22}}\Bigr]&=0
        \end{aligned}
    \]
    市场出清条件为
    \[
        \begin{aligned}
            \theta_{11}+\theta_{21}&=0\\
            \theta_{12}+\theta_{22}&=0
        \end{aligned}
    \]
    配置是 Pareto 最优的一个条件是边际效用之比在各状态之间相等，考虑
    \[
        \begin{aligned}
            \frac{\theta_{11}+\theta_{12}}{1+\theta_{11}+\theta_{12}}\Big/\frac{1+\theta_{21}+\theta_{22}}{\theta_{21}+\theta_{22}}
            &=\frac{\theta_{11}+\theta_{12}}{1+\theta_{11}+\theta_{12}}\Big/\frac{1-\theta_{11}-\theta_{12}}{-\theta_{11}-\theta_{12}}\\
            &=1\big/\bigl(1-\tfrac{1}{(\theta_{11}+\theta_{12})^{2}}\bigr)\ne 1
        \end{aligned}
    \]
    因而，均衡配置不是 Pareto 最优的。
\end{enumerate}
\end{solution}

\begin{problem}
\problemno{11.2*} 经济中有 $K$ 个具有 CARA 偏好的参与者。经济中只有两个企业，它们的资产在 1 期的支付具有联合正态分布。相关的参数如下：
\begin{center}
\begin{tabular}{lccc}
    \toprule
     & 均值 & 标准差 & 相关系数 \\
    \midrule
    企业 1 & 10 & 8 & \multirow{2}{*}{0.6} \\
    企业 2 & 12 & 11 & \\
    \bottomrule
\end{tabular}
\end{center}
利率为 $8\%$。推导两个企业 0 期市场价值之间的关系。
\end{problem}
\begin{solution}
考虑参与者 $k$，假设他的效用函数为 $-e^{-a_{k}\tilde{w}}$，那么他的优化问题是
\[
    \begin{aligned}
        &\max_{\theta_{k,1},\,\theta_{k,2}}\ \mathrm{E}\bigl[-e^{-a_{k}\tilde{w}_{k}}\bigr]\\
        &\quad\text{s.t.}\quad \tilde{w}_{k}=(w_{k}-P_{1}\theta_{k,1}-P_{2}\theta_{k,2})(1+r_{F})+\tilde{r}_{1}\theta_{k,1}+\tilde{r}_{2}\theta_{k,2}
    \end{aligned}
\]
其中 $P_{1}$，$P_{2}$ 分别是企业 1、2 的价格。支付联合正态则 $\tilde{w}_{k}$ 正态，与练习 9.7 相同：最大化 $E[\tilde{w}_{k}]-\dfrac{a_{k}}{2}\operatorname{Var}(\tilde{w}_{k})$。记超额支付 $\mu_{i}=\bar{r}_{i}-P_{i}(1+r_{F})$，$\mu=\begin{bmatrix}\mu_{1}\\\mu_{2}\end{bmatrix}$，协方差阵 $\Sigma=\begin{bmatrix}\sigma_{1}^{2}&\rho\sigma_{1}\sigma_{2}\\\rho\sigma_{1}\sigma_{2}&\sigma_{2}^{2}\end{bmatrix}$。一阶条件 $\mu=a_{k}\Sigma\theta_{k}$ 给出 $\theta_{k}=\dfrac{1}{a_{k}}\Sigma^{-1}\mu$。展开
\[
    \Sigma^{-1}=\frac{1}{1-\rho^{2}}\begin{bmatrix}1/\sigma_{1}^{2}&-\rho/(\sigma_{1}\sigma_{2})\\-\rho/(\sigma_{1}\sigma_{2})&1/\sigma_{2}^{2}\end{bmatrix}
\]
即
\[
    \begin{aligned}
        \theta_{k,1}&=\frac{1}{a_{k}(1-\rho^{2})}\Bigl[\frac{\bar{r}_{1}-P_{1}(1+r_{F})}{\sigma_{1}^{2}}-\frac{\rho}{\sigma_{1}\sigma_{2}}\bigl(\bar{r}_{2}-P_{2}(1+r_{F})\bigr)\Bigr]\\
        \theta_{k,2}&=\frac{1}{a_{k}(1-\rho^{2})}\Bigl[\frac{\bar{r}_{2}-P_{2}(1+r_{F})}{\sigma_{2}^{2}}-\frac{\rho}{\sigma_{1}\sigma_{2}}\bigl(\bar{r}_{1}-P_{1}(1+r_{F})\bigr)\Bigr]
    \end{aligned}
\]
在均衡条件下，企业资产的总需求都是 1，因而我们有
\[
    \frac{\bar{r}_{1}-P_{1}(1+r_{F})}{\sigma_{1}^{2}}-\frac{\rho}{\sigma_{1}\sigma_{2}}\bigl(\bar{r}_{2}-P_{2}(1+r_{F})\bigr)
    =\frac{\bar{r}_{2}-P_{2}(1+r_{F})}{\sigma_{2}^{2}}-\frac{\rho}{\sigma_{1}\sigma_{2}}\bigl(\bar{r}_{1}-P_{1}(1+r_{F})\bigr)
\]
即
\[
    \frac{\bar{r}_{1}-P_{1}(1+r_{F})}{\sigma_{1}^{2}+\rho\sigma_{1}\sigma_{2}}=\frac{\bar{r}_{2}-P_{2}(1+r_{F})}{\sigma_{2}^{2}+\rho\sigma_{1}\sigma_{2}}
\]
代入数值即得 $10-1.08P_{1}=0.77(12-1.08P_{2})$。
\end{solution}

\begin{problem}
\problemno{11.3**} 考虑一个可以以 0 成本生产煤的企业。
\begin{itemize}
    \item 煤的未来价格 $\tilde{p}$ 在当前看来是不确定的。令 $\mathrm{E}[\tilde{p}]=\bar{p}$。
    \item 销售额取决未来价格：$S(\tilde{p})$。
    \item 有一个煤的远期市场，远期价格为 $f$。企业可以用远期合约来对冲未来价格中的风险。比如，如果企业售出 $h$ 份远期合约，远期头寸的未来支付为 $h(f-\tilde{p})$。如果 $\tilde{p}$ 的实现值很低，企业从销售中得到的利润很低但是从远期头寸得到利润则很高。
    \item 公司经理是风险厌恶的，他的期望效用为
    \[
        \mathrm{E}[u(\tilde{\pi})],\qquad \tilde{\pi}=\tilde{p}S(\tilde{p})+h(f-\tilde{p})
    \]
    （$u'>0$，$u''<0$）。
    \item 市场上还存在投机者，他们是作为企业所售远期合约的买方而出现的。他们也具有风险厌恶度的偏好：$v(\cdot)$（$v'>0$，$v''<0$）。他们的效用函数为：
    \[
        \mathrm{E}\bigl[v\bigl(h(\tilde{p}-f)\bigr)\bigr]
    \]
    其中 $h$ 是他们买（卖）的远期合约数。
\end{itemize}
\begin{enumerate}[label=(\alph*)]
    \item 证明当且仅当 $f<\mathrm{E}[\tilde{p}]$（$f>\mathrm{E}[\tilde{p}]$）时投机者才会买入（卖出）远期合约。
    \item 假设 $f=\mathrm{E}[\tilde{p}]$。在什么时候企业会使用远期合约来对冲风险？它是买入还是卖出远期合约？[提示：考虑销售的弹性，$-(d\ln S)/(d\ln P)$。]
    \item 假设 $u(x)=v(x)=-e^{-ax}$，$S(p)=s_{0}+k/p$ 且 $\tilde{p}$ 具有如下分布：
    \[
        \text{概率密度函数}=e^{-(p-d)},\quad p\in[d,+\infty).
    \]
    假设共有 $N$ 个投机者。在远期市场上，企业和投机者都表现得像是在竞争市场上一样。均衡远期价格是多少？远期溢价/折扣：$f-\mathrm{E}[\tilde{p}]$ 是多少？解释所得到的结论。
\end{enumerate}
\end{problem}
\begin{solution}
\begin{enumerate}[label=(\alph*)]
    \item 投机者的最优化问题的一阶条件是
    \[
        \mathrm{E}\bigl[v'\bigl(h(\tilde{p}-f)\bigr)(\tilde{p}-f)\bigr]=0
    \]
    因而，
    \[
        \mathrm{Cov}\bigl(v'\bigl(h(\tilde{p}-f)\bigr),(\tilde{p}-f)\bigr)=-\mathrm{E}\bigl[v'\bigl(h(\tilde{p}-f)\bigr)\bigr]\mathrm{E}[\tilde{p}-f]
    \]
    而 $v'>0$，$v''<0$，因此
    \[
        h>0\,(<0)\Leftrightarrow \mathrm{E}[\tilde{p}]>(<)\,f
    \]
    即当且仅当 $f<\mathrm{E}[\tilde{p}]$（$f>\mathrm{E}[\tilde{p}]$）时投机者才会买入（卖出）远期合约。
    \item 厂商最优化问题的一阶条件是
    \[
        \mathrm{E}\bigl[u'\bigl(\tilde{p}S(\tilde{p})+h(\tilde{p}-f)\bigr)(\tilde{p}-f)\bigr]=0
    \]
    又由假设 $\mathrm{E}[\tilde{p}]=f$ 有
    \[
        \mathrm{Cov}\bigl(u'\bigl(\tilde{p}S(\tilde{p})+h(\tilde{p}-f)\bigr),(\tilde{p}-f)\bigr)=0
    \]
    令 $e(p)=-\dfrac{d\log S}{d\log p}$，那么 $d\bigl(pS(p)+h(f-p)\bigr)/dp=S(p)\bigl(1-e(p)\bigr)-h$；因此，如果 $e(p)<(>)1$，那么 $h>(<)0$。
    \item 给定 $f$，$h$，厂商的期望效用是
    \[
        \begin{aligned}
            \mathrm{E}[u(\tilde{\pi})]
            &=\mathrm{E}\bigl[u\bigl(\tilde{p}S(\tilde{p})+h(f-\tilde{p})\bigr)\bigr]\\
            &=\frac{e^{-a[S_{0}d+k+h(f-d)]}}{a(S_{0}-h)+1}
        \end{aligned}
    \]
    求导可得厂商的远期合约售出量为：$h_{1}^{*}=S_{0}+\dfrac{1}{a}-\dfrac{1}{a(f-d)}$。投机者的期望效用是
    \[
        \mathrm{E}\bigl[v\bigl(h(\tilde{p}-f)\bigr)\bigr]=\frac{e^{ah(f-d)}}{ah+1}
    \]
    求导可得投机者的远期合约需求量为：$h_{2}^{*}=\dfrac{1}{a(f-d)}-\dfrac{1}{a}$。市场出清时
    \[
        S_{0}+\frac{1}{a}-\frac{1}{a(f-d)}=N\Bigl(\frac{1}{a(f-d)}-\frac{1}{a}\Bigr)
    \]
    求得均衡远期价格为 $f^{*}=d+1/\bigl(1-\dfrac{aS_{0}}{N+1}\bigr)$，期权溢价为 $f-\mathrm{E}[\tilde{p}]=1/\bigl(1-\dfrac{aS_{0}}{N+1}\bigr)-1$。当风险厌恶系数增大时，溢价变大（或折扣变小），这是很显然的。当投机者增多时，溢价（折扣）的绝对值变小，即使得远期合约的价格向期望值靠拢，这是投机市场竞争加剧时：如果 $f<\mathrm{E}[\tilde{p}]$，此时投机者买入远期，厂商可以抬高价格；如果 $f>\mathrm{E}[\tilde{p}]$，此时投机者卖出远期，厂商可以压低价格。
\end{enumerate}
\end{solution}

\begin{problem}
\problemno{11.4*} 证券市场只含有一只风险证券（股票）。一单位股票的支付为 $\tilde{x}$，它服从均值为 $\bar{x}$、波动率为 $\sigma$ 的正态分布。市场上有参与者 1 和 2。参与者 1 的禀赋是 $x_{0}$ 单位的 0 期消费而参与者 2 的禀赋是一单位股票。两个参与者都具有如下形式的期望效用：
\[
    -e^{-\alpha c_{0}}-\rho\mathrm{E}\bigl[-e^{-\alpha\tilde{c}_{1}}\bigr].
\]
\begin{enumerate}[label=(\alph*)]
    \item 证明证券市场是不完全的。
    \item 写出均衡条件。
    \item 均衡配置是 Pareto 最优的吗？证明结论。
\end{enumerate}
\end{problem}
\begin{solution}
\begin{enumerate}[label=(\alph*)]
    \item 因为有无穷多个状态（支付是正态分布的），而只有一只证券，因此证券市场是不完全的；
    \item 均衡条件是存在股票价格 $S$ 使得
    \[
        \begin{aligned}
            \theta_{1}^{*}&=\operatorname{Argmax}_{\theta}\bigl[-e^{-\alpha(x_{0}-\theta S)}+\rho\mathrm{E}\bigl[-e^{-\alpha\theta\tilde{x}}\bigr]\bigr]\\
            \theta_{2}^{*}&=\operatorname{Argmax}_{\theta}\bigl[-e^{-\alpha(S-\theta S)}+\rho\mathrm{E}\bigl[-e^{-\alpha\theta\tilde{x}}\bigr]\bigr]\\
            \theta_{1}^{*}+\theta_{2}^{*}&=1
        \end{aligned}
    \]
    \item 参与者 1、2 的最优化问题的一阶条件分别是
    \[
        \begin{aligned}
            e^{-\alpha(x_{0}-\theta_{1}^{*}S)}S&=\rho e^{-\alpha\theta_{1}^{*}\bar{x}+\frac12\alpha^{2}(\theta_{1}^{*})^{2}\sigma^{2}}(\bar{x}-\alpha\theta_{1}^{*}\sigma^{2})\\
            e^{-\alpha(S-\theta_{2}^{*}S)}S&=\rho e^{-\alpha\theta_{2}^{*}\bar{x}+\frac12\alpha^{2}(\theta_{2}^{*})^{2}\sigma^{2}}(\bar{x}-\alpha\theta_{2}^{*}\sigma^{2})
        \end{aligned}
    \]
    如果均衡配置是 Pareto 最优的，那么两个参与者的边际效用之比应该相等，即
    \[
        \frac{-\rho e^{-\alpha\theta_{1}^{*}\tilde{x}}(-a)}{-e^{-\alpha(x_{0}-\theta_{1}^{*}S)}(-a)}=\frac{-\rho e^{-\alpha\theta_{2}^{*}\tilde{x}}(-a)}{-e^{-\alpha(S-\theta_{2}^{*}S)}(-a)}
    \]
    又因为 $\theta_{1}^{*}+\theta_{2}^{*}=1$，因而 $\theta_{1}^{*}=\theta_{2}^{*}=\tfrac12$，$S=x_{0}$。代入一阶条件有
    \[
        e^{-\frac12\alpha\tilde{x}}=\rho e^{-\frac12\alpha\tilde{x}+\frac18\alpha^{2}\sigma^{2}}(\tilde{x}-\tfrac12\alpha\sigma^{2})
    \]
    上述等式不一定成立，因而均衡配置不一定是 Pareto 最优的（实际上，只有当上述等式成立时才是 Pareto 最优的）。
\end{enumerate}
\end{solution}

\begin{problem}
\problemno{11.5**} 考虑练习 11.4 中的经济。现在引入一只无风险证券（债券），它在 1 期支付 1 单位消费品。参与者在债券上的禀赋都为 0。
\begin{enumerate}[label=(\alph*)]
    \item 证券市场完全吗？解释原因。
    \item 写出这个经济的均衡条件。
    \item 求解均衡。
    \item 比较存在无风险证券时的均衡配置和在练习 11.4 中不存在无风险证券时得到的均衡配置。它们一样吗？为什么一样或不一样？
    \item 均衡配置是 Pareto 最优的吗？证明并解释所得到的结论。
\end{enumerate}
\end{problem}
\begin{solution}
\begin{enumerate}[label=(\alph*)]
    \item 因为状态数是无穷的，但只有两只证券，因而市场不完全；
    \item 这个经济的均衡条件是存在债券价格 $B$ 和股票价格 $S$，使得
    \[
        \begin{aligned}
            \begin{bmatrix}\theta_{1,1}^{*} \\ \theta_{1,2}^{*}\end{bmatrix}&=\operatorname{Argmax}_{\{\theta_{1,1},\theta_{1,2}\}}\bigl[-e^{-\alpha(x_{0}-\theta_{1,1}B-\theta_{1,2}S)}+\rho\mathrm{E}\bigl[-e^{-\alpha(\theta_{1,1}+\theta_{1,2}\tilde{x})}\bigr]\bigr]\\
            \begin{bmatrix}\theta_{2,1}^{*} \\ \theta_{2,2}^{*}\end{bmatrix}&=\operatorname{Argmax}_{\{\theta_{2,1},\theta_{2,2}\}}\bigl[-e^{-\alpha(S-\theta_{2,1}B-\theta_{2,2}S)}+\rho\mathrm{E}\bigl[-e^{-\alpha(\theta_{2,1}+\theta_{2,2}\tilde{x})}\bigr]\bigr]\\
            \theta_{1,1}^{*}+\theta_{2,1}^{*}=0;&\quad\theta_{1,2}^{*}+\theta_{2,2}^{*}=1
        \end{aligned}
    \]
    \item 对于参与者 1 来说，最优化的一阶条件是
    \begin{align}
            -e^{-\alpha(x_{0}-\theta_{1,1}^{*}B-\theta_{1,2}^{*}S)}(\alpha B)
            &=\rho e^{-\alpha(\theta_{1,1}^{*}+\theta_{1,2}^{*}\tilde{x})+\frac12\alpha^{2}(\theta_{1,2}^{*})^{2}\sigma^{2}}(-\alpha)
            \tag{11.1}\\
            -e^{-\alpha(x_{0}-\theta_{1,1}^{*}B-\theta_{1,2}^{*}S)}(\alpha S)
            &=\rho e^{-\alpha(\theta_{1,1}^{*}+\theta_{1,2}^{*}\tilde{x})+\frac12\alpha^{2}(\theta_{1,2}^{*})^{2}\sigma^{2}}(-\alpha\bar{x}+\alpha^{2}\theta_{1,2}^{*}\sigma^{2})
            \tag{11.2}
    \end{align}
    参与者 2 的最优化的一阶条件是
    \begin{align}
            -e^{-\alpha(S-\theta_{2,1}^{*}B-\theta_{2,2}^{*}S)}(\alpha B)
            &=\rho e^{-\alpha(\theta_{2,1}^{*}+\theta_{2,2}^{*}\tilde{x})+\frac12\alpha^{2}(\theta_{2,2}^{*})^{2}\sigma^{2}}(-\alpha)
            \tag{11.3}\\
            -e^{-\alpha(S-\theta_{2,1}^{*}B-\theta_{2,2}^{*}S)}(\alpha S)
            &=\rho e^{-\alpha(\theta_{2,1}^{*}+\theta_{2,2}^{*}\tilde{x})+\frac12\alpha^{2}(\theta_{2,2}^{*})^{2}\sigma^{2}}(-\alpha\bar{x}+\alpha^{2}\theta_{2,2}^{*}\sigma^{2})
            \tag{11.4}
    \end{align}
    (11.2) 式/(11.1) 式，(11.4) 式/(11.3) 式可得
    \[
        \begin{aligned}
            \frac{S}{B}&=\bar{x}-\alpha\theta_{1,2}^{*}\sigma^{2}\\
            \frac{S}{B}&=\bar{x}-\alpha\theta_{2,2}^{*}\sigma^{2}
        \end{aligned}
    \]
    因此 $\theta_{1,2}^{*}=\theta_{2,2}^{*}=\tfrac12$。\\
    (11.3) 式/(11.1) 式可得
    \[
        e^{\alpha(x_{0}-S)+\alpha B(\theta_{2,1}^{*}-\theta_{1,1}^{*})}=e^{\alpha(\theta_{1,1}^{*}-\theta_{2,1}^{*})}
    \]
    因此，$\theta_{1,1}^{*}=-\theta_{2,1}^{*}=\dfrac{x_{0}}{2(1+B)}\widehat{S}$。利用 $\dfrac{\widehat{S}}{B}=\bar{x}-\tfrac12\alpha\sigma^{2}$，$\theta_{1,1}^{*}=-\theta_{2,1}^{*}=\dfrac{x_{0}-\widehat{S}}{2(1+B)}$ 以及 (11.1) 式就可求得债券价格 $B$ 以及股票价格 $S$。
    \item 因为存在了无风险债券，投资者的投资机会增加，使得投资者的选择不一样了；另外，这也改变了股票的投资价值，从而改变了参与者 2 的禀赋的相对值，这也使得均衡配置不一样。最明显的是在第 4 题中两个参与者的股票头寸不相等，而在本题当中则相等；
    \item 考虑两人的边际效用之间的比值
    \[
        \begin{aligned}
            &\frac{-\rho(-a)e^{-\alpha(\theta_{1,1}^{*}+\theta_{1,2}^{*}\tilde{x})}}{-(-a)e^{-\alpha(x_{0}-\theta_{1,1}^{*}B-\theta_{1,2}^{*}S)}}
            \Big/
            \frac{-\rho(-a)e^{-\alpha(\theta_{2,1}^{*}+\theta_{2,2}^{*}\tilde{x})}}{-(-a)e^{-\alpha(S-\theta_{2,1}^{*}B-\theta_{2,2}^{*}S)}}\\
            &\qquad=e^{-\alpha\bigl(x_{0}-S-(1+B)(\theta_{1,1}^{*}-\theta_{2,1}^{*})\bigr)}=1
        \end{aligned}
    \]
    因而均衡配置是 Pareto 最优的。
\end{enumerate}
\end{solution}

\begin{problem}
\problemno{11.6**} 在练习 11.5 描述的经济中，我们引入一只新的证券，一份以股票为标的资产、执行价格为 $\bar{x}$ 的期权。
\begin{enumerate}[label=(\alph*)]
    \item 现在，由股票、债券和期权所构成的证券市场是完全的吗？解释原因。
    \item 假设均衡配置和练习 11.5 中的一样。用参与者 1 的 Euler 方程为期权定价。同样地，用参与者 2 的 Euler 方程为期权定价。
    \item 使用前面得到的结论，分析练习 11.5 中得到的均衡配置是否会因为引入期权而改变。
    \item 计算期权的均衡价格。
\end{enumerate}
\end{problem}
\begin{solution}
\begin{enumerate}[label=(\alph*)]
    \item 由于状态数目大于证券数目，因而市场是不完全的；
    \item 由参与者 1 的 Euler 方程，期权价格为
    \[
        c_{1}=\mathrm{E}\Bigl[\frac{u_{1,1}'(\tilde{c}_{1,1})}{u_{1,0}'(c_{1,0})}[\tilde{x}-\bar{x}]\Bigr]
        =\mathrm{E}\Bigl[\frac{\rho e^{-\alpha(\theta_{1,1}^{*}+\theta_{1,2}^{*}\tilde{x})}}{e^{-\alpha(x_{0}-\theta_{1,1}^{*}B-\theta_{1,2}^{*}S)}}[\tilde{x}-\bar{x}]\Bigr]
    \]
    由参与者 2 的 Euler 方程，期权价格为
    \[
        c_{2}=\mathrm{E}\Bigl[\frac{u_{2,1}'(\tilde{c}_{2,1})}{u_{2,0}'(c_{2,0})}[\tilde{x}-\bar{x}]\Bigr]
        =\mathrm{E}\Bigl[\frac{\rho e^{-\alpha(\theta_{2,1}^{*}+\theta_{2,2}^{*}\tilde{x})}}{e^{-\alpha(S-\theta_{2,1}^{*}B-\theta_{2,2}^{*}S)}}[\tilde{x}-\bar{x}]\Bigr]
    \]
    \item 因为第五题中的均衡是 Pareto 最优的，因而边际效用之比相等，$c_{1}=c_{2}$，即参与者对期权的定价是一致的，因而他们不会交易期权，引入期权不会改变第 5 题中的均衡配置。
    \item 由前面的分析可得 $c=c_{1}=c_{2}$。
\end{enumerate}
\end{solution}



% !TeX root = ../../金融经济学习题讲义.tex
% 源 PDF：references/金融经济学/史上最全最完整的金融经济学答案(王江版).pdf，页 62–67（印刷页 61–66）

\zhchapter{在均值-方差偏好下的投资组合选择}

\ifdefined\WJHandoutKnowledge
\fi

% !TeX root = ../../金融经济学习题讲义.tex
% 12_均值方差偏好 课后习题（由章节文件拆分）

\begin{problem}
\problemno{12.1} 令交易证券的价格向量为 $p$、总收益率向量为 $\tilde{x}=\begin{bmatrix}\tilde{x}_{1} \\ \dots \\ \tilde{x}_{N}\end{bmatrix}$。假设 $\tilde{x}=\iota+\tilde{r}$ 服从均值为 $\iota+\bar{r}$、协方差阵为 $\Sigma$ 的正态分布。对于初始财富为 $w$ 的参与者，他的 1 期财富是 $\tilde{w}=wz^{\top}\tilde{x}$，其中 $z$ 是他的权重向量（$\iota^{\top}z=1$）。假设参与者的绝对风险厌恶系数恒为 $a$。
\begin{enumerate}[label=(\alph*)]
    \item 计算他的间接效用函数
    \[
        v=\mathrm{E}[u(\tilde{w})]
    \]
    其中 $u$ 是他的直接效用函数。
    \item 证明间接效用函数具有均值--方差表示，即它可以表示成 $\bar{w}=\mathrm{E}[\tilde{w}]$ 和 $\sigma_{w}^{2}=\mathrm{Var}[\tilde{w}]$ 的函数，即 $v(\bar{w},\sigma_{w})$。
    \item 证明均值方差效用函数 $v(\bar{w},\sigma_{w})$ 满足如下条件：$\partial_{1}v>0$，$\partial_{2}v<0$ 且 $\partial_{11}^{2}v<0$，$\partial_{22}^{2}v<0$ 且 $v$ 是凹的。
\end{enumerate}
\end{problem}
\begin{solution}
不妨假设效用函数为 $u(w)=-e^{-aw}$，那么 $v=\mathrm{E}[u(\tilde{w})]=\mathrm{E}[-e^{-a\tilde{w}}]$。因为 $\tilde{w}$ 是正态分布的，因而
\[
    v=-e^{-a\mathrm{E}\tilde{w}+\frac12 a^{2}\mathrm{Var}\tilde{w}}=-e^{-a\bar{w}+\frac12 a^{2}\sigma_{w}^{2}}=v(\bar{w},\sigma_{w})
\]
即间接效用函数具有均值方差表示，对 $v$ 求导可得
\[
    \begin{aligned}
        \partial_{1}v&=e^{-a\bar{w}+\frac12 a^{2}\sigma_{w}^{2}}a>0\\
        \partial_{2}v&=-e^{-a\bar{w}+\frac12 a^{2}\sigma_{w}^{2}}a^{2}\sigma_{w}<0\\
        \partial_{11}v&=e^{-a\bar{w}+\frac12 a^{2}\sigma_{w}^{2}}(-a^{2})<0\\
        \partial_{22}v&=-e^{-a\bar{w}+\frac12 a^{2}\sigma_{w}^{2}}(a^{4}\sigma_{w}^{2}+a^{2})<0\\
        \partial_{12}v&=e^{-a\bar{w}+\frac12 a^{2}\sigma_{w}^{2}}a^{3}\sigma_{w},\qquad
        \partial_{11}v\,\partial_{22}v-(\partial_{12}v)^{2}=v^{2}(\bar{w},\sigma_{w})a^{2}>0
    \end{aligned}
\]
因而，$v$ 还是凹的。
\end{solution}

\begin{problem}
\problemno{12.2} 考虑证券 1 和 2，它们的收益率分别是 $\tilde{r}_{1}$ 和 $\tilde{r}_{2}$。它们的波动率分别是 $\sigma_{1}$ 和 $\sigma_{2}$，并且它们的相关系数是 $\rho$。
\begin{enumerate}[label=(\alph*)]
    \item 假设 $\rho\ne 1$。用波动率度量风险，证明分散化即用证券形成组合可以降低风险。（也就是说，总是存在波动率小于两只单个证券的波动率的组合。）
    \item 找出收益率具有最小波动率的组合。
    \item 证明当 $\rho=1$ 时，分散化不能降低风险。
    \item 假设 $\sigma_{1}=\sigma_{2}$。证明无论 $\rho$ 的值是多少，等权重组合的风险是最小的。
\end{enumerate}
\end{problem}
\begin{solution}
\begin{enumerate}[label=(\alph*)]
    \item 设组合权重为 $z=\begin{bmatrix}z_{1} \\ 1-z_{1}\end{bmatrix}$，那么 $\tilde{r}_{p}=z_{1}\tilde{r}_{1}+(1-z_{1})\tilde{r}_{2}$，
    \[
        \sigma_{p}^{2}=z_{1}^{2}\sigma_{1}^{2}+2z_{1}(1-z_{1})\rho\sigma_{1}\sigma_{2}+(1-z_{1})^{2}\sigma_{2}^{2}
    \]
    对 $z_{1}$ 求导，$\partial\sigma_{p}^{2}/\partial z_{1}=0$ 给出
    \[
        z_{1}\sigma_{1}^{2}+(1-2z_{1})\rho\sigma_{1}\sigma_{2}-(1-z_{1})\sigma_{2}^{2}=0
    \]
    因而全局最小方差权重为
    \[
        z_{1}^{*}=\frac{\sigma_{2}^{2}-\rho\sigma_{1}\sigma_{2}}{\sigma_{1}^{2}-2\rho\sigma_{1}\sigma_{2}+\sigma_{2}^{2}}
    \]
    代回并配方，
    \[
        \sigma_{p}^{2}=\sigma_{2}^{2}-\frac{(\sigma_{2}^{2}-\rho\sigma_{1}\sigma_{2})^{2}}{\sigma_{1}^{2}-2\rho\sigma_{1}\sigma_{2}+\sigma_{2}^{2}}=\sigma_{1}^{2}-\frac{(\sigma_{1}^{2}-\rho\sigma_{1}\sigma_{2})^{2}}{\sigma_{1}^{2}-2\rho\sigma_{1}\sigma_{2}+\sigma_{2}^{2}}
    \]
    当 $\rho\ne 1$ 时分子为正，故 $\sigma_{p}<\min(\sigma_{1},\sigma_{2})$，分散化可以降低风险；
    \item 即上一问中的 $z^{*}$；
    \item 当 $\rho=1$ 时，$\sigma_{p}^{2}=(\sigma_{1}z_{1}+\sigma_{2}(1-z_{1}))^{2}\in[\min(\sigma_{1}^{2},\sigma_{2}^{2}),\,\max(\sigma_{1}^{2},\sigma_{2}^{2})]$，$z_{1}\in[0,1]$，即分散化不能降低风险。
    \item 当 $\sigma_{1}=\sigma_{2}$ 时，$z_{1}=(\sigma_{2}^{2}-\rho\sigma_{1}\sigma_{2})/(\sigma_{1}^{2}-2\rho\sigma_{1}\sigma_{2}+\sigma_{2}^{2})=\tfrac12$，即无论 $\rho$ 多少，等权重组合的风险最小。
\end{enumerate}
\end{solution}

\begin{problem}
\problemno{12.3*} 假设市场上有 $N$ 只证券，它们收益率波动率均为 $\sigma$，而协方差为 $c=\mathrm{Cov}[\tilde{r}_{i},\tilde{r}_{j}]$，$\forall\, i,j=1,\dots,N$ 且 $i\ne j$。
\begin{enumerate}[label=(\alph*)]
    \item 证明由 $N$ 只证券组成的等权组合的收益率具有最小波动率。
    \item 证明当 $N\to\infty$ 时等权重组合的方差趋近于 $c$。这和单个证券的波动率 $\sigma$ 无关。
    \item 解释上面得到的结论。
\end{enumerate}
\end{problem}
\begin{solution}
\begin{enumerate}[label=(\alph*)]
    \item 假设组合权重为 $z=\begin{bmatrix}z_{1} \\ z_{2} \\ \dots \\ z_{N}\end{bmatrix}$，那么组合的方差为
    \[
        \sigma_{p}^{2}=\mathrm{Var}(z_{1}\tilde{r}_{1}+z_{2}\tilde{r}_{2}+\dots+z_{N}\tilde{r}_{N})=\sigma^{2}\sum_{i}z_{i}^{2}+c\sum_{i\ne j}z_{i}z_{j}
    \]
    利用 $\sum_{i}z_{i}=1$ 改写为 $\sigma_{p}^{2}=(\sigma^{2}-c)\sum_{i}z_{i}^{2}+c$。拉格朗日函数为
    \[
        \mathcal{L}=(\sigma^{2}-c)\sum_{i}z_{i}^{2}+\lambda\Bigl(1-\sum_{i}z_{i}\Bigr)
    \]
    一阶条件 $2(\sigma^{2}-c)z_{i}=\lambda$（$\sigma^{2}-c>0$）迫使 $z_{i}$ 全相等，因此
    \[
        z_{i}=\frac{1}{N},\quad\forall i
    \]
    \item 把 $z_{i}=\dfrac{1}{N}$ 代入组合方差可得 $\sigma_{p}^{2}=\dfrac{1}{N}\sigma^{2}+\bigl(1-\dfrac{1}{N}\bigr)c\to c$，$N\to\infty$；
    \item 当 $N\to\infty$ 时，证券间不相关的风险都被分散掉了，只剩下了其中相关的部分，因而最终的组合风险只与协方差项有关。
\end{enumerate}
\end{solution}

\begin{problem}
\problemno{12.4} 考虑风险证券 1 和 2，它们的收益率为 $\tilde{r}_{1}$ 和 $\tilde{r}_{2}$。它们的期望收益率和协方差阵分别为 $\bar{r}_{1}$，$\bar{r}_{2}$，以及 $\Sigma=\begin{bmatrix}\sigma_{1}^{2} & \sigma_{12} \\ \sigma_{12} & \sigma_{2}^{2}\end{bmatrix}$。
\begin{enumerate}[label=(\alph*)]
    \item 推导由这两只证券构成的均值--方差有效前沿。
    \item 证明这两只证券都在前沿上。
    \item 考虑如下三种特别形式的前沿：$\rho_{12}=-1,\,0,\,1$，其中 $\rho_{12}$ 是两只证券收益率之间的相关系数。在 $\bar{r}$--$\sigma$ 平面画出这三种情形下的组合前沿。
\end{enumerate}
\end{problem}
\begin{solution}
\begin{enumerate}[label=(\alph*)]
    \item 为简便起见，假设 $\bar{r}_{1}\ne\bar{r}_{2}$；给定组合的期望收益率 $\bar{r}_{p}$ 有 $z\bar{r}_{1}+(1-z)\bar{r}_{2}=\bar{r}_{p}$，因而 $z=\dfrac{\bar{r}_{p}-\bar{r}_{2}}{\bar{r}_{1}-\bar{r}_{2}}$，此时，组合的波动率为
    \[
        \begin{aligned}
            \sigma_{p}^{2}&=\frac{1}{(\bar{r}_{1}-\bar{r}_{2})^{2}}\bigl[(\sigma_{1}^{2}+\sigma_{2}^{2}-2\sigma_{12})\bar{r}_{p}^{2}-(2\sigma_{1}^{2}\bar{r}_{2}+2\sigma_{2}^{2}\bar{r}_{1}-2(\bar{r}_{1}+\bar{r}_{2})\sigma_{12})\bar{r}_{p}\\
            &\qquad+(\sigma_{1}^{2}\bar{r}_{2}^{2}+\sigma_{2}^{2}\bar{r}_{1}^{2}-2\bar{r}_{1}\bar{r}_{2}\sigma_{12})\bigr]
        \end{aligned}
    \]
    $(\bar{r}_{p},\sigma_{p})$ 即为均值--方差前沿组合；
    \item 当 $\bar{r}_{p}=\bar{r}_{1}$ 时，前沿组合即为证券 1；当 $\bar{r}_{p}=\bar{r}_{2}$ 时，前沿组合即为证券 2；
    \item 当 $\rho_{12}=-1$ 时
    \[
        \sigma_{p}=\left|\frac{(\sigma_{1}+\sigma_{2})\bar{r}_{p}-(\sigma_{1}\bar{r}_{2}+\sigma_{2}\bar{r}_{1})}{\bar{r}_{1}-\bar{r}_{2}}\right|
    \]
    当 $\rho_{12}=1$ 时
    \[
        \sigma_{p}=\left|\frac{(\sigma_{1}-\sigma_{2})\bar{r}_{p}-(\sigma_{1}\bar{r}_{2}+\sigma_{2}\bar{r}_{1})}{\bar{r}_{1}-\bar{r}_{2}}\right|
    \]
    在这两种情形下，前沿组合由两条起点一直的射线构成。当 $\rho_{12}=0$ 时，
    \[
        \sigma_{p}^{2}=\frac{1}{(\bar{r}_{1}-\bar{r}_{2})^{2}}\bigl[(\sigma_{1}^{2}+\sigma_{2}^{2}-2\sigma_{12})\bar{r}_{p}^{2}-(2\sigma_{1}^{2}\bar{r}_{2}+2\sigma_{2}^{2}\bar{r}_{1})\bar{r}_{p}+(\sigma_{1}^{2}\bar{r}_{2}^{2}+\sigma_{2}^{2}\bar{r}_{1}^{2})\bigr]
    \]
    此时的前沿组合是 $\bar{r}$--$\sigma$ 平面中的一条抛物线。
\end{enumerate}
\end{solution}

\begin{problem}
\problemno{12.5} 令市场上风险证券的收益率向量为 $\tilde{r}$，它的均值为 $\bar{r}$，它的协方差阵为 $\Sigma$，$\Sigma$ 是正定的，而 $\iota$ 为同阶的元素全为 1 的向量。另外，令 $a=\iota^{\top}\Sigma^{-1}\bar{r}$，$b=\bar{r}^{\top}\Sigma^{-1}\bar{r}$，$c=\iota^{\top}\Sigma^{-1}\iota$ 且 $d=bc-a^{2}$。证明 $d=0$ 等价于 $\bar{r}=(b/a)\iota$，即所有风险证券具有相同的期望收益 $(b/a)$。
\end{problem}
\begin{solution}
当 $d=0$ 时，$bd=(a\bar{r}-b\iota)^{\top}\Sigma^{-1}(a\bar{r}-b\iota)=0$，因为 $\Sigma$，$\Sigma^{-1}$ 都是正定的，因而 $a\bar{r}-b\iota=0$，$\bar{r}=\dfrac{b}{a}\iota$，即所有风险证券具有相同的期望收益 $\dfrac{b}{a}$。
\end{solution}

\begin{problem}
\problemno{12.6} 证明最小方差组合（MVP）的期望收益率为 $a/c$，这里的 $a$ 和 $c$ 由练习 12.5 给定。
\end{problem}
\begin{solution}
我们知道最小方差组合的期望收益率为 $r_{\mathrm{MVP}}=-\dfrac{(z_{1}-z_{0})^{\top}\Sigma^{-1}z_{0}}{(z_{1}-z_{0})^{\top}\Sigma^{-1}(z_{1}-z_{0})}$，$z_{0}$，$z_{1}$ 由正文 2.3 节定义，把 $z_{0}$，$z_{1}$ 代入得
\[
    \begin{aligned}
        r_{\mathrm{MVP}}
        &=-\frac{d^{-2}(c\Sigma^{-1}\bar{r}-a\Sigma^{-1}\iota)^{\top}\Sigma(b\Sigma^{-1}\iota-a\Sigma^{-1}\bar{r})}{d^{-2}(c\Sigma^{-1}\bar{r}-a\Sigma^{-1}\iota)^{\top}\Sigma(c\Sigma^{-1}\bar{r}-a\Sigma^{-1}\iota)}\\
        &=-\frac{a(a^{2}-bc)}{c(bc-a^{2})}=\frac{a}{c}
    \end{aligned}
\]
\end{solution}

\begin{problem}
\problemno{12.7} 考虑 $\bar{r}=(b/a)\iota$ 的特殊情形。导出均值--方差前沿。在 $\bar{r}$--$\sigma$ 平面上画出这个组合前沿。
\end{problem}
\begin{solution}
当期望收益率都相等时，所有组合的期望收益率都是常数，因而均值--方差前沿组合退化成为一点即最小方差组合，此时的方差为 $\sigma_{\mathrm{MVP}}^{2}=\dfrac{1}{\iota^{\top}\Sigma^{-1}\iota}$，对应的组合权重为 $z_{\mathrm{MVP}}=\dfrac{\Sigma^{-1}\iota}{\iota^{\top}\Sigma^{-1}\iota}$。
\end{solution}

\begin{problem}
\problemno{12.8**} 令证券收益率向量为 $\tilde{r}$，它的均值为 $\bar{r}$、协方差阵为 $\Sigma$。在 $\bar{r}$--$\sigma$ 平面上，任一组合 $p$ 由一点 $(\bar{r}_{p},\sigma_{p})$ 给定。令所有可能组合的集合为 $A=\{(\bar{r}_{p},\sigma_{p}),\forall\, p\}$。
\begin{enumerate}[label=(\alph*)]
    \item 证明 $A$ 是闭集。
    \item 证明 $A$ 是凸集。
\end{enumerate}
\end{problem}
\begin{solution}
\begin{enumerate}[label=(\alph*)]
    \item 假设有一序列 $\bigl\{(r_{p}^{(n)},\sigma_{p}^{(n)})\bigr\}_{1}^{\infty}$，其极限为 $(r_{p},\sigma_{p})$。那么存在 $\theta^{(n)}\in S^{N-1}=\{\theta\in\mathrm{R}^{N}:\theta^{\top}\iota=1\}$ 使得
    \[
        \theta^{(n)\top}\bar{r}=r_{p}^{(n)},\qquad\theta^{(n)\top}\Sigma\theta^{(n)}=\sigma_{p}^{(n)2}
    \]
    显然，$\theta^{(n)}$ 是 $S^{N-1}$ 中的有界序列，存在收敛子列 $\{\theta^{(k_{n})}\}_{1}^{\infty}$，记其极限为 $\theta\in S^{N-1}$，那么
    \[
        \begin{aligned}
            \theta^{\top}\bar{r}&=\lim_{n}\theta^{(k_{n})\top}\bar{r}=\lim_{n}r_{p}^{(k_{n})}=r_{p}\\
            \theta^{\top}\Sigma\theta&=\lim_{n}\theta^{(k_{n})\top}\Sigma\theta^{(k_{n})}=\lim_{n}\sigma_{p}^{(k_{n})2}=\sigma_{p}^{2}
        \end{aligned}
    \]
    因而 $(r_{p},\sigma_{p})\in A$，$A$ 是闭集；
    \item 对于 $(r_{1},\sigma_{1})$，$(r_{2},\sigma_{2})$，我们可以按照以下的规则选择相应的权重向量 $\theta$。
    \[
        \begin{aligned}
            \theta_{1},\,\theta_{2}&\in S^{N-1}\\
            \theta_{1}^{\top}\bar{r}&=r_{1},\quad\theta_{1}^{\top}\Sigma\theta_{1}=\sigma_{1}^{2}\\
            \theta_{2}^{\top}\bar{r}&=r_{2},\quad\theta_{2}^{\top}\Sigma\theta_{2}=\sigma_{2}^{2}\\
            \theta_{1}^{\top}\Sigma\theta_{2}&=\sigma_{1}\sigma_{2}
        \end{aligned}
    \]
    那么，对于任意的 $\alpha\in[0,1]$，$\bigl(\alpha r_{1}+(1-\alpha)r_{2},\alpha\sigma_{1}+(1-\alpha)\sigma_{2}\bigr)$ 对应的权重向量是 $\alpha\theta_{1}+(1-\alpha)\theta_{2}$，因此 $\alpha r_{1}+(1-\alpha)r_{2}$，$\alpha\sigma_{1}+(1-\alpha)\sigma_{2}\in A$，$A$ 是凸集。
\end{enumerate}
Notes：凸性的命题不一定成立，比如当 $N=2$ 时，$A$ 是 $\bar{r}$--$\sigma$ 平面中的一条抛物线，显然不是凸的。
\end{solution}

\begin{problem}
\problemno{12.9**} 令 $p$ 为 MVF 组合，$zcp$ 是与 $p$ 的协方差为零的 MVF。
\begin{enumerate}[label=(\alph*)]
    \item 证明对任意证券 $i$，我们有
    \[
        \tilde{r}_{i}=\alpha+\beta_{i,zcp}\tilde{r}_{zcp}+\beta_{i,p}\tilde{r}_{p}+\tilde{u}_{i}
    \]
    其中 $\mathrm{Cov}[\tilde{r}_{zcp},\tilde{u}_{i}]=0$，$\mathrm{Cov}[\tilde{r}_{p},\tilde{u}_{i}]=0$，$\mathrm{E}[\tilde{u}_{i}]=0$，$\beta_{i,zcp}=\mathrm{Cov}[\tilde{r}_{i},\tilde{r}_{zcp}]/\mathrm{Var}[\tilde{r}_{zcp}]$ 且 $\beta_{i,p}=\mathrm{Cov}[\tilde{r}_{i},\tilde{r}_{p}]/\mathrm{Var}[\tilde{r}_{p}]$。
    \item 证明 $\beta_{i,zcp}=1-\beta_{i,p}$。
    \item 证明 $\alpha=0$ 及 $\bar{r}_{i}=(1-\beta_{i,p})\bar{r}_{zcp}+\beta_{i,p}\bar{r}_{p}$。
\end{enumerate}
\end{problem}
\begin{solution}
\begin{enumerate}[label=(\alph*)]
    \item 令 $\tilde{r}_{i}=\alpha+\beta_{i,1}\tilde{r}_{zcp}+\beta_{i,2}\tilde{r}_{p}+\tilde{u}_{i}$，由
    \[
        \mathrm{Cov}(\tilde{u}_{i},\tilde{r}_{p})=\mathrm{Cov}(\tilde{r}_{i}-\alpha-\beta_{i,1}\tilde{r}_{zcp}-\beta_{i,2}\tilde{r}_{p},\tilde{r}_{p})=0
    \]
    可得 $\beta_{i,2}=\beta_{i,p}$；同样地有 $\beta_{i,1}=\beta_{i,zcp}$，由 $\mathrm{E}\tilde{u}_{i}=0$ 可以得到 $\alpha=\bar{r}_{i}-\beta_{i,p}\bar{r}_{p}-\beta_{i,zcp}\bar{r}_{zcp}$；
    \item 由文中 (12.6) 式有
    \[
        \begin{aligned}
            \bar{r}_{i}-\bar{r}_{zcp}&=\beta_{i,p}(\bar{r}_{p}-\bar{r}_{zcp})\\
            \bar{r}_{i}-\bar{r}_{p}&=\beta_{i,zcp}(\bar{r}_{zcp}-\bar{r}_{p})
        \end{aligned}
    \]
    两式相减可以得到 $\beta_{i,p}+\beta_{i,zcp}=1$；
    \item 另外，由 (12.6) 式有
    \[
        \bar{r}_{i}=\bar{r}_{zcp}+\beta_{i,p}(\bar{r}_{p}-\bar{r}_{zcp})=\beta_{i,p}\bar{r}_{p}+(1-\beta_{i,p})\bar{r}_{zcp}
    \]
    因此 $\alpha=0$。
\end{enumerate}
\end{solution}

\begin{problem}
\problemno{12.10} 证明 MVE 组合的凸组合（即权重非负的组合）仍然是 MVE 组合。
\end{problem}
\begin{solution}
首先，我们知道 MVE 组合的组合是 MVF 组合；另外，由于是凸组合，因而所得 MVF 组合在均值--方差前沿的上半部，即为 MVE 组合。
\end{solution}

\begin{problem}
\problemno{12.11} 当存在一只收益率为 $r_{F}$ 的无风险债券时，记 $\bar{\eta}=\bar{r}-r_{F}\iota$ 是风险证券的超额收益向量。证明 $\bar{\eta}^{\top}\Sigma^{-1}\bar{\eta}=0$ 等价于 $\bar{r}=r_{F}\iota$。
\end{problem}
\begin{solution}
因为 $\Sigma^{-1}$ 是正定的，
\[
    \bar{\eta}^{\top}\Sigma^{-1}\bar{\eta}=0\Leftrightarrow\bar{\eta}=\bar{r}-r_{F}\iota=0\Leftrightarrow\bar{r}=r_{F}\iota
\]
\end{solution}

\begin{problem}
\problemno{12.12} 当存在无风险证券时，切点组合的期望超额收益为正等价于 $\bar{r}_{\mathrm{mvp}}-r_{F}>0$，即最小方差组合的期望收益为正。
\end{problem}
\begin{solution}
当切点组合的期望收益率为正时，
\[
    \begin{aligned}
        \iota^{\top}\Sigma^{-1}\bar{\eta}>0
        &\Leftrightarrow\iota^{\top}\Sigma^{-1}(\bar{r}-r_{F}\iota)=\iota^{\top}\Sigma^{-1}\bar{r}-r_{F}\iota^{\top}\Sigma^{-1}\iota=\iota^{\top}\Sigma^{-1}\iota\Bigl(\frac{\iota^{\top}\Sigma^{-1}\bar{r}}{\iota^{\top}\Sigma^{-1}\iota}-r_{F}\Bigr)\\
        &=\iota^{\top}\Sigma^{-1}\iota\Bigl(\frac{a}{c}-r_{F}\Bigr)=\iota^{\top}\Sigma^{-1}\iota(\bar{r}_{\mathrm{mvp}}-r_{F})>0\\
        &\Leftrightarrow\bar{r}_{\mathrm{mvp}}-r_{F}>0.
    \end{aligned}
\]
\end{solution}

\begin{problem}
\problemno{12.13*} 当存在无风险证券时，证明切点组合的 Sharpe 比最大。
\end{problem}
\begin{solution}
考虑任意组合 Sharpe 比的平方
\[
    \frac{z^{\top}\bar{r}+(1-z^{\top}\iota)r_{F}-r_{F}}{z^{\top}\Sigma z}=\frac{z^{\top}\bar{\eta}}{z^{\top}\Sigma z}
\]
其中 $z$ 是风险组合的权重，最大化的一阶条件是
\[
    \bar{\eta}\,z^{\top}\Sigma z=z\,\bar{\eta}^{\top}\Sigma z
\]
可以验证 $z_{T}=\dfrac{1}{\iota^{\top}\Sigma^{-1}\bar{\eta}}\Sigma^{-1}\bar{\eta}$ 即为其解，因而切点组合的 Sharpe 比最大。
\end{solution}



% !TeX root = ../../金融经济学习题讲义.tex
% 源 PDF：references/金融经济学/史上最全最完整的金融经济学答案(王江版).pdf，页 68–75（印刷页 67–74）

\zhchapter{资本资产定价模型（CAPM）}

\ifdefined\WJHandoutKnowledge
\fi

% !TeX root = ../../金融经济学习题讲义.tex
% 13_CAPM 课后习题（由章节文件拆分）

\begin{problem}
\problemno{13.1*} 如下的经济有两个时期，0 和 1：
\begin{itemize}
    \item 有 $N$ 只风险证券（股票）和一只无风险证券。
    \item 无风险利率为 0。
    \item 风险证券的 1 期支付向量为 $\tilde{v}$，$\tilde{v}$ 是均值为 $\bar{v}$、协方差阵为 $\Sigma$ 的 $N$ 维正态随机变量。假设 $\Sigma$ 是正定的。
    \item 股票市场组合的 1 期支付为 $\tilde{m}=\iota^{\mathrm{T}}\tilde{v}$，其中 $\iota$ 是元素全为 1 的 $(N\times 1)$ 列向量。
    \item 经济中有 $K$ 个投资者，所有投资者都具有相同的 CARA 偏好：$u(w)=-e^{-w}$。
    \item 每一投资者的禀赋是 $\bar{\theta}_{k}$ 股股票，$k=1,\cdots,K$。（注意 $\bar{\theta}_{k}$ 是由股票份额组成的列向量）且 $\sum_{k}\bar{\theta}_{k}=\iota$。
    \item 每一投资者在 1 期也有劳动收入 $y_{k}(\tilde{m}+\tilde{\varepsilon})$，其中 $\tilde{\varepsilon}$ 服从均值为 0、方差为 $\sigma^{2}$ 的正态分布，而 $y_{k}$ 是给定的。
\end{itemize}
假设 $\tilde{\varepsilon}$ 和 $\tilde{v}$ 不相关。
\begin{enumerate}[label=(\alph*)]
    \item 求解投资者（在给定价格向量下）的股票需求函数。
    \item 求解均衡股票价格。
    \item 证明 CAPM 成立。
    \item 均衡配置是 Pareto 最优吗？为什么？证明答案。如果 $\tilde{\varepsilon}=0$，则答案如何变化？[提示：在 Pareto 最优配置当中，投资者的相对边际效用在任何状态下都是相等的。]
\end{enumerate}
\end{problem}
\begin{solution}
\begin{enumerate}[label=(\alph*)]
    \item 假设股票价格向量为 $S$，参与者持有的股票头寸为 $\theta_{k}$，那么参与者 $k$ 的优化问题是
    \[
        \begin{aligned}
            \max_{\theta_{k}}\;&\mathrm{E}[-e^{-\tilde{w}_{k}}]\\
            \text{s.t.}\;&\tilde{w}_{k}=w_{k}+\bar{\theta}_{k}^{\mathrm{T}}S-\theta_{k}^{\mathrm{T}}S+\theta_{k}^{\mathrm{T}}\tilde{v}+y_{k}(\iota^{\mathrm{T}}\tilde{v}+\tilde{\varepsilon})
        \end{aligned}
    \]
    CARA 且 $\tilde{v}$、$\tilde{\varepsilon}$ 正态，最大化 $E[-e^{-\tilde{w}_{k}}]$ 等价于最大化 $E[\tilde{w}_{k}]-\tfrac12\operatorname{Var}(\tilde{w}_{k})$。记常数项 $A_{k}=w_{k}+\bar{\theta}_{k}^{\mathrm{T}}S$，则
    \[
        \begin{aligned}
            E[\tilde{w}_{k}]&=A_{k}+\theta_{k}^{\mathrm{T}}(\bar{v}-S)+y_{k}\iota^{\mathrm{T}}\bar{v}\\
            \operatorname{Var}(\tilde{w}_{k})&=(\theta_{k}+y_{k}\iota)^{\mathrm{T}}\Sigma(\theta_{k}+y_{k}\iota)+y_{k}^{2}\sigma^{2}
        \end{aligned}
    \]
    对 $\theta_{k}$ 求导：$-(\bar{v}-S)+\Sigma(\theta_{k}+y_{k}\iota)=0$，解得
    \[
        \theta_{k}^{*}=\Sigma^{-1}(\bar{v}-S)-y_{k}\iota
    \]

    \item 市场出清 $\sum_{k}\theta_{k}^{*}=\iota$ 给出
    \[
        K\Sigma^{-1}(\bar{v}-S)-\Bigl(\sum_{k}y_{k}\Bigr)\iota=\iota
    \]
    即 $\bar{v}-S=\dfrac{1}{K}\bigl(\sum_{k}y_{k}+1\bigr)\Sigma\iota$，因而股价为
    \[
        S=\bar{v}-\frac{1}{K}\Bigl(\sum y_{k}+1\Bigr)\Sigma\iota
    \]

    \item 因为
    \[
        \begin{aligned}
            \frac{\bar{r}_{i}}{\beta_{i}\bar{r}_{m}}
            &=\Bigl(\frac{\bar{v}_{i}}{S_{i}}-1\Bigr)\Big/\Biggl[\frac{\operatorname{Cov}(\tilde{v}_{i},\,\iota^{\mathrm{T}}\tilde{v})(\iota^{\mathrm{T}}S)}{\operatorname{Var}(\iota^{\mathrm{T}}\tilde{v})S_{i}}\Bigl(\frac{\iota^{\mathrm{T}}\bar{v}}{\iota^{\mathrm{T}}S}-1\Bigr)\Biggr]\\
            &=(\bar{v}_{i}-S_{i})\big/\bigl[\operatorname{Cov}(\tilde{v}_{i},\,\iota^{\mathrm{T}}\tilde{v})(\iota^{\mathrm{T}}\bar{v}-\iota^{\mathrm{T}}S)\bigr]\\
            &=\frac{1}{K}\Bigl(\sum y_{k}+1\Bigr)(\Sigma\iota)_{i}\Big/\frac{(\Sigma\iota)_{i}\bigl(\tfrac{1}{K}(\sum y_{k}+1)(\iota^{\mathrm{T}}\Sigma\iota)\bigr)}{\iota^{\mathrm{T}}\Sigma\iota}\\
            &=1
        \end{aligned}
    \]
    即 $\bar{r}_{i}=\beta_{i}\bar{r}_{m}$，因而 CAPM 成立；

    \item 对于参与者 $k$ 来说，1 期边际效用为：
    \[
        e^{-\tilde{w}_{k}}=e^{-w_{k}+\bar{\theta}_{k}^{\mathrm{T}}S+(\Sigma^{-1}(\bar{v}-S))^{\mathrm{T}}(\bar{v}-S)+y_{k}\tilde{\varepsilon}}
    \]
    除非 $y_{k}$ 全部相等，否则边际效用之比（1 期边际效用除以 0 期边际效用）是不会相等的，因而均衡配置不是 Pareto 最优的；如果 $\tilde{\varepsilon}=0$，那么参与者的边际效用之比最多相差一个常数，但是在均衡条件下，他们对股票的定价是相同的，因而边际效用之比相等，均衡配置是 Pareto 最优的。
\end{enumerate}
\end{solution}

\begin{problem}
\problemno{13.2**} 在练习 13.1 定义的经济中，现在假设 $\tilde{\varepsilon}$ 与 $\tilde{v}$ 相关。令 $\operatorname{Cov}[\tilde{v},\tilde{\varepsilon}]=c$，这里 $c$ 是一 $(N\times 1)$ 维向量。
\begin{enumerate}[label=(\alph*)]
    \item 求解投资者的股票需求函数。是否得到两基金分离？证明答案。如果两基金分离不成立，能否得到 $K$ 基金分离？$K$ 应为多少？
    \item 求解均衡股票价格。
    \item CAPM 成立吗？解释它为什么成立或不成立。列出导致不同股票的期望收益不同的原因并做出解释。
    \item 讨论均衡配置的 Pareto 最优性。
\end{enumerate}
\end{problem}
\begin{solution}
\begin{enumerate}[label=(\alph*)]
    \item 用和第 1 题类似的方法可以求得参与者 $k$ 的股票需求函数为
    \[
        \theta_{k}=\Sigma^{-1}(\bar{v}-S)-y_{k}(\iota+c)
    \]
    此时两基金分离不成立，但是三基金分离成立：一是组合头寸向量为 $\Sigma^{-1}(\bar{v}-S)$ 的股票组合，另一是组合头寸向量为 $\iota+c$ 的股票组合，最后是无风险债券；

    \item 类似地，可以求得均衡股价为
    \[
        S=\bar{v}-\frac{1}{K}\Biggl[\Bigl(\sum y_{k}+1\Bigr)\Sigma\iota+\sum y_{k}\Sigma c\Biggr]
    \]

    \item 考虑
    \[
        \begin{aligned}
            \frac{\bar{r}_{i}}{\beta_{i}\bar{r}_{m}}
            &=\Bigl(\frac{\bar{v}_{i}}{S_{i}}-1\Bigr)\Big/\Biggl[\frac{\operatorname{Cov}(\tilde{v}_{i},\,\iota^{\mathrm{T}}\tilde{v})(\iota^{\mathrm{T}}S)}{\operatorname{Var}(\iota^{\mathrm{T}}\tilde{v})S_{i}}\Bigl(\frac{\iota^{\mathrm{T}}\bar{v}}{\iota^{\mathrm{T}}S}-1\Bigr)\Biggr]\\
            &=(\bar{v}_{i}-S_{i})\big/\bigl[\operatorname{Cov}(\tilde{v}_{i},\,\iota^{\mathrm{T}}\tilde{v})(\iota^{\mathrm{T}}\bar{v}-\iota^{\mathrm{T}}S)\bigr]\\
            &=\frac{1}{K}\Biggl(\Bigl(\sum y_{k}+1\Bigr)(\Sigma\iota)+\sum y_{k}\Sigma c\Biggr)_{i}\Big/\frac{(\Sigma\iota)_{i}\bigl(\tfrac{1}{K}\bigl\{(\sum y_{k}+1)(\iota^{\mathrm{T}}\Sigma\iota)+\sum y_{k}\Sigma c\bigr\}\bigr)}{\iota^{\mathrm{T}}\Sigma\iota}\\
            &\ne 1
        \end{aligned}
    \]
    因而 CAPM 成立；这是因为需要用股票来对冲与股票相关的其它风险。股票 $i$ 的期望收益率为
    \[
        \bar{r}_{i}=\frac{\bar{v}_{i}}{S_{i}}-1=\frac{1}{1-\frac{1}{K\bar{v}_{i}}\bigl((\sum y_{k}+1)(\Sigma\iota)+\sum y_{k}\Sigma c\bigr)_{i}}-1
    \]
    因而导致股票期望收益率不同的因素有：
    \begin{itemize}
        \item 期望支付；
        \item 与其它股票的相关性以及自身的风险（方差）；
        \item 与收入中的风险 $\tilde{\varepsilon}$ 的相关性。
    \end{itemize}

    \item 考虑参与者 $k$ 的 1 期边际效用：
    \[
        e^{-\tilde{w}_{k}}=e^{-w_{k}+\bar{\theta}_{k}^{\mathrm{T}}S+(\Sigma^{-1}(\bar{v}-S))^{\mathrm{T}}(\bar{v}-S)+y_{k}(\tilde{\varepsilon}-c^{\mathrm{T}}\tilde{v})}
    \]
    除非 $\tilde{\varepsilon}=c^{\mathrm{T}}\tilde{v}$ 或者 $y_{k}$ 全部相等，否则边际效用之比（1 期边际效用除以 0 期边际效用）是不会相等的，均衡配置也不是 Pareto 最优的。
\end{enumerate}
\end{solution}

\begin{problem}
\problemno{13.3**} 考虑一个由两个国家（$a$ 和 $b$）组成的、两个时期（0 和 1）的经济。这两个国家使用相同的货币。
\begin{itemize}
    \item 两个国家中的无风险资产是一样的，无风险利率为 $r_{F}$。
    \item 国家 $a$ 中有 $N$ 只风险证券，国家 $b$ 中也有 $N$ 只风险证券。$\tilde{v}_{a}=\bar{v}+\tilde{z}_{a}$、$\tilde{v}_{b}=\bar{v}+\tilde{z}_{b}$ 分别是 $a$-资产、$b$-资产的 1 期支付。假设 $\bar{v}\sim N(0,\sigma_{w}^{2}I)$，这里 $I$ 是 $N\times N$ 单位阵，$\tilde{v}_{k}\sim N(\bar{v},\Sigma_{k})$，$k=a,b$，且 $\bar{v}$、$\tilde{z}_{a}$ 和 $\tilde{z}_{b}$ 是独立的。
    \item 每一风险资产的总供给都为 1。
    \item 国家 $a$ 有 $K_{a}$ 个投资者而国家 $b$ 有 $K_{b}$ 个投资者。所有投资都具有风险厌恶系数为 $\alpha$、依赖于 1 期财富的 CARA 效用函数。
\end{itemize}
令 $p_{a}$ 和 $p_{b}$ 为国家 $a$ 和国家 $b$ 的资产价格向量。
\begin{enumerate}[label=(\alph*)]
    \item 假设投资者可以持有两个国家的资产。计算国家 $a$ 和国家 $b$ 的投资者的资产需求。
    \item 计算均衡价格。
    \item 计算所有资产的均衡风险溢价。
    \item 证明国际化的 CAPM 成立，即
    \[
        \bar{r}-r_{F}=\beta(\bar{r}_{wm}-r_{F}),\quad(\bar{r}=\mathrm{E}[\tilde{r}]),
    \]
    其中 $\beta$ 是相对于世界市场组合的 beta 值，世界市场组合是由国家 $a$ 和国家 $b$ 的市场组合构成的等权重组合（也即两国所有资产的组合）。
\end{enumerate}
\end{problem}
\begin{solution}
\begin{enumerate}[label=(\alph*)]
    \item 考虑国家 $a$ 中的参与者 $k$，他的优化问题是
    \[
        \begin{aligned}
            \max_{\theta_{k,a},\,\theta_{k,b}}\;&\mathrm{E}[-e^{-\alpha\tilde{w}_{k}}]\\
            \text{s.t.}\;&\tilde{w}_{k}=(w_{k}-\theta_{k,a}P_{a}-\theta_{k,b}P_{b})(1+r_{F})+\theta_{k,a}^{\mathrm{T}}\tilde{v}_{a}+\theta_{k,b}^{\mathrm{T}}\tilde{v}_{b}
        \end{aligned}
    \]
    求得最优的股票需求量为
    \[
        \begin{aligned}
            \theta_{k,a}&=\frac{1}{\alpha}\Sigma_{a}^{-1}(\bar{v}-P_{a}(1+r_{F}))-\frac{\sigma_{w}^{2}}{\alpha}\Sigma_{a}^{-1}(\sigma_{w}^{2}(\Sigma_{a}^{-1}+\Sigma_{b}^{-1})+I)^{-1}\\
            &\quad\times\bigl(\Sigma_{a}^{-1}(\bar{v}-P_{a}(1+r_{F}))+\Sigma_{b}^{-1}(\bar{v}-P_{b}(1+r_{F}))\bigr)\\
            \theta_{k,b}&=\frac{1}{\alpha}\Sigma_{b}^{-1}(\bar{v}-P_{b}(1+r_{F}))-\frac{\sigma_{w}^{2}}{\alpha}\Sigma_{b}^{-1}(\sigma_{w}^{2}(\Sigma_{a}^{-1}+\Sigma_{b}^{-1})+I)^{-1}\\
            &\quad\times\bigl(\Sigma_{a}^{-1}(\bar{v}-P_{a}(1+r_{F}))+\Sigma_{b}^{-1}(\bar{v}-P_{b}(1+r_{F}))\bigr)
        \end{aligned}
    \]
    类似地可以求得国家 $b$ 中参与者 $k$ 的需求函数有
    \[
        \theta_{k,a}'=\theta_{k,a},\qquad\theta_{k,b}'=\theta_{k,b}
    \]

    \item 市场出清条件是
    \[
        \begin{aligned}
            \sum\theta_{k,a}+\sum\theta_{k,a}'&=\iota\\
            \sum\theta_{k,b}+\sum\theta_{k,b}'&=\iota
        \end{aligned}
    \]
    因而可以求得均衡股价为
    \[
        \begin{aligned}
            P_{a}(1+r_{F})&=\bar{v}-\frac{\alpha}{K}\sigma_{w}^{2}\iota-\frac{\alpha}{2K}\Sigma_{a}\iota\\
            P_{b}(1+r_{F})&=\bar{v}-\frac{\alpha}{K}\sigma_{w}^{2}\iota-\frac{\alpha}{2K}\Sigma_{b}\iota
        \end{aligned}
    \]

    \item 国家 $a$ 中资产 $k$ 的风险溢价为
    \[
        \bar{r}_{k}-r_{F}=\frac{\bar{v}_{k}}{P_{a,k}}-(1+r_{F})=\frac{\frac{\alpha}{2K}(2\sigma_{w}^{2}\iota+\Sigma_{a}\iota)_{k}(1+r_{F})}{\bar{v}_{k}-\frac{\alpha}{2K}(2\sigma_{w}^{2}\iota+\Sigma_{a}\iota)_{k}}
    \]
    同样地，对于国家 $b$ 中资产 $i$ 有
    \[
        \bar{r}_{i}-r_{F}=\frac{\bar{v}_{i}}{P_{b,i}}-(1+r_{F})=\frac{\frac{\alpha}{2K}(2\sigma_{w}^{2}\iota+\Sigma_{a}\iota)_{i}(1+r_{F})}{\bar{v}_{i}-\frac{\alpha}{2K}(2\sigma_{w}^{2}\iota+\Sigma_{a}\iota)_{i}}
    \]

    \item 对于国家 $a$ 中的资产 $k$，
    \[
        \begin{aligned}
            \beta_{k}&=\operatorname{Cov}\Bigl(\frac{\tilde{v}_{k}}{P_{a,k}}-1,\,\frac{\iota^{\mathrm{T}}(\tilde{v}_{a}+\tilde{v}_{b})}{\iota^{\mathrm{T}}(P_{a}+P_{b})}-1\Bigr)\Big/\operatorname{Var}\Bigl(\frac{\iota^{\mathrm{T}}(\tilde{v}_{a}+\tilde{v}_{b})}{\iota^{\mathrm{T}}(P_{a}+P_{b})}\Bigr)\\
            &=\frac{(2\sigma_{w}^{2}\iota+\Sigma_{a}\iota)_{k}(\iota^{\mathrm{T}}(P_{a}+P_{b}))}{\iota^{\mathrm{T}}(4\sigma_{w}^{2}I+\Sigma_{a}+\Sigma_{b})\iota}\\
            \bar{r}_{wm}&=\frac{\tfrac12\iota^{\mathrm{T}}(\tilde{v}_{a}+\tilde{v}_{b})}{\iota^{\mathrm{T}}(P_{a}+P_{b})}\\
            \bar{r}_{wm}-r_{F}&=\frac{\frac{\alpha}{2K}\iota^{\mathrm{T}}(4\sigma_{w}^{2}I+\Sigma_{a}+\Sigma_{b})\iota}{\iota^{\mathrm{T}}(P_{a}+P_{b})}
        \end{aligned}
    \]
    因此，$\bar{r}_{k}-r_{F}=\beta_{k}(\bar{r}_{wm}-r_{F})$，可以验证对于国家 $b$ 中的资产此式也成立，因而国际化的 CAPM 成立。
\end{enumerate}
\end{solution}

\begin{problem}
\problemno{13.4**} 继续考虑练习 13.3 所描述的经济。假设国家 $a$ 中的投资者可以持有国家 $b$ 的资产而国家 $b$ 中的投资者不能持有国家 $a$ 的资产。
\begin{enumerate}[label=(\alph*)]
    \item 重复 (a)--(c) 中的计算。
    \item 现在证明
    \[
        \bar{r}-r_{F}=\gamma_{a}(\bar{r}_{a}-r_{F})+\gamma_{b}(\bar{r}_{b}-r_{F}).
    \]
    其中 $\bar{r}_{a}$ 和 $\bar{r}_{b}$ 分别是国家 $a$ 和 $b$ 的市场组合的期望收益。计算 $\gamma_{a}$ 和 $\gamma_{b}$ 并讨论所得到的结论。
\end{enumerate}
\end{problem}
\begin{solution}
略
\end{solution}

\begin{problem}
\problemno{13.5*} 考虑如下的具有两个时期 0 和 1 的 Arrow--Debreu 经济：
\begin{itemize}
    \item 单一的消费品（也作为计量单位）。
    \item 经济在 $t=1$ 期的状态为 $\omega\in\Omega=\{1,\cdots,\Omega\}$。
    \item $\pi_{\omega}>0$ 是状态 $\omega$，$\omega\in\Omega$ 的概率。
    \item 证券市场上存在 Arrow--Debreu 状态或有证券的完全集合。令 $\phi_{\omega}$ 为状态 $\omega$ 或有要求权（只在状态 $\omega$ 下支付 1 单位消费品）的价格。
    \item 总禀赋是 $\{e_{\omega}:\omega\in\Omega\}$，假设 $e_{\omega}>0$，$\forall\,\omega\in\Omega$。
\end{itemize}
\begin{enumerate}[label=(\alph*)]
    \item 写出 $\phi_{\omega}$，$\omega\in\Omega$ 的无套利条件。
    \item 一个状态或有要求权的组合是 $w=\begin{bmatrix}w_{1} \\ \cdots \\ w_{\omega} \\ \cdots \\ w_{\Omega}\end{bmatrix}$，其中 $w_{\omega}$ 是投资于状态 $\omega$ 状态或有要求权的份额，$\sum_{\omega}w_{\omega}=1$。
    \begin{itemize}
        \item 计算组合收益率的均值和方差。
        \item 计算两个组合 $w$ 和 $w'$ 之间的协方差。
    \end{itemize}
    \item 考虑如下的两个组合 $w^{a}$ 和 $w^{b}$
    \[
        \forall\,\omega\in\Omega:\quad w_{\omega}^{a}=\frac{\phi_{\omega}}{\sum_{\omega}\phi_{\omega}},\quad w_{\omega}^{b}=\frac{\phi_{\omega}q_{\omega}}{\sum_{\omega}\phi_{\omega}q_{\omega}}
    \]
    其中 $q_{\omega}=\phi_{\omega}/\pi_{\omega}$。
    \begin{itemize}
        \item 验证 $w^{a}$ 是 $t=0$ 期价格为 1 的无风险组合。
        \item 证明所有的前沿组合都是组合 $w^{a}$ 和 $w^{b}$ 的凸组合。
    \end{itemize}
\end{enumerate}
\end{problem}
\begin{solution}
\begin{enumerate}[label=(\alph*)]
    \item 无套利的条件是 $\phi_{\omega}>0$，$\sum\phi_{\omega}\le 1$；
    \item
    \begin{itemize}
        \item 收益率 $\tilde{r}(\omega)=\frac{w_{\omega}}{\phi_{\omega}}-1$，因而收益率的期望值是 $\bar{r}=\sum\pi_{\omega}\frac{w_{\omega}}{\phi_{\omega}}-1$，方差是 $\operatorname{Var}(\tilde{r})=\sum\pi_{\omega}\bigl(\frac{w_{\omega}}{\phi_{\omega}}\bigr)^{2}-\bigl(\sum\pi_{\omega}\frac{w_{\omega}}{\phi_{\omega}}\bigr)^{2}$
        \item
        \[
            \operatorname{Cov}(\tilde{r},\,\tilde{r}')=\sum\pi_{\omega}\frac{w_{\omega}}{\phi_{\omega}}\frac{w_{\omega}'}{\phi_{\omega}}-\sum\pi_{\omega}\frac{w_{\omega}}{\phi_{\omega}}\sum\pi_{\omega}\frac{w_{\omega}'}{\phi_{\omega}}
        \]
    \end{itemize}
    \item
    \begin{itemize}
        \item $w^{a}$ 的 1 期支付为 $\frac{w_{\omega}^{a}}{\phi_{\omega}}=\frac{1}{\phi_{0}}$ 是一常数；而它的价格为 1，即它是 0 期价格为 1 的无风险组合。
        \item 考虑下面的优化问题
        \[
            \begin{aligned}
                \max_{w_{\omega}}\;&\frac12\sum\pi_{\omega}\Bigl(\frac{w_{\omega}}{\phi_{\omega}}\Bigr)^{2}\\
                \text{s.t.}\;&\sum\pi_{\omega}\frac{w_{\omega}}{\phi_{\omega}}-1=\bar{r}_{p}\\
                &\sum w_{\omega}=1
            \end{aligned}
        \]
        构造 Lagrangian 函数
        \[
            \mathcal{L}=\frac12\sum\pi_{\omega}\Bigl(\frac{w_{\omega}}{\phi_{\omega}}\Bigr)^{2}+\lambda_{r}\Bigl(\bar{r}_{p}+1-\pi_{\omega}\frac{w_{\omega}}{\phi_{\omega}}\Bigr)+\lambda_{1}\Bigl(1-\sum w_{\omega}\Bigr)
        \]
        一阶条件为
        \[
            \pi_{\omega}\frac{w_{\omega}}{\phi_{\omega}^{2}}=\lambda_{r}\frac{w_{\omega}}{\phi_{\omega}}+\lambda_{1}
        \]
        \[
            \begin{aligned}
                \Rightarrow w_{\omega}&=\lambda_{r}\phi_{\omega}+\lambda_{1}\phi_{\omega}q_{\omega}\\
                &=\frac{(1+\bar{r}_{p})\sum\phi_{\omega}-1}{\sum\phi_{\omega}q_{\omega}-(\sum\phi_{\omega})^{2}}\phi_{\omega}q_{\omega}+\frac{\sum\phi_{\omega}-(1+\bar{r}_{p})\sum\phi_{\omega}q_{\omega}}{\sum\phi_{\omega}q_{\omega}-(\sum\phi_{\omega})^{2}}\phi_{\omega}\\
                &=\delta\frac{\phi_{\omega}q_{\omega}}{\sum\phi_{\omega}q_{\omega}}+(1-\delta)\frac{\phi_{\omega}}{\sum\phi_{\omega}}
            \end{aligned}
        \]
        其中，
        \[
            \delta=\frac{((1+\bar{r}_{p})\sum\phi_{\omega}-1)\sum\phi_{\omega}q_{\omega}}{\sum\phi_{\omega}q_{\omega}-(\sum\phi_{\omega})^{2}}
        \]
        因为约束都起限制作用，因而 $\lambda_{r}>0$，$\lambda_{1}>0$，因而任意前沿组合都是 $w^{a}$，$w^{b}$ 的凸组合。
    \end{itemize}
\end{enumerate}
\end{solution}

\begin{problem}
\problemno{13.6**} 继续考虑练习 13.5 中的经济。
\begin{enumerate}[label=(\alph*)]
    \item 令 $u(\cdot)$ 为代表性参与者的效用函数。证明为使 CAPM 成立 $u(\cdot)$ 必须具有平方函数的形式。
    \item 假设所有的资产都可由 CAPM 定价。证明一份在状态 $\omega$（$\omega\in\Omega$）下支付 $x_{\omega}$ 单位消费品的证券价格由下式给出
    \[
        p_{x}=\alpha\mathrm{E}[\tilde{x}]-\beta\operatorname{Cov}[\tilde{x},\tilde{e}].
    \]
    其中 $\alpha$ 和 $\beta$ 是两个和 $x$ 无关的常数。
\end{enumerate}
\end{problem}
\begin{solution}
\begin{enumerate}[label=(\alph*)]
    \item 对于支付为 $\tilde{x}$ 的资产而言，其价格为 $P_{x}=\mathrm{E}\bigl[\frac{u'(\tilde{e})}{u'(e_{0})}\tilde{x}\bigr]=\operatorname{Cov}\bigl(\frac{u'(\tilde{e})}{u'(e_{0})},\,\tilde{x}\bigr)+\frac{\bar{x}}{1+r_{F}}$，如果 CAPM 成立，那么有
    \[
        \begin{aligned}
            \frac{\bar{x}}{P_{x}}-(1+r_{F})&=\frac{\operatorname{Cov}(\tilde{x},\,\tilde{e})P_{e}}{P_{x}\operatorname{Var}(\tilde{e})}\Bigl(\frac{\bar{e}}{P_{e}}-(1+r_{F})\Bigr)\\
            \bar{x}-P_{x}(1+r_{F})&=\frac{\operatorname{Cov}(\tilde{x},\,\tilde{e})}{\operatorname{Var}(\tilde{e})}\bigl(\bar{e}-P_{e}(1+r_{F})\bigr)
        \end{aligned}
    \]
    因而对于任意的 $\tilde{x}$ 都有
    \[
        \operatorname{Cov}\Biggl(\tilde{x},\,\frac{\bar{e}-P_{e}(1+r_{F})}{\operatorname{Var}(\tilde{e})}\tilde{e}+\frac{u'(\tilde{e})(1+r_{F})}{u'(e_{0})}\Biggr)=0
    \]
    由此可得 $u'(\tilde{e})$ 是 $\tilde{e}$ 的线性函数，$u$ 是二次函数；

    \item 由上面的
    \[
        \bar{x}-P_{x}(1+r_{F})=\frac{\operatorname{Cov}(\tilde{x},\,\tilde{e})}{\operatorname{Var}(\tilde{e})}\bigl(\bar{e}-P_{e}(1+r_{F})\bigr)
    \]
    可得，
    \[
        \begin{aligned}
            P_{x}&=\frac{1}{1+r_{F}}\bar{x}-\frac{\bar{e}-P_{e}(1+r_{F})}{(1+r_{F})\operatorname{Var}(\tilde{e})}\operatorname{Cov}(\tilde{x},\,\tilde{e})\\
            &=\alpha\bar{x}-\beta\operatorname{Cov}(\tilde{x},\,\tilde{e})
        \end{aligned}
    \]
    其中，$\alpha=\frac{1}{1+r_{F}}$，$\beta=\frac{\bar{e}-P_{e}(1+r_{F})}{(1+r_{F})\operatorname{Var}(\tilde{e})}$ 是两个与 $\tilde{x}$ 无关的常数。
\end{enumerate}
\end{solution}

\begin{problem}
\problemno{13.7*} 考虑练习 13.5 中的经济。假设只有两个可能状态 $a$ 和 $b$，$e_{b}>e_{a}$。称支付总禀赋的证券为股票。考虑一份以股票为标的资产、执行价格为 $K\in(e_{a},e_{b})$ 的看涨期权。
\begin{enumerate}[label=(\alph*)]
    \item 能用 CAPM（代表性参与者具有一般形式的 $u(\cdot)$）为期权定价吗？解释所得的结论。
    \item 能把上面的结论推广到多余两个状态的情形吗？给出解释。
\end{enumerate}
\end{problem}
\begin{solution}
\begin{enumerate}[label=(\alph*)]
    \item 如果 1 期只有两个状态，并且 $e_{b}>e_{a}$，那么可由
    \[
        \bar{e}=\phi_{a}e_{a}+\phi_{b}e_{b},\quad\sigma_{e}^{2}=\phi_{a}(e_{a}-\bar{e})^{2}+\phi_{b}(e_{b}-\bar{e})^{2}
    \]
    求得 $e_{a}=e_{a}(\bar{e},\,\sigma_{e}^{2})$，$e_{b}=e_{b}(\bar{e},\,\sigma_{e}^{2})$，因此，
    \[
        \mathrm{E}[u(\tilde{e})]=\phi_{a}u(e_{a})+\phi_{b}u(e_{b})=v(\bar{e},\,\sigma_{e}^{2})
    \]
    即代表性参与者具有均值一方差偏好，按照正文中的结论，我们可以知道 CAPM 成立，因此可以用 CAPM 为期权定价；
    \item 如果状态数大于 2，我们不能用均值、方差表示出各个状态的禀赋，代表性参与者不一定具有均值方差偏好（如果效用函数是二次函数则具有均值一方差偏好），CAPM 不一定成立。
\end{enumerate}
\end{solution}



% !TeX root = ../../金融经济学习题讲义.tex
% 源 PDF：references/金融经济学/史上最全最完整的金融经济学答案(王江版).pdf，页 76–79（印刷页 75–78）

\zhchapter{套利定价理论（APT）}

\ifdefined\WJHandoutKnowledge
\fi

% !TeX root = ../../金融经济学习题讲义.tex
% 14_APT 课后习题（由章节文件拆分）

\begin{problem}
\problemno{14.1} 市场上有无穷多的证券。它们的收益率由单因子模型给定：
\[
    \tilde{r}_{n}=\bar{r}_{n}+b_{n}\tilde{f}+\tilde{u}_{n},\quad n=1,2,\ldots
\]
这里 $\mathrm{E}[\tilde{f}]=0$；$\forall\,n,n'$，$\mathrm{E}[\tilde{u}_{n}]=0$，$\mathrm{E}[\tilde{f}\tilde{u}_{n}]=0$；$\operatorname{Var}[\tilde{u}_{n}]\le\sigma^{2}$；以及 $\mathrm{E}[\tilde{u}_{n}\tilde{u}_{n'}]=0$。一个组合由 $\{w_{1},w_{2},\ldots\}$ 给定。如当 $N$ 很大时，
\[
    \sum_{n}^{N}w_{n}^{2}\simeq O(N^{-1})
\]
则称它是充分分散的，其中 $N$ 代表组合中绝对权重最大的股票数目而 $O(\cdot)$ 表示同阶或更高阶。
\begin{enumerate}[label=(\alph*)]
    \item 计算充分分散组合的方差。
    \item 证明它不受单个证券的特殊风险的影响。
    \item 更一般地讨论什么因素决定了充分分散组合的风险，这里风险的度量是收益率的方差。
\end{enumerate}
\end{problem}
\begin{solution}
\begin{enumerate}[label=(\alph*)]
    \item 组合收益率为
    \[
        \tilde{r}_{p}=\sum w_{n}\tilde{r}_{n}=\sum w_{n}\bar{r}_{n}+\Bigl(\sum w_{n}b_{n}\Bigr)\tilde{f}+\sum w_{n}\tilde{u}_{n}
    \]
    组合风险为
    \[
        \operatorname{Var}(\tilde{r}_{p})=\Bigl(\sum w_{n}b_{n}\Bigr)^{2}\sigma_{f}^{2}+\sum w_{n}^{2}\operatorname{Var}(\tilde{u}_{n})
    \]
    因为组合是充分分散的，故
    \[
        \sum w_{n}^{2}\operatorname{Var}(\tilde{u}_{n})\le\sigma^{2}\sum w_{n}^{2}\to 0,\quad N\to\infty
    \]
    因此
    \[
        \operatorname{Var}(\tilde{r}_{p})=\Bigl(\sum w_{n}b_{n}\Bigr)^{2}\sigma_{f}^{2}
    \]
    \item $\operatorname{Var}(\tilde{r}_{p})$ 中没有 $\operatorname{Var}(\tilde{u}_{n})$ 项，即组合风险不受特殊风险的影响；
    \item 影响充分分散组合的风险有两个：一是因子的风险，另一个是组合的载荷。
\end{enumerate}
\end{solution}

\begin{problem}
\problemno{14.2} 假设资产 $A$、$B$、$C$ 的收益服从精确两因子模型：
\begin{align}
    r_{A}&=a_{A}+b_{A,1}\tilde{f}_{1}+b_{A,2}\tilde{f}_{2}\tag{14.1}\\
    r_{B}&=a_{B}+b_{B,1}\tilde{f}_{1}+b_{B,2}\tilde{f}_{2}\tag{14.2}\\
    r_{C}&=a_{C}+b_{C,1}\tilde{f}_{1}+b_{C,2}\tilde{f}_{2}\tag{14.3}
\end{align}
其中 $\tilde{f}_{1}$ 和 $\tilde{f}_{2}$ 是两个均值为 0、方差为 1 的因子。$r_{F}$ 是无风险利率。为消除套利机会，$a$、$b$ 以及 $r_{F}$ 应该满足什么条件？
\end{problem}
\begin{solution}
构造一个组合，组合权重为 $\begin{bmatrix}x_{1} \\ x_{2} \\ 1-x_{1}-x_{2}\end{bmatrix}$，那么组合的收益率为
\[
    \begin{aligned}
        \tilde{r}_{p}&=\bigl[(a_{A}-a_{C})x_{1}+(a_{B}-a_{C})x_{2}\bigr]+\bigl[(b_{A,1}-b_{C,1})x_{1}+(b_{B,1}-b_{C,1})x_{2}\bigr]\tilde{f}_{1}\\
        &\quad+\bigl[(b_{A,2}-b_{C,2})x_{1}+(b_{B,2}-b_{C,2})x_{2}\bigr]\tilde{f}_{2}
    \end{aligned}
\]
如果不存在套利机会，那么
\[
    \begin{aligned}
        (b_{A,1}-b_{C,1})x_{1}+(b_{B,1}-b_{C,1})x_{2}&=0\\
        (b_{A,2}-b_{C,2})x_{1}+(b_{B,2}-b_{C,2})x_{2}&=0\\
        (a_{A}-a_{C})x_{1}+(a_{B}-a_{C})x_{2}&=r_{F}
    \end{aligned}
\]
也就是说，若 $x\in\{y\in\mathrm{R}^{2}:y\perp\begin{bmatrix}b_{A,1}-b_{C,1} \\ b_{B,1}-b_{C,1}\end{bmatrix},\,y\perp\begin{bmatrix}b_{A,2}-b_{C,2} \\ b_{B,2}-b_{C,2}\end{bmatrix}\}$，那么 $x^{\mathrm{T}}\begin{bmatrix}a_{A}-a_{C} \\ a_{B}-a_{C}\end{bmatrix}=r_{F}$。
\end{solution}

\begin{problem}
\problemno{14.3} 考虑一个有无穷多资产和 $K$ 个因子的因子模型：
\[
    \tilde{r}_{n}=\bar{r}_{n}+b_{n}\tilde{f}+\tilde{u}_{n},\quad n=1,2,\ldots
\]
其中 $\tilde{f}=\begin{bmatrix}\tilde{f}_{1} \\ \ldots \\ \tilde{f}_{K}\end{bmatrix}$ 是因子向量，$b_{n}=\begin{bmatrix}b_{n,1} \\ \ldots \\ b_{n,K}\end{bmatrix}$ 资产 $n$ 的因子载荷向量，而 $\forall\,n,n'$，$\mathrm{E}[\tilde{f}]=0$，$\mathrm{E}[\tilde{u}_{n}]=0$，$\mathrm{E}[\tilde{f}\tilde{f}^{\mathrm{T}}]=I$，$\mathrm{E}[\tilde{f}\tilde{u}_{n}]=0$，$\mathrm{E}[\tilde{u}_{n}\tilde{u}_{n'}]=0$，其中 $I$ 是单位阵。考虑另外一组因子 $\tilde{f}'$，它们是 $\tilde{f}$ 的线性变换：
\[
    \tilde{f}'=\mathrm{A}\tilde{f}
\]
这里 $\mathrm{A}$ 是一个 $(K\times K)$ 的矩阵。
\begin{enumerate}[label=(\alph*)]
    \item 证明存在 $\tilde{f}'$ 使得
    \[
        \tilde{r}_{n}=\bar{r}_{n}+b_{n}\tilde{f}'+\tilde{u}_{n},\quad n=1,2,\ldots
    \]
    其中 $\mathrm{E}[\tilde{f}']=0$，$\mathrm{E}[\tilde{f}'\tilde{f}'^{\mathrm{T}}]=I$，$\mathrm{E}[\tilde{f}'\tilde{u}_{n}]=0$，$\forall\,n$。因此，因子模型不是唯一的。
    \item 证明因子模型有许多因子不同但却等价的表述形式。
\end{enumerate}
\end{problem}
\begin{solution}
\begin{enumerate}[label=(\alph*)]
    \item 若 $A$ 为正交阵（即 $A^{\mathrm{T}}A=I$），那么
    \[
        \begin{aligned}
            \tilde{r}_{n}&=\bar{r}_{n}+b_{n}^{\mathrm{T}}\tilde{f}+\tilde{u}_{n}\\
            &=\bar{r}_{n}+b_{n}^{\mathrm{T}}A^{\mathrm{T}}A\tilde{f}+\tilde{u}_{n}\\
            &=\bar{r}_{n}+(Ab_{n})^{\mathrm{T}}(A\tilde{f})+\tilde{u}_{n}\\
            &=\bar{r}_{n}+b_{n}'^{\mathrm{T}}\tilde{f}'+\tilde{u}_{n}
        \end{aligned}
    \]
    可以验证：
    \[
        \mathrm{E}[\tilde{f}']=0,\quad\mathrm{E}[\tilde{f}'\tilde{f}'^{\mathrm{T}}]=I,\quad\mathrm{E}[\tilde{f}'\tilde{u}_{n}]=0
    \]
    \item 因为 $K$ 阶正交阵有很多，因而等价的表述形式也有很多。
\end{enumerate}
\end{solution}

\begin{problem}
\problemno{14.4*} 给定练习 14.3 描述的因子模型，APT 要求
\[
    \bar{r}_{n}-r_{F}=\lambda^{\mathrm{T}}b_{n}
\]
其中 $\lambda=\begin{bmatrix}\lambda_{1} \\ \lambda_{2} \\ \ldots \\ \lambda_{K}\end{bmatrix}$ 是因子风险溢价向量。证明存在一个等价的因子模型，它的因子向量为 $\tilde{f}^{*}$ 而对应的因子风险溢价向量为 $\lambda^{*}=\begin{bmatrix}\lambda_{1}^{*} \\ 0 \\ \ldots \\ 0\end{bmatrix}$，即只有第一个因子有风险溢价而其它因子虽然代表了不可分散风险，却没有风险溢价。
\end{problem}
\begin{solution}
考虑这样的正交阵：$A=\bigl[\frac{1}{\sqrt{\lambda^{\mathrm{T}}\lambda}}\lambda,\,U\bigr]$，其中
\[
    U=[U_{1},\,U_{2},\,\cdots,\,U_{K-1}],\quad U_{i}^{\mathrm{T}}U_{j}=\delta_{ij},\quad U_{i}^{\mathrm{T}}\lambda=0
\]
\[
    \delta_{ij}=\begin{cases}
        1&i=j\\
        0&i\neq j
    \end{cases}
\]
可以验证：$AA^{\mathrm{T}}=I$，即 $A$ 是正交阵，且 $A^{\mathrm{T}}\lambda=\begin{bmatrix}\sqrt{\lambda^{\mathrm{T}}\lambda} \\ \,0 \\ \,0 \\ \,\cdots \\ \,0\end{bmatrix}$，那么
\[
    \begin{aligned}
        \tilde{r}_{n}&=\bar{r}_{n}+b_{n}^{\mathrm{T}}AA^{\mathrm{T}}\tilde{f}+\tilde{u}_{n}\\
        &=\bar{r}_{n}+(A^{\mathrm{T}}b_{n})^{\mathrm{T}}(A^{\mathrm{T}}\tilde{f})+\tilde{u}_{n}\\
        &=\bar{r}_{n}+b_{n}^{*\,\mathrm{T}}\tilde{f}^{*}+\tilde{u}_{n}
    \end{aligned}
\]
是一个等价表述。并且 $\bar{r}_{n}-r_{F}=\lambda^{*\,\mathrm{T}}b_{n}^{*}$，其中，$\lambda^{*}=A^{\mathrm{T}}\lambda=\begin{bmatrix}\sqrt{\lambda^{\mathrm{T}}\lambda} \\ \,0 \\ \,0 \\ \,\cdots \\ \,0\end{bmatrix}$。
\end{solution}



% !TeX root = ../../金融经济学习题讲义.tex
% 源 PDF：references/金融经济学/史上最全最完整的金融经济学答案(王江版).pdf，页 80–86（印刷页 79–85）

\zhchapter{完全市场下的公司财务}

\ifdefined\WJHandoutKnowledge
\fi

% !TeX root = ../../金融经济学习题讲义.tex
% 15_完全市场下的公司财务 课后习题（由章节文件拆分）

\begin{problem}
\problemno{15.1*} 考虑一个具有两个时期 0 和 1 的经济。在 1 期有两个可能状态 $a$ 和 $b$，发生概率分别为 $\pi_{a}$ 和 $\pi_{b}$。存在生产技术如下：如果在 0 期投入 $I$，那么在 1 期的产出将是
\[
    Y=\begin{cases}
        \lambda_{a}I^{\alpha},&\text{状态 }a\\
        \lambda_{b}I^{\alpha},&\text{状态 }b
    \end{cases}
\]
其中 $0<\lambda_{b}<\lambda_{a}<1$ 且 $0<\alpha<1$。假设参与者 1 拥有这项技术但却没有任何其它资产。他具有如下形式的效用函数：
\[
    \log(c_{0})+\pi_{a}\log c_{1a}+\pi_{b}\log c_{1b}.
\]
证券市场上只存在无风险债券，它的利率为 $r_{F}$。
\begin{enumerate}[label=(\alph*)]
    \item 写出参与者 1 的优化问题。
    \item 求解他的最优生产、消费和投资选择。
    \item 说明他的生产决策如何依赖于利率。并给出解释。
    \item 计算在最优选择下他的期望效用。
\end{enumerate}
\end{problem}
\begin{solution}
\begin{enumerate}[label=(\alph*)]
    \item 假设参与者的 0 期消费为 $c_{0}$、投资为 $I$，那么参与者的优化问题是
    \[
        \begin{aligned}
            \max_{c_{0},\,I}\;&\log c_{0}+\pi_{a}\log(c_{a})+\pi_{b}\log(c_{b})\qquad(*)\\
            \text{s.t.}\;&c_{a}=\lambda_{a}I^{\alpha}-(c_{0}+I)(1+r_{F})\\
            &c_{b}=\lambda_{b}I^{\alpha}-(c_{0}+I)(1+r_{F})
        \end{aligned}
    \]
    \item 上述优化问题的一阶条件是
    \[
        \begin{aligned}
            \frac{1}{c_{0}}+\pi_{a}\frac{-(1+r_{F})}{\lambda_{a}I^{\alpha}-(c_{0}+I)(1+r_{F})}+\pi_{b}\frac{-(1+r_{F})}{\lambda_{b}I^{\alpha}-(c_{0}+I)(1+r_{F})}&=0\\
            \pi_{a}\frac{\lambda_{a}I^{\alpha-1}\alpha-(1+r_{F})}{\lambda_{a}I^{\alpha}-(c_{0}+I)(1+r_{F})}+\pi_{b}\frac{\lambda_{b}I^{\alpha-1}\alpha-(1+r_{F})}{\lambda_{b}I^{\alpha}-(c_{0}+I)(1+r_{F})}&=0
        \end{aligned}
    \]
    第一式是消费 Euler：$1/c_{0}=(1+r_{F})E[1/\tilde{c}_{1}]$。第二式说：用 1 期边际效用加权后，边际产出 $\alpha\lambda_{\omega}I^{\alpha-1}$ 的期望等于 $1+r_{F}$。两式联立决定最优 $c_{0}^{*}$、$I^{*}$（一般无更简的显式解）。另外参与者须在 0 期卖空 $(c_{0}^{*}+I^{*})$ 单位的无风险债券；
    \item 由生产一阶条件，边际效用加权的 $E[\mathrm{MPK}]$ 必须等于 $1+r_{F}$。$\alpha<1$ 时 MPK 随 $I$ 递减，故 $r_{F}$ 上升迫使 $I^{*}$ 下降：借款为生产融资的成本提高。
    \item 把求得的 $c_{0}^{*}$，$I^{*}$ 代入 $(*)$ 式即可求得期望效用，记作 $EU_{1}$。
\end{enumerate}
\end{solution}

\begin{problem}
\problemno{15.2*} 在练习 15.1 描述的经济中，现在假设证券市场还包括一只证券，它在状态 $a$ 支付 $\lambda_{a}$、在状态 $b$ 支付 $\lambda_{b}$ 而它的价格为 $S$。
\begin{enumerate}[label=(\alph*)]
    \item 证明现在证券市场是完全的。
    \item 用 $r_{F}$ 和 $S$ 表示状态价格。
    \item 描述参与者 1 的优化问题。
    \item 计算他的总财富。证明它总是正的。
    \item 求解他的最优生产、消费和投资选择。
    \item 用 NPV 法则解释他的投资选择。
    \item 他的生产技术的 NPV 是多少？这个如何随着状态价格的改变而改变？解释所得到的结论。
    \item 计算引入风险证券参与者 1 得到的期望效用的增加值，用 0 期财富的确定性等价表示。
\end{enumerate}
\end{problem}
\begin{solution}
\begin{enumerate}[label=(\alph*)]
    \item 两个状态，两只支付独立的证券，因而证券市场是完全的；
    \item 由
    \[
        \begin{aligned}
            \frac{1}{1+r_{F}}&=\phi_{a}+\phi_{b}\\
            S&=\phi_{a}\lambda_{a}+\phi_{b}\lambda_{b}
        \end{aligned}
    \]
    可以求得状态价格为
    \[
        \begin{aligned}
            \phi_{a}&=\frac{S(1+r_{F})-\lambda_{b}}{(\lambda_{a}-\lambda_{b})(1+r_{F})}\\
            \phi_{b}&=\frac{\lambda_{a}-S(1+r_{F})}{(\lambda_{a}-\lambda_{b})(1+r_{F})}
        \end{aligned}
    \]
    \item 假设参与者的 0 期消费为 $c_{0}$、股票头寸为 $\theta_{S}$、生产投入为 $I$，那么他的优化问题是
    \[
        \begin{aligned}
            \max_{c_{0},\,\theta_{S},\,I}\;&\log c_{0}+\pi_{a}\log(c_{a})+\pi_{b}\log(c_{b})\\
            \text{s.t.}\;&c_{a}=\lambda_{a}I^{\alpha}+\theta_{S}\lambda_{a}-(\theta_{S}S+c_{0}+I)(1+r_{F})\\
            &c_{b}=\lambda_{b}I^{\alpha}+\theta_{S}\lambda_{b}-(\theta_{S}S+c_{0}+I)(1+r_{F})
        \end{aligned}
    \]
    \item 参与者的总财富为
    \[
        \begin{aligned}
            w&=c_{0}+\phi_{a}\bigl(\lambda_{a}I^{\alpha}+\theta_{S}\lambda_{a}-(\theta_{S}S+c_{0}+I)(1+r_{F})\bigr)\\
            &\quad+\phi_{b}\bigl(\lambda_{b}I^{\alpha}+\theta_{S}\lambda_{b}-(\theta_{S}S+c_{0}+I)(1+r_{F})\bigr)\\
            &=SI^{\alpha}-I
        \end{aligned}
    \]
    因为当 $I=0$ 时，$w=0$，因而在参与者选择最优生产投入时，$w>0$；
    \item     完全市场下生产与消费分离。把财富约束写成 $c_{0}+\phi_{a}c_{a}+\phi_{b}c_{b}=SI^{\alpha}-I$。拉格朗日函数为
    \[
        \mathcal{L}=\log c_{0}+\pi_{a}\log c_{a}+\pi_{b}\log c_{b}+\lambda\bigl(SI^{\alpha}-I-c_{0}-\phi_{a}c_{a}-\phi_{b}c_{b}\bigr)
    \]
    对 $I$ 的一阶条件即 NPV 法则
    \[
        \alpha SI^{\alpha-1}=1\implies I^{*}=(\alpha S)^{\frac{1}{1-\alpha}}
    \]
    消费侧一阶条件：$1/c_{0}=\lambda$，$\pi_{\omega}/c_{\omega}=\lambda\phi_{\omega}$。代入预算并注意 $\pi_{a}+\pi_{b}=1$，得 $c_{0}^{*}=w^{*}/2$，$c_{\omega}^{*}=\pi_{\omega}w^{*}/(2\phi_{\omega})$，其中 $w^{*}=SI^{*\alpha}-I^{*}$。再用 $c_{a}$、$c_{b}$ 反解股票头寸，得到
    \[
        I^{*}=(\alpha S)^{\frac{1}{1-\alpha}},\quad w^{*}=SI^{*\alpha}-I^{*},\quad c_{0}^{*}=\tfrac12 w^{*},\quad\theta_{S}^{*}=\frac{\frac{\pi_{a}}{2\phi_{a}}-\frac{\pi_{b}}{2\phi_{b}}}{\lambda_{a}-\lambda_{b}}w^{*}-I^{*\alpha}
    \]
    另外，参与者须在 0 期卖空 $\theta_{S}^{*}S+c_{0}^{*}+I^{*}$ 单位的无风险债券；
    \item
    \[
        NPV=(\phi_{a}\lambda_{a}I^{\alpha}+\phi_{b}\lambda_{b}I^{\alpha})-I=SI^{\alpha}-I
    \]
    最优的生产投入应是 $(NPV)'=0$，即 $I=(\alpha S)^{\frac{1}{1-\alpha}}$，恰是上面所求得的最优生产投入；
    \item 代入计算得
    \[
        \begin{aligned}
            NPV^{*}&=S(\alpha S)^{\frac{\alpha}{1-\alpha}}-(\alpha S)^{\frac{1}{1-\alpha}}\\
            &=(\alpha^{\frac{\alpha}{1-\alpha}}-\alpha^{\frac{1}{1-\alpha}})(\phi_{a}\lambda_{a}+\phi_{b}\lambda_{b})^{\frac{1}{1-\alpha}}
        \end{aligned}
    \]
    由此可以看到，NPV 会随着状态价格的上升而增加：一方面是在相同生产投入条件下产出的净现值增加；另外，根据 NPV 法则，参与者会增加生产投入，这也会使得 NPV 增加；
    \item 在最优选择下，我们可计算得期望效用为 $EU_{2}$。在第 1 题中，0 期禀赋为 0 时计算得到的期望效用是 $EU_{1}$，我们可以通过 $EU_{1}(\Delta w)=EU_{2}$ 求效用增加值的确定性等价 $\Delta w$。
\end{enumerate}
\end{solution}

\begin{problem}
\problemno{15.3**} 考虑练习 15.1 描述的经济。现在假设经济中除了那里描述的参与者 1 以外还有一个参与者 2。参与者 2 没有生产技术，但经济赋予了他 $e_{0}$ 单位的 0 期消费。他的效用函数和参与者 1 的一样。证券市场和练习 15.2 中的一样。
\begin{enumerate}[label=(\alph*)]
    \item 求解两个参与者在给定状态价格下的最优生产、消费和投资选择。
    \item 求解均衡状态价格、总消费和消费配置。
    \item 证明均衡配置是 Pareto 最优的。
\end{enumerate}
\end{problem}
\begin{solution}
\begin{enumerate}[label=(\alph*)]
    \item 在第 2 题中已经求解了参与者 1 的最优化问题，现在考虑参与者 2 的最优化问题（假设他持有的无风险债券、股票的头寸分别为 $\theta_{B}$，$\theta_{S}$）
    \[
        \begin{aligned}
            \max_{\theta_{B},\,\theta_{S}}\;&\log c_{0}+\pi_{a}\log(c_{a})+\pi_{b}\log(c_{b})\\
            \text{s.t.}\;&c_{0}=e_{0}-\theta_{B}-S\theta_{S}\\
            &c_{a}=\theta_{B}(1+r_{F})+\lambda_{a}\theta_{S}\\
            &c_{b}=\theta_{B}(1+r_{F})+\lambda_{b}\theta_{S}
        \end{aligned}
    \]
    完全市场下与练习 15.2 同一拉格朗日，Euler 给出 $c_{0}^{*}=e_{0}/2$，$c_{\omega}^{*}=\pi_{\omega}e_{0}/(2\phi_{\omega})$。由 1 期支付反解头寸：
    \[
        \theta_{S}^{*}=\frac{c_{a}^{*}-c_{b}^{*}}{\lambda_{a}-\lambda_{b}}=\frac{1}{\lambda_{a}-\lambda_{b}}\Bigl(\frac{\pi_{a}}{2\phi_{a}}-\frac{\pi_{b}}{2\phi_{b}}\Bigr)e_{0},\qquad
        \theta_{B}^{*}=\frac{c_{a}^{*}-\lambda_{a}\theta_{S}^{*}}{1+r_{F}}
    \]
    \item 市场出清条件是
    \[
        \begin{aligned}
            \theta_{1,S}^{*}S+c_{1,0}^{*}+I^{*}&=\theta_{2,b}^{*}\\
            \theta_{1,S}^{*}+\theta_{2,S}^{*}&=0
        \end{aligned}
    \]
    再加上 $\frac{1}{1+r_{F}}=\phi_{a}+\phi_{b}$，$S=\phi_{a}\lambda_{a}+\phi_{b}\lambda_{b}$ 可以求得均衡状态价格为
    \[
        \begin{aligned}
            \phi_{a}&=\frac{\pi_{a}e_{0}^{1-\alpha}}{\lambda_{a}\alpha^{\alpha}(1+\alpha)^{1-\alpha}}\\
            \phi_{b}&=\frac{\pi_{b}e_{0}^{1-\alpha}}{\lambda_{b}\alpha^{\alpha}(1+\alpha)^{1-\alpha}}
        \end{aligned}
    \]
    最优的生产投入为 $I^{*}=\frac{\alpha e_{0}}{1+\alpha}$，最优消费为
    \[
        \begin{aligned}
            c_{1,0}^{*}&=\frac{1-\alpha}{2(1+\alpha)}e_{0},\quad c_{2,0}^{*}=\tfrac12 e_{0},\quad C_{0}^{*}=\frac{1}{1+\alpha}e_{0}\\
            c_{1,a}^{*}&=\frac{\pi_{a}}{2\phi_{a}}\frac{1-\alpha}{1+\alpha}e_{0},\quad c_{2,a}^{*}=\frac{\pi_{a}}{2\phi_{a}}e_{0},\quad C_{a}^{*}=\lambda_{a}\Bigl(\frac{\alpha e_{0}}{1+\alpha}\Bigr)^{\alpha}\\
            c_{1,b}^{*}&=\frac{\pi_{b}}{2\phi_{b}}\frac{1-\alpha}{1+\alpha}e_{0},\quad c_{2,b}^{*}=\frac{\pi_{b}}{2\phi_{b}}e_{0},\quad C_{b}^{*}=\lambda_{b}\Bigl(\frac{\alpha e_{0}}{1+\alpha}\Bigr)^{\alpha}
        \end{aligned}
    \]
    \item 考虑两个参与者的边际效用之比：
    \[
        \begin{aligned}
            \frac{\pi_{a}c_{1,0}^{*}}{c_{1,a}^{*}}&=\phi_{a}=\frac{\pi_{a}c_{2,0}^{*}}{c_{2,a}^{*}}\\
            \frac{\pi_{b}c_{1,0}^{*}}{c_{1,b}^{*}}&=\phi_{b}=\frac{\pi_{b}c_{2,0}^{*}}{c_{2,b}^{*}}
        \end{aligned}
    \]
    因而均衡配置是 Pareto 最优的。
\end{enumerate}
\end{solution}

\begin{problem}
\problemno{15.4**} 考虑练习 15.3 中的经济。
\begin{enumerate}[label=(\alph*)]
    \item 构建代表性参与者。
    \item 通过求解代表性参与者的优化问题而得到均衡配置。
    \item 用代表性参与者的边际效用计算状态价格。
    \item 讨论 $\lambda_{b}$、$\lambda_{a}$ 以及 $\alpha$ 如何影响状态价格和均衡配置。
\end{enumerate}
\end{problem}
\begin{solution}
\begin{enumerate}[label=(\alph*)]
    \item 我们知道，对于对数形式的效用函数，代表性参与者的效用函数也是对数形式的，且与财富在参与者之间的分布无关；
    \item 对于代表性参与者，他的优化问题是
    \[
        \max_{I}\ \log(e_{0}-I)+\pi_{a}\log(\lambda_{a}I^{\alpha})+\pi_{b}\log(\lambda_{b}I^{\alpha})
    \]
    目标函数等于 $\log(e_{0}-I)+\alpha\log I+$ 常数。一阶条件
    \[
        -\frac{1}{e_{0}-I}+\frac{\alpha}{I}=0
    \]
    即 $I=\alpha(e_{0}-I)$，故 $I^{*}=\dfrac{\alpha e_{0}}{1+\alpha}$，$c_{0}^{*}=\dfrac{e_{0}}{1+\alpha}$。1 期总消费为 $C_{\omega}^{*}=\lambda_{\omega}(I^{*})^{\alpha}$；
    \item 状态价格为
    \[
        \begin{aligned}
            \phi_{a}&=\frac{\pi_{a}C_{0}}{C_{a}}=\pi_{a}\frac{\frac{e_{0}}{1+\alpha}}{\lambda_{a}\bigl(\frac{\alpha e_{0}}{1+\alpha}\bigr)^{\alpha}}=\frac{\pi_{a}e_{0}^{1-\alpha}}{\lambda_{a}\alpha^{\alpha}(1+\alpha)^{1-\alpha}}\\
            \phi_{b}&=\frac{\pi_{b}e_{0}^{1-\alpha}}{\lambda_{b}\alpha^{\alpha}(1+\alpha)^{1-\alpha}}
        \end{aligned}
    \]
    \item 均衡配置只受到 $\alpha$ 的影响，$\alpha$ 增大，相同生产投入的产出增加，因而参与者会减少当前消费，增加生产投入；当 $\alpha$ 增大时，状态价格均变小，而当 $\lambda_{a}$，$\lambda_{b}$ 增加时，相应状态的状态价格变小，因而 $\alpha$，$\lambda_{a}$，$\lambda_{b}$ 增加会增加产出，因而状态价格变小。
\end{enumerate}
\end{solution}

\begin{problem}
\problemno{15.5**} 与练习 15.2 中的经济一样。在这里我们不假设有一组给定的可交易证券，而是假设拥有生产技术的参与者 1 可以成立一个公司、而公司可以到市场上发行债券和股票，但没有其它可交易的证券。一单位（无风险）公司债在 1 期有 1 单位支付，而一股股票在 1 期有 $\hat{\lambda}I^{\alpha}$ 单位支付，在状态 $a$ 和 $b$ 下 $\hat{\lambda}=\lambda_{a}$，$\lambda_{b}$，$I$ 是公司的生产投资。
\begin{enumerate}[label=(\alph*)]
    \item 给定利率 $r_{F}$ 和股价 $S$，求解参与者 1 的优化问题，即最优的投资额、在市场发行的债券和股票的数量以及他的消费。
    \item 在这种情况下，有无最优负债水平？这和 Modigliani--Miller 定理是否一致？解释原因。
    \item 求解市场均衡，特别地，求解均衡利率、股价、生产水平、发行的债券和股票的数量以及两个参与者的消费配置。
\end{enumerate}
\end{problem}
\begin{solution}
\begin{enumerate}[label=(\alph*)]
    \item 假设参与者的投资额为 $I$，在市场上发行的债券和股票分别为 $\theta_{B}$，$\theta_{S}$，那么他的优化问题是
    \[
        \begin{aligned}
            \max_{\theta_{B},\,\theta_{S},\,I}\;&\log c_{0}+\pi_{a}\log(c_{a})+\pi_{b}\log(c_{b})\\
            \text{s.t.}\;&c_{0}=\theta_{B}+S\theta_{S}-I\\
            &c_{a}=\lambda_{a}I^{\alpha}(1-\theta_{S})-\theta_{B}(1+r_{F})\\
            &c_{b}=\lambda_{b}I^{\alpha}(1-\theta_{S})-\theta_{B}(1+r_{F})
        \end{aligned}
    \]
    一阶条件为
    \[
        \begin{aligned}
            \frac{1}{c_{0}}+\pi_{a}\frac{-(1+r_{F})}{c_{a}}+\pi_{b}\frac{-(1+r_{F})}{c_{b}}&=0\\
            \frac{S}{c_{0}}-\frac{\pi_{a}\lambda_{a}I^{\alpha}}{c_{a}}-\frac{\pi_{b}\lambda_{b}I^{\alpha}}{c_{b}}&=0\\
            \frac{1}{c_{0}}-\frac{\pi_{a}\lambda_{a}\alpha I^{\alpha-1}(1-\theta_{S})}{c_{a}}+\frac{\pi_{b}\lambda_{b}\alpha I^{\alpha-1}(1-\theta_{S})}{c_{b}}&=0
        \end{aligned}
    \]
    由上述方程组可以解得最优的 $\theta_{B}^{*}$，$\theta_{S}^{*}$，$I^{*}$。
    \item 给定 $r_{F}$，$S$，参与者的债券和股票的最优发行量是定值，与 MM 定理不一致，这是因为此时的证券市场是不完全的。
    \item 对于参与者 2 来说，他的最优化问题是
    \[
        \begin{aligned}
            \max_{\theta_{B}',\,\theta_{S}'}\;&\log c_{0}+\pi_{a}\log(c_{a})+\pi_{b}\log(c_{b})\\
            \text{s.t.}\;&c_{0}=e_{0}-\theta_{B}'-S\theta_{S}'\\
            &c_{a}=\lambda_{a}I^{\alpha}\theta_{S}'+\theta_{B}'(1+r_{F})\\
            &c_{b}=\lambda_{b}I^{\alpha}\theta_{S}'+\theta_{B}'(1+r_{F})
        \end{aligned}
    \]
    求解出最优的 $\theta_{B}'^{*}$，$\theta_{S}'^{*}$ 以后，再加上市场出清条件
    \[
        \theta_{B}^{*}=\theta_{B}'^{*},\quad\theta_{S}^{*}=\theta_{S}'^{*}
    \]
    就可以求解出均衡利率、均衡股价，从而可以得到均衡的生产水平、债券和股票的发行量、以及消费配置。
\end{enumerate}
\end{solution}

\begin{problem}
\problemno{15.6*} 在 15.5 现在假设市场上存在一只无风险证券可供参与者交易。Modigliani--Miller 定理成立，即公司的融资决策不影响它的投资决策和均衡配置。解释这种情形和上述情形之间的差别。
\end{problem}
\begin{solution}
如果市场上有无风险证券可以交易，假设参与者 1 购买的债券数量为 $\theta_{B}'$，那么他的优化问题就变成了
\[
    \begin{aligned}
        \max_{\theta_{B},\,\theta_{B}',\,\theta_{S},\,I}\;&\log c_{0}+\pi_{a}\log(c_{a})+\pi_{b}\log(c_{b})\\
        \text{s.t.}\;&c_{0}=\theta_{B}+\theta_{B}'+S\theta_{S}-I\\
        &c_{a}=\lambda_{a}I^{\alpha}(1-\theta_{S})-(\theta_{B}+\theta_{B}')(1+r_{F})\\
        &c_{b}=\lambda_{b}I^{\alpha}(1-\theta_{S})-(\theta_{B}+\theta_{B}')(1+r_{F})
    \end{aligned}
\]
这时我们只能得到 $\theta_{B}+\theta_{B}'$ 的值，而不能确定债券的发行量，因而公司债券发行量的大小（融资决策）不影响公司的投资决策和均衡配置。这是因为公司除了发行债券外，还可以到市场上买卖无风险债券，可以通过调整无风险债券的头寸来补充公司债券发行量的增减，因此债券的发行量不再影响到公司的投资决策以及消费配置，
\end{solution}




\end{document}
